% ============================================================================
%  광학 격자에서 제타 풍경까지의 위상 양자화:
%  통합 프레임워크
%  ── 한국어판 ──
%  컴파일: xelatex unified_ko.tex (2회 실행)
% ============================================================================

% ── 논문 구조 맵 ──────────────────────────────────────────────────────────
%
% 카테고리 시스템 (scripts/agents/paper_categories.md 참조):
%   Paper A (ξ-다발 기하학) — 메인 논문, 이 파일
%   Paper B (스펙트럼 통계) — §구조적 대응에 통합
%   Paper C (GL(2) 확장) — 부록, 향후 xi-bundle-gl2로 분리
%   Paper D (양자 시뮬레이션) — 별도 저장소 xi-bundle-quantum
%
% ──────────────────────────────────────────────────────────────────────────
\documentclass[11pt, a4paper]{article}

% ── 한글 지원 (XeLaTeX + kotex) ────────────────────────────────────────────
\usepackage{kotex}
\usepackage{fontspec}

% ── Layout ─────────────────────────────────────────────────────────────────
\usepackage[top=2.5cm, bottom=2.5cm, left=2.8cm, right=2.8cm]{geometry}
\usepackage{setspace}
\onehalfspacing

% ── Mathematics ────────────────────────────────────────────────────────────
\usepackage{amsmath, amssymb, amsthm, mathtools}
\usepackage{bm}

% ── Graphics & Diagrams ───────────────────────────────────────────────────
\usepackage{graphicx}
\usepackage{tikz}
\usetikzlibrary{
  arrows.meta, positioning, calc, decorations.pathreplacing,
  shapes.geometric, fit, backgrounds, chains, matrix
}
\usepackage{pgfplots}
\pgfplotsset{compat=1.18}

% ── Tables ─────────────────────────────────────────────────────────────────
\usepackage{booktabs}
\usepackage{array}
\usepackage{multirow}
\usepackage{colortbl}
\usepackage{longtable}

% ── Colors ─────────────────────────────────────────────────────────────────
\usepackage[dvipsnames,svgnames,x11names]{xcolor}
\definecolor{txblue}{HTML}{2563EB}
\definecolor{chgreen}{HTML}{059669}
\definecolor{rxorange}{HTML}{D97706}
\definecolor{lossred}{HTML}{DC2626}
\definecolor{decodepurple}{HTML}{7C3AED}
\definecolor{secgray}{HTML}{374151}
\definecolor{lightbg}{HTML}{F8FAFC}
\definecolor{accentblue}{HTML}{3B82F6}
\definecolor{zetaviolet}{HTML}{6D28D9}
\definecolor{primegold}{HTML}{B45309}

% ── Algorithms ────────────────────────────────────────────────────────────
\usepackage[ruled, vlined, linesnumbered]{algorithm2e}
\SetKwComment{Comment}{$\triangleright$\ }{}

% ── Cross-references & Links ──────────────────────────────────────────────
\usepackage[colorlinks=true, linkcolor=accentblue, citecolor=chgreen, urlcolor=txblue]{hyperref}
\usepackage[capitalise, noabbrev]{cleveref}

% ── Theorem Environments (한국어) ─────────────────────────────────────────
\theoremstyle{definition}
\newtheorem{definition}{정의}[section]
\newtheorem{property}{성질}[section]
\newtheorem{remark}{비고}[section]
\newtheorem{example}{예제}[section]
\theoremstyle{plain}
\newtheorem{theorem}{정리}[section]
\newtheorem{lemma}[theorem]{보조정리}
\newtheorem{proposition}[theorem]{명제}
\newtheorem{corollary}[theorem]{따름정리}
\newtheorem{conjecture}[theorem]{추측}
\newtheorem{observation}[theorem]{관찰}
\crefname{observation}{관찰}{관찰}
\Crefname{observation}{관찰}{관찰}

% ── Custom Commands ───────────────────────────────────────────────────────
\newcommand{\Zq}{\mathbb{Z}_q}
\newcommand{\Rq}{\mathbb{Z}_q[X]/(X^n+1)}
\newcommand{\norm}[1]{\left\|#1\right\|}
\newcommand{\abs}[1]{\left|#1\right|}
\newcommand{\ceil}[1]{\left\lceil#1\right\rfloor}
\newcommand{\floor}[1]{\left\lfloor#1\right\rfloor}
\newcommand{\FrFT}{\mathrm{FrFT}}
\newcommand{\STE}{\mathrm{STE}}
\newcommand{\ie}{\textit{i.e.}}
\newcommand{\eg}{\textit{e.g.}}
\newcommand{\sectionline}{\noindent\rule{\textwidth}{0.4pt}\vspace{0.3em}}

% Resonance-specific commands
\newcommand{\xif}{\xi}
\newcommand{\Phiz}{\Phi_\xi}
\newcommand{\Ampl}{A}
\newcommand{\Psip}{\Psi}
\newcommand{\pis}{\pi_{\mathrm{s}}}
\newcommand{\slift}{\sigma^*\!-\!\tfrac{1}{2}}
\newcommand{\Lhat}{\hat{L}}
\newcommand{\GateA}{\mathsf{G}_{\mathrm{A}}}
\newcommand{\GateB}{\mathsf{G}_{\mathrm{B}}}
\newcommand{\Lpqo}{\mathcal{L}_{\mathrm{pqo}}}
\newcommand{\Ltgt}{\mathcal{L}_{\mathrm{tgt}}}
\newcommand{\Lcurv}{\mathcal{L}_{\mathrm{curv}}}
\newcommand{\Lres}{\mathcal{L}_{\mathrm{res}}}
\newcommand{\PQO}{\mathsf{PQO}}
\newcommand{\Ftwo}{F_{2}}

% ── Header Box Style ─────────────────────────────────────────────────────
\usepackage{tcolorbox}
\tcbuselibrary{skins, breakable}

% ============================================================================
\begin{document}

% ── 제목 ──────────────────────────────────────────────────────────────────
\begin{center}
  {\LARGE\bfseries 광학 격자에서 제타 풍경까지의\\[4pt]
   위상 양자화}\\[14pt]
  {\large\itshape 격자 암호, 프레넬 광학, 해석적 정수론을\\
   연결하는 통합 프레임워크}\\[20pt]
  {\large 김영후}\\[4pt]
  {\normalsize 독립 연구자}\\[4pt]
  {\small\today}
\end{center}

\vspace{1.5em}

% ── 초록 ──────────────────────────────────────────────────────────────────
\begin{tcolorbox}[
  colback=lightbg, colframe=zetaviolet!60!black,
  title={\bfseries 초록}, fonttitle=\large,
  arc=3pt, boxrule=0.8pt, toptitle=3pt, bottomtitle=3pt
]
본 논문은 세 가지 겉보기에 이질적인 영역을 연결하는 공유 수학적 문법---\textit{대수적 구조 위의 위상 양자화}---을 규명한다. 그 세 영역은 CRYSTALS-Kyber 격자 기반 양자후 암호, 자유 공간 프레넬 회절, 그리고 임계선 위 리만 $\xif$-함수의 위상 구조이다. 구체적 핵은 QROP-Net으로, Kyber 암호문을 광학 위상 패턴으로 인코딩하고 딥러닝 위상 복원을 통해 복원하는 완전 미분 가능 파이프라인이다. 연속 위상 예측을 이산 시클로토믹 격자 수준으로 강제하는 $\sin^2$ 기반 양자화 손실이, 공명소수법칙 프레임워크에서 ``공명 소수''를 특성화하는 보어--좀머펠트 위상 양자화 조건 $\Delta\nu = m\pi$와 구조적으로 동형임을 보인다. 나아가 Kyber의 LWE 오차 임계치 $\lfloor q/4\rfloor = 1920$과 임계선으로부터의 $\sigma$-리프트 편차 $\sigma^*-\tfrac{1}{2}$가, 신호 복원이 실패하는 경성 오차 경계로서 동일한 역할을 공유함을 증명한다. 외부 정답 없이 자가 교정하는 다원군(Polyadic) ODE 기반 영점 탐지기가 QROP-Net의 블라인드 테스트 시간 최적화를 반영한다. 7개의 구조적 대응이 형식화되며, 유비와 증명 사이의 경계가 명시적으로 서술된다.

\smallskip
\noindent 일곱 가지 검증 축---국소 곡률, 전역 스펙트럼 통계, 위상적 불변량,
디리클레 $L$-함수 확장, 음성 대조군, GL(2) 확장, GL(3) 확장, 교차-차수 모노드로미---에 걸친 62개의 수치 결과가
이 프레임워크의 내적 정합성을 지지한다.

\smallskip
\noindent\textit{개정 후 부기 (2026-04-08).} 임계선 위의 시클로토믹 PQO 프로그램을 경험적으로 실현하려는 시도 — $L=4$ 코사인 PQO 와 \cref{eq:psi_n}의 합성 이상 위상 $\Psip(n)$ 을 감독 타겟으로 사용하는 판별 네트워크 — 는 결정적 음성을 반환한다: 세 시드 모두에서 결과 위상 분포가 리만 영점과 비영점 종좌표를 구별하지 못하며 ($\chi^2$ p값 $0.13/0.53/0.41$; 섹터 $k=2,3$ 정확히 비어 있음). 진단 (\cref{rem:cos2_pqo,rem:psi_defect}) 은 합성 $\Psip$ 가 $\|X\|$에만 의존하고 $\xif$-영점 정보를 전혀 포함하지 않으므로 4-항 손실의 어느 항도 $t$-국소 영점 라벨을 학습에 주입하지 못하며, \cref{eq:pqo_cos} (학습 포텐셜) 와 \cref{eq:cyclo_label} (측정) 가 운영적으로 분리되어 있다는 것이다. 이하 본문은 합성 $\Psip$ 를 역사적 객체로, 시클로토믹 PQO 를 대수적 타겟으로 보존하되, 양자 모두 감독 교체식 $\Psip(t) := \arg\xif(\tfrac{1}{2}+it)$ (\cref{eq:psi_supervised}) 을 사용하는 실험을 기다리는 항목으로 표시된다.
\end{tcolorbox}

\vspace{0.5em}
\noindent\textbf{핵심어:} 위상 양자화, 시클로토믹 다항식환, CRYSTALS-Kyber,
분수 푸리에 변환, 리만 $\xif$-함수, 보어--좀머펠트 조건, 공명 소수,
$\sigma$-리프트, 베리 위상, 디리클레 $L$-함수

\tableofcontents
\newpage


% ============================================================================
\section{서론}
\label{sec:intro}
% ============================================================================

다항식환 $R_q = \Rq$에서 $n = 256$, $q = 7681$은 NIST 표준화 양자후 키 캡슐화 메커니즘 CRYSTALS-Kyber~\cite{kyber2021}의 대수적 뼈대이다. 정의 관계식 $X^n + 1 = 0$---동치적으로 $X^n \equiv -1$, 본 프로젝트 코드명 $k^n = -1$의 기원---은 시클로토믹 다항식이며, 그 근은 $2n$차 단위근이다:
\begin{equation}
  \omega_j = e^{i\pi(2j+1)/n}, \quad j = 0, 1, \ldots, n-1.
  \label{eq:cyclo_roots}
\end{equation}
이 $n$개의 위상은 단위원 위에 등간격으로 배치되어 $\Delta\theta = \pi/n$만큼 분리된다. 이 이산적이고 주기적인 위상 구조는 단순히 수론적 변환(NTT)을 통한 빠른 다항식 곱셈의 대수적 편의가 아니다; 우리는 이것이 적어도 두 가지 다른 영역에서 나타나는 보다 일반적인 \textit{위상 양자화 원리}의 한 사례라고 주장한다:

\begin{enumerate}
  \item \textbf{프레넬 광학}: 분수 푸리에 변환(FrFT)은 결맞음 파면 $U_0 = e^{i\phi}$를 자유 공간 회절을 통해 전파시키며, 분수차 $\alpha$가 전파 거리를 위상 매개변수로 인코딩한다.
  \item \textbf{해석적 정수론}: 완성 리만 $\xif$-함수 $\xif(s) = \tfrac{1}{2}s(s-1)\pi^{-s/2}\Gamma(s/2)\zeta(s)$는 임계선 $s = \tfrac{1}{2} + it$ 위에서 진폭 $\Ampl(t) = |\xif(\tfrac{1}{2}+it)|$과 위상 $\Phiz(t) = \arg\,\xif(\tfrac{1}{2}+it)$으로 분해되며, 위상 양자화 조건은 소수 분포 및 리만 가설과 밀접하게 연결된다.
\end{enumerate}

\textbf{논제.} QROP-Net의 $\sin^2$ 기반 양자화 손실, 임계선 위의 보어--좀머펠트 조건 $\Delta\nu = m\pi$, 그리고 시클로토믹 근 구조 $X^n \equiv -1$은 하나의 수학적 문법의 세 가지 발현이다: \textit{대수적 구조 위의 주기적 위상 벌칙이 연속 신호를 이산 격자 수준으로 강제하며, 해독이 실패하는 경성 오차 경계가 존재한다}.

\textbf{기여.}
\begin{enumerate}
  \item Kyber 격자 암호, 프레넬 광학, 딥러닝 위상 복원을 통합하는 완전 미분 가능 파이프라인 QROP-Net을 제시한다 (제1부).
  \item 위상 양자화 메커니즘을 이산 및 연속 푸리에 해석을 잇는 일반 원리로 추상화한다 (제~\ref{sec:bridge}절).
  \item 이 원리를 $\xif(s)$의 위상 구조에 적용하여, 소수 탐지를 위한 2-게이트 공명 기준을 공식화한다 (제2부, 제~\ref{sec:xi_phase}--\ref{sec:gates}절).
  \item LWE 오차 임계치와 임계선으로부터의 $\sigma$-리프트 사이의 구조적 대응을 확립한다 (제~\ref{sec:sigma_lift}절).
  \item 정답 없이 자가 교정하는 다원군 ODE 영점 탐지기를 제시하며, 이것이 QROP-Net의 블라인드 테스트 시간 최적화를 반영함을 보인다 (제~\ref{sec:polyadic}절).
  \item 7개의 구조적 대응을 형식화하고, 유비와 증명 사이의 경계를 엄격하게 구분한다 (제3부, 제~\ref{sec:correspondence}절).
  \item \textbf{(음성 경험 결과, 2026-04-08)} 시클로토믹 PQO 프로그램을 $t \in [100, 200]$ 위에서 $L=4$로 경험적으로 시도하고, 세 시드 전체에서 영점/비영점 판별의 완전한 실패를 문서화하며, 그 원인을 \cref{eq:psi_n}의 \emph{$\Psip$-디커플링 결함}으로 추적하고, 향후 양성 결과의 필요 조건으로 감독 교체식 \cref{eq:psi_supervised}를 제안한다 (제~\ref{part:failed}부, Remark~\ref{rem:cos2_pqo}, \ref{rem:psi_defect}).
  \item \textbf{(고립 실험, 2026-04-09)} 타겟 주입, 손실 가중, 읽기 채널을 모두 분리하여, 세 가지를 동시에 차단했을 때 감독 사슬 $\arg\xif(\tfrac{1}{2}+it)\to\mathcal{L}_{\mathrm{tgt}}\to\hat\Psip$이 비영점 샘플링 지점에서 건전함을 보인다 ($p=1.18\times10^{-15}$, 비율 $18.2\times$). 이로써 잔존 결함은 영점 근방의 네트워크 출력 표현에 국한된다 (Remark~\ref{rem:isolation}).
  \item 다섯 개의 독립적 축(국소 곡률, 전역 스펙트럼 통계, 위상적 불변량,
        디리클레 $L$-함수 확장, 음성 대조군, GL(2) 확장, GL(3) 확장)에 걸친 62개의 수치 결과를 통해
        $\xif$-함수의 파이버 다발 해석을 검증한다.
        종합 요약은 표~\ref{tab:summary}에 수록된다.
        Stirling 기반 분석을 통해 \{모노드로미, $\kappa$-집중도, 블라인드 예측\}은
        \emph{차수-보편적(degree-universal)} 진단임을, $\sigma$-유일성은
        차수-1-특이적 성질임을 규명하며, GL(2) FAIL은
        정확한 항등식 $f_{\Gamma(s)} = f_\mathrm{zero}$로 설명된다
        (\cref{sec:sigma_uniqueness_mechanism}).
\end{enumerate}

\subsection*{독해 안내: 증거 계층}
\label{sec:tiers}

\noindent 본 논문은 \emph{통합 연구 노트}이다. 4개월간의 연구 궤적 전체를 기록하며, 여전히 추측 상태인 아이디어와 실패로 끝난 경험적 시도까지 포함한다. 따라서 모든 주장에 세 가지 계층 중 하나를 명시적으로 부여하고, 이를 섹션 제목 및 목차에 표시한다:

\begin{description}
  \item[\textbf{\textcolor{ForestGreen}{\rule{0.9ex}{0.9ex}}\ \ 계층 1 --- 검증됨.}]
        완전한 유도와 성공적인 수치 실험으로 뒷받침된 주장. 예:
        $F_2$-판별식 영점 탐지기 (제~\ref{sec:f2_hardy}절),
        고립 실험을 통한 감독 동일시
        $\Psip(t)=\arg\xif(\tfrac{1}{2}+it)$ (Remark~\ref{rem:isolation}),
        $t\in[100,200]$ 고위영점 확장 (제~\ref{sec:high_zeros}절).
        이들 섹션은 성숙한 결과로 읽어야 한다.

  \item[\textbf{\textcolor{orange}{\rule{0.9ex}{0.9ex}}\ \ 계층 2 --- 아이디어와 부분 결과.}]
        해석적 대응, 이론적 구성, 또는 시사적이지만 불완전한 증거를 가진 실험. 예: Gate B 정류성 (제~\ref{sec:gates}절),
        $\sigma$-리프트 (제~\ref{sec:sigma_lift}절), 7가지 구조적 대응 (제~\ref{sec:correspondence}절), $|F_2|$ 공간 관측량 및 2026-04-08 재해석 (제~\ref{sec:f2_residual}절). 이들은 연구 강령으로 읽어야 하며, 확정된 이론이 아니다.

  \item[\textbf{\textcolor{red}{\rule{0.9ex}{0.9ex}}\ \ 계층 3 --- 실패한 시도와 그 진단.}]
        성공하지 못한 경험적 프로그램을 완전한 진단과 함께 보존한 것. 예: 시클로토믹 $\cos^2$ $L\in\{2,4\}$ PQO 학습 (제~\ref{sec:cos2_failure}절), 위상-게이지 경사하강(PGGD) 옵티마이저 (제~\ref{sec:pggd_failure}절). 이들 섹션을 압축하되 유지하는 이유는 진단 사슬 자체가 기여이기 때문이다: 타겟 설계, 손실 가중, 측정, 모델 표현력의 어떤 조합이 호환되지 않는지 정확히 기록함으로써 후속 작업이 동일한 실패 모드를 피할 수 있게 한다. 원인이 이해된 실패 방향은 부분적 양성 결과이다.
\end{description}

\noindent 섹션이나 비고의 계층은 \textcolor{ForestGreen}{\rule{0.9ex}{0.9ex}}\,/\textcolor{orange}{\rule{0.9ex}{0.9ex}}\,/\textcolor{red}{\rule{0.9ex}{0.9ex}}\ 기호로 표시한다. 독자는 계층을 특정 진술에 대한 신뢰도의 눈금으로 사용할 수 있다. 제~\ref{part:qropnet}--\ref{part:synthesis}부의 구조적 분할은 이 계층화를 따른다: 제~\ref{part:qropnet}부와 제~\ref{part:resonance}부는 계층~1과 계층~2 자료를, 제~\ref{part:failed}부는 계층~3 실패 진단을, 제~\ref{part:synthesis}부는 세 가지 모두를 통합 그림으로 종합한다.


% ############################################################################
% 제1부 — 구체적 시스템
% ############################################################################

\part{구체적 시스템: QROP-Net}
\label{part:qropnet}


% ============================================================================
\section{QROP-Net: 시스템과 수학}
\label{sec:qropnet}
% ============================================================================

\subsection{다항식환과 NTT}

Kyber는 $R_q = \Rq$ 위에서 동작하며, $q = 7681$은 $q \equiv 1 \pmod{2n}$을 만족하도록 선택되어 $O(n\log n)$ 다항식 곱셈을 위한 수론적 변환(NTT)을 가능하게 한다. 환 조건 $X^n + 1 = 0$은 다음과 같이 인수분해된다:
\begin{equation}
  X^n + 1 = \prod_{j=0}^{n-1}(X - \omega_j), \quad
  \omega_j = e^{i\pi(2j+1)/n} \in \mathbb{C}.
  \label{eq:factorisation}
\end{equation}
모든 원소 $a(X) \in R_q$는 벡터 $\mathbf{a} = (a_0, \ldots, a_{n-1}) \in \Zq^n$에 대응하며, 곱셈 $a \cdot b \pmod{X^n+1}$는 NTT 영역에서 성분별 곱셈이 된다.

\begin{definition}[Kyber IND-CPA 암호화]
\label{def:kyber}
공개키 $(\mathbf{A}, \mathbf{t})$가 주어지고 $\mathbf{A} \in R_q^{k \times k}$, $\mathbf{t} \in R_q^k$, 메시지 $m \in \{0,1\}^n$일 때, 암호문 $(\mathbf{u}, v)$는:
\begin{align}
  \mathbf{u} &= \mathbf{A}^\top \mathbf{r} + \mathbf{e}_1,
  \label{eq:kyber_u}\\
  v &= \mathbf{t}^\top \mathbf{r} + e_2 + \ceil{q/2} \cdot m,
  \label{eq:kyber_v}
\end{align}
여기서 $\mathbf{r}, \mathbf{e}_1, e_2 \sim \beta_\eta$는 중심 이항 분포에서 추출된다:
\begin{equation}
  \beta_\eta: \quad x = \sum_{i=1}^{\eta}(a_i - b_i), \quad
  a_i, b_i \sim \mathrm{Bernoulli}(\tfrac{1}{2}), \quad
  x \in [-\eta, \eta].
  \label{eq:cbd}
\end{equation}
\end{definition}

\begin{definition}[압축과 역압축]
\label{def:compress}
\begin{align}
  \mathrm{Compress}_q(x, d) &= \ceil{\frac{2^d}{q} \cdot x}
  \bmod 2^d, \label{eq:compress}\\
  \mathrm{Decompress}_q(x, d) &= \ceil{\frac{q}{2^d} \cdot x},
  \label{eq:decompress}
\end{align}
여기서 $\ceil{\cdot}$는 반올림을 나타낸다: $\floor{x + 0.5}$.
매개변수: $d_u = 11$ ($\mathbf{u}$에 대해 2048 수준), $d_v = 3$ ($v$에 대해 8 수준).
\end{definition}

\begin{theorem}[LWE 정확성 기준]
\label{thm:lwe}
복호화가 성공하려면 총 오차 $w = v - \mathbf{s}^\top\mathbf{u} \pmod{q}$가 다음을 만족해야 한다:
\begin{equation}
  \norm{w - \ceil{q/2}\hat{m}}_\infty < \ceil{q/4} = 1920.
  \label{eq:lwe_threshold}
\end{equation}
단 하나의 계수라도 이 경계를 초과하면 완전한 복호화 실패가 발생한다 ($\mathrm{DFR} = 1$). 환 거리는:
\begin{equation}
  |e|_{\Zq} = \min(e, \; q - e), \quad e \in \{0, 1, \ldots, q-1\}.
  \label{eq:ring_distance}
\end{equation}
\end{theorem}


\subsection{위상 변조: 정수에서 광학으로}
\label{sec:phase_mod}

각 압축된 암호문 계수 $c_i \in \{0, \ldots, 2^{d_u}-1\}$는 위상 각도로 변환되어 단위 진폭 복소 파면으로 인코딩된다:
\begin{equation}
  \phi_i = \frac{2\pi \cdot c_i}{2^{d_u}}, \qquad
  U_0(x,y) = e^{i\phi(x,y)}, \quad |U_0| \equiv 1.
  \label{eq:phase_mod}
\end{equation}
$2^{d_u} = 2048$개의 이산 암호문 수준은 단위원 위 2048개의 등간격 위상에 대응하며, $\Delta\phi = 2\pi/2048$만큼 분리된다. 이는 대수적 구조의 직접적인 물리적 실현이다: 다항식환의 이산 계수가 SLM 위의 이산 광학 위상이 된다.

차원 제약 조건이 암호학적 영역과 광학적 영역을 연결한다:
\begin{equation}
  N = \sqrt{k \cdot n + n} \times M = \sqrt{1024} \times 8 = 256,
  \label{eq:dim_constraint}
\end{equation}
여기서 $M = 8$은 잡음 내성을 위한 매크로 픽셀 확장 계수이다.


\subsection{미분 가능 광학 채널: FrFT}
\label{sec:frft}

파면 $U_0$가 자유 공간을 통해 거리 $d$만큼 전파될 때, 회절 적분은 차수 $p$의 분수 푸리에 변환과 동치이다:
\begin{equation}
  p = \frac{2}{\pi}\arctan\!\left(\frac{\lambda d}{N \Delta x^2}\right),
  \qquad \alpha = \frac{p\pi}{2},
  \label{eq:frft_order}
\end{equation}
여기서 $\lambda = 500\,$nm은 파장, $\Delta x = 8\,\mu$m은 픽셀 피치이다. 전파 거리 $d$는 학습 가능한 매개변수로 선언되어 경사 기반 자동 초점을 가능하게 한다:
\begin{equation}
  d \;\xrightarrow{\;\partial\;}\;
  R = \frac{\lambda d}{N\Delta x^2} \;\to\;
  \alpha = \arctan(R) \;\to\;
  s_2 = \sqrt{1 + R^2}.
  \label{eq:learnable_d}
\end{equation}

IP-DFrFT는 처프--FFT--처프 분해를 통해 구현되어 복잡도를 $O(N^4)$에서 $O(N^2\log N)$으로 줄인다:
\begin{equation}
\boxed{
  U_d = \underbrace{(1 - i\cot\alpha)}_{\displaystyle A_\alpha^2}
  \cdot\;\mathcal{F}^{-1}\!\Big[
    \mathcal{F}\!\big[\underbrace{U_0 \cdot
    e^{i\pi(\cot\alpha - \csc\alpha)\,r^2}}_{\text{입력 처프}}\big]
    \;\cdot\;
    \mathcal{F}\!\big[\underbrace{e^{i\pi\,\csc\alpha\,r'^2}}_{\text{커널 처프}}\big]
  \Big] \cdot \Delta x^2
}
\label{eq:frft_full}
\end{equation}
여기서 $r^2 = (m^2+n^2)\Delta x^2$는 $N \times N$, $r'^2$는 순환 앨리어싱을 방지하기 위해 제로 패딩된 $2N \times 2N$ 위에 정의된다.

센서는 세기를 캡처하고 물리적 잡음을 추가한다:
\begin{align}
  I_{\mathrm{clean}} &= \frac{|U_d|^2}{1 + \cot^2\!\alpha},
  \label{eq:intensity}\\
  I_{\mathrm{shot}} &\sim \mathrm{Poisson}(I_{\mathrm{clean}} \cdot \mathrm{FWC}),
  \label{eq:poisson}\\
  I_{\mathrm{meas}} &= \mathrm{ADC}_{12\text{-bit}}
  (I_{\mathrm{shot}} + \mathcal{N}(0, \sigma_{\mathrm{read}}^2)).
  \label{eq:adc}
\end{align}
직통 추정기(STE)는 미분 불가능한 포아송 표본 추출과 ADC 반올림을 통과하는 경사 흐름을 보존한다:
\begin{equation}
  \STE(x_{\text{noisy}}, x_{\text{clean}}) =
  x_{\text{clean}} + \mathrm{stopgrad}(x_{\text{noisy}} - x_{\text{clean}}).
  \label{eq:ste}
\end{equation}


\subsection{혼합 손실 오케스트레이션}
\label{sec:loss}

네 가지 손실 항이 교과 일정 스케줄러 하에서 결합된다.

\subsubsection{데이터 충실도 손실}
예측된 파면 $\hat{U}_0$는 FrFT를 통해 전파되어 $\hat{I}$를 생성한다. 진폭 영역에서의 비교가 균일한 경사 흐름을 보장한다:
\begin{equation}
  \mathcal{L}_{\mathrm{fid}} = \frac{1}{N^2}\sum_{x,y}
  \left(\sqrt{\hat{I}(x,y) + \epsilon} -
  \sqrt{I_{\mathrm{target}}(x,y) + \epsilon}\right)^{\!2}.
  \label{eq:loss_fid}
\end{equation}

\subsubsection{주기적 연성 양자화 손실}
$\sin^2$ 기반 벌칙이 연속 위상을 이산 Kyber 격자 수준으로 강제한다:
\begin{equation}
\boxed{
  \mathcal{L}_{\mathrm{quant}} = \frac{1}{2}\!\left[
    \frac{1}{kn}\sum_{i=1}^{kn}
    \sin^2\!\Big(\frac{2^{d_u}}{2}\,\phi_i^{(u)}\Big) +
    \frac{1}{n}\sum_{j=1}^{n}
    \sin^2\!\Big(\frac{2^{d_v}}{2}\,\phi_j^{(v)}\Big)
  \right]
}
\label{eq:loss_quant}
\end{equation}
$\sin^2$ 함수는 \textit{골짜기--능선 풍경}을 생성한다: 정확한 격자 위치 $\phi = 2\pi k / 2^d$에서 영 손실, 중점에서 최대 벌칙. 이것이 해석적 정수론으로 가교를 놓는 핵심 방정식이다 (제~\ref{sec:bridge}절).

\subsubsection{$L^\infty$ 지수 장벽 손실}
지수 장벽이 가장 위험한 이상값에 대해 LWE 임계치를 강제한다:
\begin{equation}
  \mathcal{L}_{\max} = \frac{1}{|S_K|}\sum_{i \in S_K}
  \exp\!\big(\gamma \cdot \mathrm{ReLU}(|e_i|_{\Zq} -
  \tau_{\mathrm{safe}})\big),
  \label{eq:loss_max}
\end{equation}
여기서 $\tau_{\mathrm{safe}} = 0.8 \times 1920 = 1536$, $\gamma = 10$, $S_K$는 상위 5\% 오차 픽셀 집합이다.

\subsubsection{총 변이 손실}
복소 파면에 대해 (래핑 아티팩트를 피하기 위해 위상 각도가 아님):
\begin{equation}
  \mathcal{L}_{\mathrm{TV}} = \frac{1}{N^2}\sum_{x,y}
  \big(|\hat{U}_0^{x+1,y} - \hat{U}_0^{x,y}|^2 +
  |\hat{U}_0^{x,y+1} - \hat{U}_0^{x,y}|^2\big).
  \label{eq:loss_tv}
\end{equation}


\subsection{수신기 신경망과 복호화}
\label{sec:rx}

LeWin U-Net (윈도우 기반 다중 헤드 자기 주의 + 지역 강화 피드포워드)이 세기 측정값 $I_{\mathrm{meas}}$를 2-채널 $(\mathrm{Re}, \mathrm{Im})$ 출력으로 변환하며, 위상은 다음을 통해 추출된다:
\begin{align}
  \hat{\phi} &= \mathrm{atan2}(\mathrm{Im},\; \mathrm{Re}) \bmod 2\pi,
  \label{eq:phase_extract}\\
  \hat{U}_0 &= e^{i\hat{\phi}}, \quad |\hat{U}_0| \equiv 1.
  \label{eq:wavefront}
\end{align}

매크로 픽셀 풀링은 위상 평균화 재앙을 피하기 위해 원형 통계를 사용한다:
\begin{equation}
  \bar{\phi} = \mathrm{atan2}\!\left(
    \frac{1}{M^2}\sum \sin\phi_{ij},\;
    \frac{1}{M^2}\sum \cos\phi_{ij}
  \right) \bmod 2\pi.
  \label{eq:circular_mean}
\end{equation}

최종 Kyber 역캡슐화:
\begin{align}
  \hat{\mathbf{u}}, \hat{v} &= \mathrm{Decompress}_q(c_{\mathrm{pred}},
  d_u, d_v), \label{eq:decompress_final}\\
  w &= \hat{v} - \mathbf{s}^\top\hat{\mathbf{u}} \pmod{q},
  \label{eq:inner_product}\\
  \hat{m}_j &= \ceil{\frac{2}{q} \cdot w_j} \bmod 2.
  \label{eq:decode_msg}
\end{align}


\subsection{2단계 훈련과 물리층 보안}
\label{sec:training}

\subsubsection{1단계: 상각 사전훈련}
합성 암호문이 즉석에서 생성되며; AdamW가 119개 U-Net 매개변수와 물리 거리 $d$를 동시에 갱신한다.

\subsubsection{2단계: 블라인드 테스트 시간 최적화 (TTO)}
배치 시, 물리 기반 및 수학 기반 제약만 사용 가능하다:
\begin{equation}
  \mathcal{L}_{\text{blind}} =
  \mathcal{L}_{\text{fid}} + 0.5\,\mathcal{L}_{\text{quant}} +
  0.05\,\mathcal{L}_{\text{TV}}.
  \label{eq:blind_loss}
\end{equation}
U-Net 백본은 $\mathrm{lr} = 10^{-5}$로 동결되고 $d$는 $\mathrm{lr} = 10^{-3}$을 받아, \textit{정답 암호문에 접근하지 않고} 단 15회 경사 단계만으로 빠른 자동 초점을 가능하게 한다.

\begin{theorem}[눈사태 효과 --- 물리층 보안]
\label{thm:avalanche}
동일한 U-Net 가중치, PRNG 시드, Kyber 비밀키를 보유하지만 $d' = d + 0.1\,\mathrm{mm}$을 추정하는 도청자는 $\mathrm{DFR} = 1$을 경험한다. FrFT 각도 $\alpha$는 $d$의 비선형 함수이며; 0.1\,mm 오프셋은 처프--FFT--처프 파이프라인을 통해 연쇄 전파되어 모든 암호문 원소에 대해 $\norm{e}_\infty > 1920$을 초래한다.
\end{theorem}


% ============================================================================
\section{위상 양자화 원리}
\label{sec:bridge}
% ============================================================================

이제 $\sin^2$ 양자화 메커니즘 (\cref{eq:loss_quant})을 일반 원리로 추상화한다.

\begin{definition}[위상 양자화 연산자]
\label{def:pqo}
$\phi \in [0, 2\pi)$를 연속 위상 변수, $L \in \mathbb{N}$을 목표 이산 수준 수라 하자. \textit{위상 양자화 벌칙}은:
\begin{equation}
  \mathcal{P}(\phi; L) = \sin^2\!\left(\frac{L}{2}\,\phi\right).
  \label{eq:pqo}
\end{equation}
이 함수는 다음 성질을 갖는다:
\begin{enumerate}
  \item $\mathcal{P}(\phi; L) = 0$ $\iff$ $\phi \in \{2\pi k/L : k = 0, 1, \ldots, L-1\}$ (격자 위치).
  \item $\mathcal{P}(\phi; L) = 1$ (중점 $\phi = 2\pi(k + \tfrac{1}{2})/L$에서).
  \item $\nabla_\phi \mathcal{P} = \tfrac{L}{2}\sin(L\phi)$은 $\phi$를 가장 가까운 격자 점으로 조향하는 경사를 제공한다.
  \item 풍경은 $2\pi/L$-주기이며, 단위원 위에 $L$개의 등간격 흡인 유역을 생성한다.
\end{enumerate}
\end{definition}

QROP-Net에서 격자 수준은 $L = 2^{d_u} = 2048$ ($\mathbf{u}$-계수용)과 $L = 2^{d_v} = 8$ ($v$-계수용)이다. $\sin^2$ 손실은 각 연속 위상 예측을 가장 가까운 이산 Kyber 수준으로 유도하며, 완전 미분 가능한 \textit{연성 양자화}를 실현한다.

\begin{remark}[시클로토믹 연결]
\label{rem:cyclotomic}
$\mathcal{P}(\phi; L)$의 $L$개 영점은 정확히 $L$차 단위근 $e^{2\pi i k/L}$이다. $L = 2n = 512$ (시클로토믹 다항식 $X^n + 1$의 차수)일 때, 이것은 \cref{eq:cyclo_roots}의 근 $\omega_0, \ldots, \omega_{n-1}$이다. $\sin^2$ 손실은 따라서 대수적 조건 $X^n + 1 = 0$의 위상 영역에서의 \textit{해석적 재진술}이다.
\end{remark}

\begin{remark}[코사인 PQO와 고분해능 시클로토믹 격자]
\label{rem:cos2_pqo}
PQO는 구조적으로 동일한 \emph{코사인} 변형을 허용한다:
\begin{equation}
  \mathcal{P}_{\cos}(\phi; L) = \cos^2\!\left(\frac{L}{2}\,\phi\right),
  \label{eq:pqo_cos}
\end{equation}
그 영점 집합은 $\phi \in \{(2k+1)\pi/L : k = 0, 1, \ldots, L-1\}$이다. 이들은 정확히 \emph{원시 $2L$차 단위근}, 즉 \cref{eq:cyclo_roots}에서 $X^L + 1 = 0$의 $L$개 근 $\omega_j$이다. 따라서 $\mathcal{P}_{\cos}(\,\cdot\,; L)$는 $L$차 시클로토믹 PQO이며, \cref{rem:cyclotomic}와 평행하되 대수적 제약 $X^L + 1 = 0$을 위상 영역에 \emph{직접} 실현한다 ($\sin^2$가 매칭하는 이동된 단위근 제약 $X^L - 1 = 0$과 대조).

\smallskip\noindent\textbf{$L = 2$에서의 자명한 붕괴.} $L = 2$일 때 끌개 집합은 $\{i, -i\}$로 퇴화한다. $\mathcal{P}_{\cos}(\phi; 2)$ 항을 비자명한 가중치로 포함한 손실로 학습된 임의의 신호는 전역적으로 그 두 점 중 하나로 붕괴한다: 특히 임계선 판별 네트워크의 손실이 해당 항을 포함하면, 경험적 위상 $\arg Z(t)$는 \emph{학습된 대역의 모든 $t$에 대해} $\pi/2$로 끌려가며, $t$가 리만 영점 근처인지 여부와 무관하다. 따라서 $L = 2$ 코사인 PQO는 영점 국소화 정보를 전혀 운반하지 않는다; 단지 끌개 메커니즘이 설계대로 작동함을 확인할 뿐이다.

\smallskip\noindent\textbf{고분해능 일반화 (시도).} $L \geq 4$에서 격자는, 원리상, $S^1$ 위의 섹터 분류기로 작용할 충분한 계수를 획득한다. 각 후보 종좌표 $t$에 다음의 이산 라벨이 부여된다:
\begin{equation}
  k(t) \;=\; \arg\min_{0 \leq k < L}\,\big|\arg Z(t) - (2k{+}1)\pi/L\big|
  \pmod{2\pi},
  \label{eq:cyclo_label}
\end{equation}
학습된 위상에 가장 가까운 시클로토믹 섹터 $\omega_{k(t)}$를 명명하는 것으로, 리만 영점이 비자명한 부분 격자를 점유하여 $L=2$가 허용하는 것보다 더 미세한 산술 층화를 제공하기를 기대한 것이다.

\smallskip\noindent\textbf{$L = 4$에서의 경험적 반증 (2026-04-08).} 이 프로그램에 대한 직접 검정이 수행되었다: 임계선 판별 네트워크, $L=4$, 은닉 폭 $32$, 시드 $\{42, 123, 777\}$ 3개, 학습 에폭 $200$, 라벨된 모집단은 리만 영점 $50$개와 간격 $\geq 0.5$인 비영점 종좌표 $200$개 (학습 스크립트는 동반 \textsc{gdl\_unified} 저장소의 \texttt{scripts/phase\_L4\_experiment.py}, 출력은 \texttt{outputs/analysis/phase\_L4\_experiment.\{json,txt\}}). 결과는 세 독립 축에서 모호함 없이 음성이다:
\begin{enumerate}
  \item 세 시드 모두 영점/비영점 $k$-분포 독립성 $\chi^2$ 검정 통과 ($p = 0.127, 0.528, 0.410$; 풀링 $p = 0.154$).
  \item 네 시클로토믹 섹터 $\{k=0,1,2,3\}$ 중 $k=0,1$만 점유되며 $k=2,3$은 모든 시드에서 \emph{정확히 0개} 샘플 (비대칭 basin 붕괴).
  \item 영점/비영점의 최근접 끌개 잔차 평균은 모든 시드에서 $\sim 5\%$ 이내로 일치.
\end{enumerate}
$L=2$의 자명한 붕괴와 $L=4$의 비대칭 붕괴는 따라서 단일 근본 원인을 공유한다; 이제 이를 식별한다.

\smallskip\noindent\textbf{진단: $\Psip$-디커플링 결함.} 판별 네트워크의 학습 손실은
$\mathcal{L} = \lambda_{\mathrm{res}}\mathcal{L}_{\mathrm{res}} +
        \lambda_{\mathrm{curv}}\mathcal{L}_{\mathrm{curv}} +
        \lambda_{\mathrm{tgt}}\mathcal{L}_{\mathrm{tgt}} +
        \lambda_{\mathrm{pqo}}\mathcal{L}_{\mathrm{pqo}}$
의 합이며, 타겟 매칭 항 $\mathcal{L}_{\mathrm{tgt}}$는 \cref{def:psi}\,/\,\cref{eq:psi_n}의 $\Psip$를 기준으로 구성된다. 그러나 $\Psip$는 $\|X\|$에만 의존하는, 합성으로 선택된 $(\alpha,\beta,\gamma,\delta)$ 곡선이며 \emph{$\xif$-영점 정보를 전혀 운반하지 않는다}. 결과적으로:
\begin{itemize}
  \item $\mathcal{L}$의 어느 항도 $t$-국소 영점 라벨을 학습에 주입하지 않는다;
  \item 코사인 PQO 항 $\mathcal{L}_{\mathrm{pqo}}$는 $t$-\emph{무관} 포텐셜이며, basin 선택은 $t$가 영점에 대해 어디에 있는지가 아니라 초기화에 의해 고정된다;
  \item 라벨링 규칙 \cref{eq:cyclo_label}는 따라서 영점이 무엇인지 한 번도 들은 적이 없는 네트워크 위에 얹힌 \emph{운영적으로 분리된 측정}이다.
\end{itemize}
특히 \cref{eq:pqo_cos} (학습)와 \cref{eq:cyclo_label} (측정)는 호환되지 않는 신호 범주에 작용한다: 전자는 후자가 읽어내려는 $t$-국소 정보를 \emph{지운다}. $L$을 $2$에서 $4, 8, \ldots$로 늘려도 이 결함은 복구되지 않는다 — 끊어진 고리는 $L$의 상류에 있다.

\smallskip\noindent\textbf{시클로토믹 PQO 프로그램의 상태.} 이전에 정식화된 $L \geq 4$ 프로그램은 보류된다. 부활시키려면 \cref{eq:psi_n}을 진정한 영점 정보를 운반하는 타겟으로 \emph{교체}해야 한다 (자연스러운 후보는 $\Psip(t) := \arg \xif(\tfrac{1}{2}+it)$ 또는 \cref{eq:hardy_z}의 Hardy 위상 $\arg Z(t)$), $\mathcal{P}_{\cos}$를 \emph{학습} 포텐셜에서 결과 위상에 대한 순수 \emph{측정}으로 강등시키며, 새로운 감독 하에 $L=4$ 프로토콜을 재실행해야 한다. 그러한 실험이 존재하기 전까지, 본 연구가 리만 영점의 시클로토믹 층화에 대해 어떤 주장도 지지하지 않는다; 데이터 파일 \texttt{phase\_L4\_experiment.*} 및 \texttt{phase\_pi\_half\_verify.txt}는 결정적 음성 통제로 보존되며, 위 진단의 근거 기록의 일부이다.
\end{remark}

\begin{remark}[보어--좀머펠트 유비]
\label{rem:bohr_sommerfeld}
조건 $\mathcal{P}(\phi; L) = 0$은 $L\phi/2 = m\pi$ ($m \in \mathbb{Z}$)와 동치이다, 즉
\begin{equation}
  \Delta\nu \;=\; m\pi, \qquad m \in \mathbb{Z},
  \label{eq:bohr_sommerfeld}
\end{equation}
여기서 $\Delta\nu = L\phi/2$는 재조정된 위상 축적이다. 이는 반고전 역학의 \textit{보어--좀머펠트 양자화 조건}과 정확히 같은 형태이며, 닫힌 궤도 주위의 작용 적분이 $\pi\hbar$의 정수배여야 한다 (적절한 단위에서). 광학에서 같은 조건이 보강 간섭을 결정한다: 경로 길이가 반파장의 정수배를 수용해야 한다. 위상 양자화 원리는 따라서 다음을 통합한다:
\begin{center}
\begin{tabular}{lll}
\toprule
\textbf{영역} & \textbf{양자화되는 변수} & \textbf{조건} \\
\midrule
QROP-Net       & 암호문 위상 $\phi_i$
                & $\sin^2(2^{d_u-1}\phi_i) = 0$ \\
반고전 양자역학 & 작용 적분 $\oint p\,dq$
                & $= (n + \tfrac{1}{2})\pi\hbar$ \\
광학 (공진기)   & 광학 경로 길이 $nL_{\mathrm{cav}}$
                & $= m\lambda/2$ \\
임계선          & $\xif(s)$의 위상
                & $\Delta\nu = m\pi$ \\
\bottomrule
\end{tabular}
\end{center}
\end{remark}


% [Paper A: ξ-다발 기하학]
% ############################################################################
% 제2부 — 해석적 확장
% ############################################################################

\part{해석적 확장: 공명 이론}
\label{part:resonance}

\vspace{0.5em}
\noindent\textit{이 부분의 모든 수학적 주장은 구체적인 QROP-Net 대응(비고 환경에서 표기)이 있거나 문헌에 독립적인 근거가 있다. 두 조건 모두 충족하지 못하는 공명소수법칙 초기 버전의 주장은 생략되었다.}
\vspace{0.5em}


% ============================================================================
\section{완성 제타 함수의 위상 구조}
\label{sec:xi_phase}
% ============================================================================

\subsection{리만 $\xif$-함수}

\begin{definition}[완성 제타 함수]
\label{def:xi}
리만 $\xif$-함수는 다음과 같이 정의된다:
\begin{equation}
  \xif(s) = \frac{1}{2}s(s-1)\pi^{-s/2}\Gamma\!\left(\frac{s}{2}\right)\zeta(s),
  \label{eq:xi_def}
\end{equation}
여기서 $\zeta(s) = \sum_{n=1}^{\infty}n^{-s}$는 리만 제타 함수, $\Gamma$는 오일러 감마 함수이다. 함수 $\xif$는 전해석적(극이 없음)이며 함수 방정식을 만족한다:
\begin{equation}
  \xif(s) = \xif(1-s).
  \label{eq:func_eq}
\end{equation}
임계선 $s = \tfrac{1}{2} + it$ 위에서 $\xif$는 실수값을 가진다: $\xif(\tfrac{1}{2}+it) \in \mathbb{R}$.
\end{definition}

\begin{definition}[바이어슈트라스 곱 표현]
\label{def:weierstrass}
$\xif(s)$가 전해석적이므로, 그 영점 $\rho$ ($\zeta$의 비자명 영점)에 대한 바이어슈트라스 곱을 허용한다:
\begin{equation}
  \xif(s) = e^{g(s)}\prod_{\rho}
  \left(1 - \frac{s}{\rho}\right),
  \label{eq:weierstrass}
\end{equation}
여기서 $g(s)$는 수렴을 보장하는 전해석 함수이다. $\ln\xif$를 미분하면 \textit{로그 도함수}를 얻는다:
\begin{equation}
  \frac{\xif'(s)}{\xif(s)} = g'(s) + \sum_{\rho}
  \frac{1}{s - \rho}.
  \label{eq:log_deriv}
\end{equation}
영점에서 떨어진 곳에서 $g'(s)$는 매끄러우며, 지배적 기여는 $s = \rho$ (즉 $\zeta$의 비자명 영점)에서의 극에서 온다.
\end{definition}


\subsection{임계선 위의 진폭--위상 분해}
\label{sec:amp_phase}

일반적인 $s = \sigma + it$ ($\sigma = \tfrac{1}{2}$로 제한되지 않음)에 대해 $\xif$를 진폭과 위상으로 분해한다:
\begin{align}
  \Ampl(\sigma, t) &= |\xif(\sigma + it)|,
  \label{eq:amplitude}\\
  v(\sigma, t) &= \arg\,\xif(\sigma + it),
  \label{eq:phase_general}
\end{align}
따라서 $\xif(\sigma + it) = \Ampl(\sigma,t)\,e^{iv(\sigma,t)}$이다.
$\log\xif = u + iv$ ($u = \log|\xif|$, $v = \arg\xif$)로 쓰면, 코시--리만 방정식이 다음을 준다:
\begin{equation}
  v_\sigma = \mathrm{Re}\!\left(\frac{\xif'}{\xif}\right), \quad
  v_t = -\mathrm{Im}\!\left(\frac{\xif'}{\xif}\right), \quad
  u_t = -v_\sigma, \quad u_\sigma = v_t.
  \label{eq:cauchy_riemann}
\end{equation}

임계선 $\sigma = \tfrac{1}{2}$ 위에서:
\begin{align}
  \Phiz(t) &\;:=\; v(\tfrac{1}{2}, t) = \arg\,\xif(\tfrac{1}{2}+it),
  \label{eq:phase_crit}\\
  \Ampl(t) &\;:=\; \Ampl(\tfrac{1}{2}, t) = |\xif(\tfrac{1}{2}+it)|.
  \label{eq:amp_crit}
\end{align}
$\xif(\tfrac{1}{2}+it) \in \mathbb{R}$이므로, 위상 $\Phiz(t)$는 $\{0, \pi\}$ 값만 가지며 ($2\pi$ 모듈러), 각 영점 교차에서 $\pi$만큼 도약한다.

\begin{remark}[QROP-Net 대응: 진폭--위상 분리]
\label{rem:amp_phase}
QROP-Net에서 송신 파면 $U_0 = e^{i\phi}$는 단위 진폭($|U_0| \equiv 1$)과 이산 위상 $\phi$를 갖는다. 수신기는 세기 $I = |U_d|^2$로부터 $\phi$를 복원해야 하며, 이는 회절이 초래한 진폭--위상 혼합을 ``역전''시키는 것을 포함한다. 임계선 위의 구조 $\xif = \Ampl\,e^{iv}$는 동일한 수학적 도전을 제시한다: 진폭 포락선에서 기저 위상 구조를 분리하는 것.
\end{remark}


\subsection{이상적 위상 함수 $\Psip(n)$}

\noindent\textit{독자 경고.} 아래 \cref{eq:psi_n}의 합성 형태는 2026-04-08 감독 타겟으로서 경험적으로 반증되었다; 정의 직후의 \cref{rem:psi_defect}와 진단에 대한 \cref{rem:cos2_pqo}을 참조하라. 정의는 역사적 서술을 위해 보존된다; 이후 개정에서 사용되는 운영적 교체식은 \cref{eq:psi_supervised}이다.

\begin{definition}[매끄러운 소수 계수 위상]
\label{def:psi}
정수 $n \geq 2$에 대해 이상적 위상 함수는 다음과 같이 정의된다:
\begin{equation}
  \Psip(n) = \alpha\log n +
  \frac{\beta\sin(\gamma\log n)}{n^\delta},
  \label{eq:psi_n}
\end{equation}
여기서 $\alpha, \beta, \gamma, \delta$는 구조 상수이다:
\begin{itemize}
  \item $\alpha > 0$은 대수 상승(정렬 기울기)을 결정한다.
  \item $\beta \geq 0$은 진동 진폭(공명 강도)을 제어한다.
  \item $\gamma > 0$은 위상 회전 주파수(국소 감마 인자와 관련)이다.
  \item $\delta > 0$은 지수 감쇠를 제어한다: 진동 항은 $n^{-\delta}$로 감쇠하여, 큰 $n$에서 $\Psip(n) \to \alpha\log n$이 된다.
\end{itemize}
\end{definition}

단조 항 $\alpha\log n$은 영점 계수 함수에 대한 리만--폰 만골트 공식 $N(T) \sim \frac{T}{2\pi}\log\frac{T}{2\pi e}$를 반영하며, 진동 보정 $\beta\sin(\gamma\log n)/n^\delta$는 소수 간격에 관련된 국소 요동에 대한 휴리스틱 대용으로 의도된 것이다.

\begin{remark}[\cref{eq:psi_n}의 디커플링 결함; 2026-04-08 기록]
\label{rem:psi_defect}
\cref{eq:psi_n}의 형태는 크기 $n$에만(혹은 네트워크 구현에서는 $\|X\|$에만) 의존하는 합성 $(\alpha,\beta,\gamma,\delta)$ 곡선이며, 리만 영점의 위치에 대한 정보를 \emph{전혀} 운반하지 않는다. 이를 판별 네트워크의 타겟으로 사용할 경우 — \cref{rem:cos2_pqo}의 $L=4$ 시클로토믹 PQO 실험에서처럼 — 네 손실 항 $\mathcal{L}_{\mathrm{res}},\mathcal{L}_{\mathrm{curv}},\mathcal{L}_{\mathrm{tgt}},\mathcal{L}_{\mathrm{pqo}}$의 어느 것도 $t$-국소 영점 신호를 학습에 주입하지 않으며, 결과 네트워크는 $L$과 무관하게 영점/비영점 판별에 반드시 실패한다. 이 결함은 매개변수적이 아니라 구조적이다: $(\alpha,\beta,\gamma,\delta)$의 어떤 선택으로도 복구되지 않는다.

\smallskip
이후 개정에서 사용되는 의도된 교체는 임계선 위의 이상적 위상을 직접 재정의하는 것이다:
\begin{equation}
  \Psip(t) \;:=\; \arg \xif(\tfrac{1}{2}+it),
  \label{eq:psi_supervised}
\end{equation}
또는 동등하게 \cref{eq:hardy_z}의 Hardy $Z$-함수를 통해 $\Psip(t) := \arg Z(t)$로 잡는다. 어느 선택이든 \emph{구성에 의해} $\xif$-영점 정보를 공급하며, 감독 타겟과 시클로토믹 측정 \cref{eq:cyclo_label} 사이의 인과 사슬을 복원한다. 합성 형태 \cref{eq:psi_n}는 본문에 그 실패가 진단을 동기화시킨 역사적 객체로서만 보존된다.

\smallskip\noindent\textbf{감독-$\Psip$ 제어 실험 (2026-04-08).} 교체식 \cref{eq:psi_supervised}를 직접 시험하였다. 동일한 판별 네트워크 구조 (은닉폭 $32$, 3층, 세 시드 $\{42,123,777\}$, $200$ 에폭, $t \in [100, 200]$, 영점 $50$개 대 간격 $\geq 0.5$의 비영점 $200$개)에 대해 $L=4$ 기준선 대비 두 가지 한 줄 수정을 가하였다: (i) 모델의 $\Psip$ 출력을 미리 계산한 $\arg\xif(\tfrac{1}{2}+it)$로 덮어쓰기하여 감독 신호로 공급, (ii) $\lambda_{\mathrm{pqo}}=0$으로 코사인 PQO를 학습에서 제거하고 오직 시클로토믹 측정 \cref{eq:cyclo_label}에서만 유지 (스크립트 \texttt{scripts/phase\_supervised\_experiment.py}, 출력 \texttt{phase\_supervised\_experiment.\{json,txt\}}). 결과는 헤드라인 지표 상 다시 음성이다 --- 시드 풀링된 $\arg Z(t)$ 대 $\arg \xif(\tfrac{1}{2}+it)$의 위상 MSE가 영점과 비영점에서 사실상 동일하다 (평균 $2.256$ 대 $2.247$, 비율 $1.004$, Welch $p = 0.976$) --- 그리고 시드별 Welch 검정은 방향이 불일치한다 (한 시드에서 $p=0.034$로 기대 순서와 \emph{반대}). 그러나 두 가지 구조적 징후가 새롭고 유익하다:
\begin{enumerate}
  \item $L=4$ 기준선의 비대칭 분지 붕괴가 해제된다: 네 시클로토믹 구간 $k\in\{0,1,2,3\}$ 모두 모든 시드에서 점유된다. $\mathcal{L}_{\mathrm{pqo}}$에 의해 유도된 $t$-독립 어트랙터 고정이 삭제 실험을 통해 그 붕괴의 두 기여 요인 중 하나로 확인된 것이다.
  \item 기록된 최적 검증 손실 ($\approx 1.5$--$1.7$)이 기록된 위상 MSE ($\approx 2.25$)보다 상당히 작으며, 이는 $\mathcal{L}_{\mathrm{pqo}}$가 영으로 된 이후에도 $\mathcal{L}_{\mathrm{tgt}}$ 기여가 $\mathcal{L}$의 작은 부분임을 함의한다. 남은 $\mathcal{L}_{\mathrm{res}}$와 $\mathcal{L}_{\mathrm{curv}}$ 항이 그래디언트를 지배하므로, 네트워크는 실제로는 $\arg \xif$를 적합하라고 요구받지 않는다.
\end{enumerate}
이것은 위의 진단을 \emph{철회하지} 않고 정교화한다: $\Psip$-디커플링은 \cref{eq:psi_n}의 필요 결함이며, \cref{eq:psi_supervised}는 타겟 수준에서만 그것을 복구한다; 그러나 손실 가중 불균형과 의심되는 측정 아티팩트 (채널별 $\arg Z$의 산술 평균화, 이는 \cref{eq:phase_crit}가 예측하는 $\pm\pi$ 랩 신호를 파괴한다)는 \emph{추가적인} 끊긴 연결이다. $\lambda_{\mathrm{res}}=\lambda_{\mathrm{curv}}=\lambda_{\mathrm{pqo}}=0$ (오직 $\mathcal{L}_{\mathrm{tgt}}$만 활성) 및 원형 평균 읽기를 사용하는 두 번째 제어 실험은 아래 \cref{rem:isolation}에서 서술한다.
\end{remark}

\begin{remark}[고립 실험: 감독 사슬은 건강, $\pm\pi$ 점프 표현력은 결함; 2026-04-09 기록]
\label{rem:isolation}
\cref{rem:psi_defect}의 세 의심 결함 --- (i) \cref{eq:psi_n}를 통한 $\Psip$-디커플링, (ii) 손실 가중 불균형, (iii) $\pm\pi$ 랩을 표현할 수 없는 산술 평균 읽기 --- 를 분리하기 위해, 세 채널을 동시에 모두 차단한 최대 고립 제어를 실행하였다. 네트워크는 $\lambda_{\mathrm{res}}=\lambda_{\mathrm{curv}}=\lambda_{\mathrm{pqo}}=0$, 오직 $\mathcal{L}_{\mathrm{tgt}}$만 활성으로 재학습하였고, 감독 타겟은 \cref{eq:psi_supervised}와 같이 $\arg\xif(\tfrac{1}{2}+it)$로 하드 설정하였으며, 샘플별 위상은 채널별 산술 평균이 아니라 \emph{원형} 평균 $\arg\sum_c e^{i\hat\Psip_c}$로 산출하였다 (스크립트 \texttt{scripts/phase\_isolation\_experiment.py}, 출력 \texttt{phase\_isolation\_experiment.\{json,txt\}}; \cref{rem:psi_defect,rem:cos2_pqo}와 동일한 구조, 시드, 에폭, $t$-구간).
결과는 강하게 \emph{양성}이다:
\begin{itemize}
  \item \textbf{원형 읽기.} $\arg\xif(\tfrac{1}{2}+it)$에 대한 시드 풀링 위상 MSE가 $\xif$-영점 집합에서 $2.1708$, 비영점 집합에서 $0.1195$이며, 비율 $18.16\times$, Welch $p=1.18\times10^{-15}$이다. 시드별 비율은 각각 $29.5,\,31.9,\,9.5$이며 $p\in\{4.5,3.7,5.5\}\times10^{-6}$이다.
  \item \textbf{산술 읽기, 동일 가중.} 동일한 학습 실행을 기존 산술 평균으로 읽어내면 영점/비영점 평균이 $2.410$ 대 $1.849$, 비율 $1.30$, 풀링 $p=0.046$로 붕괴된다. 따라서 원형 대 산술의 차이는 학습 차이가 아니라 측정 차이이다.
\end{itemize}
\cref{rem:psi_defect}의 두 가설이 모두 확인된다: 손실 가중 불균형 (ii)와 산술 평균 읽기 (iii)는 각각 신호를 은폐하기에 충분하며, 둘이 합쳐져 \cref{rem:cos2_pqo}와 첫 감독 제어의 음성 결과를 완전히 설명한다. 둘이 모두 차단되면 감독 사슬 $\arg\xif(\tfrac{1}{2}+it)\mapsto\mathcal{L}_{\mathrm{tgt}}\mapsto\hat\Psip$은 \emph{건강}하다: 비영점 MSE $0.12$는 네트워크가 매끄러운 가지 위에서 $\arg\xif$를 사실상 완벽하게 학습함을 보인다.

\smallskip
잔존하는 영점 MSE $2.17\approx(\pi/2)^2$는 깔끔한 해석을 가진다. 임계선 위에서 $\xif$는 실수이며, 수치적으로 우리 캐시에서 $|\mathrm{Im}\,\xi|/|\mathrm{Re}\,\xi|\sim10^{-53}$이므로 $\arg\xif$는 사실상 모든 리만 영점에서 부호를 전환하는 $\{0,\pi\}$-값 계단 함수이다. 네트워크 출력 $\hat\Psip(t)$는 매끄러운 $\mathbb{R}$-값 장(場)이며 그러한 순간 $\pm\pi$ 점프를 표현할 수 없다; 영점 근방에서 최적의 원형 평균 적합은 상수를 제외하면 중간값이며, 정확히 $(\pi/2)^2$을 MSE에 기여한다. 따라서 남은 결함은 타겟에도 읽기에도 있지 않고 판별 네트워크의 \emph{출력 표현}에 있으며, \cref{rem:cos2_pqo,rem:psi_defect}에서는 $L$을 올려서만 간접 대응된다; 올바른 해결은 매끄러운 영점 교차 타겟 (예컨대 $\mathrm{sign}(Z(t))\cdot e^{-|Z(t)|^2/\sigma^2}$) 또는 $\xif(\tfrac{1}{2}+it)$를 복소 벡터로 예측하고 사후에 $\arg$를 읽는 헤드를 요구한다.

\smallskip
\noindent\textbf{매끄러운 영점 교차 확인 (2026-04-10).}
동일한 고립 설정 ($\lambda_{\mathrm{res}}=\lambda_{\mathrm{curv}}=\lambda_{\mathrm{pqo}}=0$, \texttt{in\_features}$=128$, 3 시드) 에서, $\pm\pi$ 불연속 없이 동일한 부호 정보를 부호화하는 두 가지 매끄러운 대리 타겟을 시험하였다:
\begin{itemize}
  \item \textbf{범프 타겟}
  $\Psi_{\mathrm{bump}}(t) = \mathrm{sign}(\mathrm{Re}\,\xif)
   \cdot e^{-(\mathrm{Re}\,\xif)^2/\sigma^2}$,
  $\sigma=\mathrm{median}(|\mathrm{Re}\,\xif|)$:
  영점 MSE $=0.935$, 비영점 MSE $=0.693$, 비율 $1.35\times$
  (갭 축소 $92.6\%$).
  \item \textbf{정규화 타겟}
  $\Psi_{\mathrm{norm}}(t) = \mathrm{Re}\,\xif(\tfrac{1}{2}+it)
   \big/\max|\mathrm{Re}\,\xif|$:
  영점 MSE $=1.8\times 10^{-5}$, 비영점 MSE $=5.8\times 10^{-4}$,
  비율 $0.03\times$ (갭 축소 $99.8\%$).
\end{itemize}
정규화 타겟은 영점/비영점 갭을 완전히 소멸시킬 뿐 아니라 \emph{역전}시키며, 영점이 비영점보다 더 쉬워진다. 이는 $\pm\pi$ 불연속이 \emph{유일한} 병목이었음을 확정한다: 타겟이 매끄러우면 네트워크 아키텍처와 학습 파이프라인에 잔존 결함이 없다. 프로그램의 해석적 내용은 보존되며; 오직 $\{0,\pi\}$ 부호화만 교체가 필요하다. (스크립트 \texttt{scripts/smooth\_zero\_crossing.py}, 출력 \texttt{smooth\_zero\_crossing.\{json,txt\}}.)

\smallskip
우리는 고립 및 매끄러운 타겟 실험을 프로그램의 해석적 내용 --- \cref{eq:psi_supervised}의 동일시 $\Psip(t)=\arg\xif(\tfrac{1}{2}+it)$ --- 이 옳은 자기 쌍대 타겟이라는 증거로 기록한다; 매끄러운 타겟 변형이 남은 표현 가능성 장벽을 해소한다.
\end{remark}


\subsection{리만--지겔 $\vartheta$ 함수와의 연결}

임계선 위 $\zeta$의 고전적 위상 분석은 리만--지겔 세타 함수를 사용한다:
\begin{equation}
  \vartheta(t) = \arg\,\Gamma\!\left(\frac{1/4 + it/2}{1}\right) -
  \frac{t}{2}\log\pi
  \;\approx\;
  \frac{t}{2}\log\frac{t}{2\pi} - \frac{t}{2} - \frac{\pi}{8} +
  O(t^{-1}).
  \label{eq:riemann_siegel_theta}
\end{equation}
하디 $Z$-함수는 다음과 같이 정의된다:
\begin{equation}
  Z(t) = e^{i\vartheta(t)}\zeta(\tfrac{1}{2}+it),
  \label{eq:hardy_z}
\end{equation}
이는 실수 $t$에 대해 실수이다. 부호 변화가 임계선 위 $\zeta$의 영점을 표시한다. 그람 점 $g_n$은 $\vartheta(g_n) = n\pi$로 정의되어, $t$-축을 따른 위상 양자화 격자를 제공한다.

\begin{definition}[그람 점]
\label{def:gram}
$n$번째 그람 점 $g_n$은 다음을 만족한다:
\begin{equation}
  \vartheta(g_n) = n\pi, \quad n \in \mathbb{Z}.
  \label{eq:gram_points}
\end{equation}
연속된 그람 점 사이에서 위상 $\vartheta(t)$는 정확히 $\pi$만큼 진행한다. 논증 원리에 의해, 높이 $T$까지의 $\zeta(\tfrac{1}{2}+it)$ 영점 수는:
\begin{equation}
  N(T) = \frac{1}{\pi}\Big[\vartheta(T) +
  \mathrm{Im}\log\zeta(\tfrac{1}{2}+iT)\Big] + O(1).
  \label{eq:zero_counting}
\end{equation}
\end{definition}

\begin{remark}[QROP-Net 대응: 위상 격자]
\label{rem:gram_grid}
그람 점 $\{\vartheta(g_n) = n\pi\}$는 임계선 위의 이산 위상 격자를 형성하며, 이는 QROP-Net의 SLM 위 $2^{d_u}$개 이산 위상 수준과 직접 유비된다. 두 시스템 모두에서 근본적 질문은 동일하다: \textit{연속 신호(광학 파면, $\xif$-함수)가 규정된 이산 위상 격자와 정렬되는가?}
\end{remark}


% ============================================================================
\section{게이트 구조: 소수 탐지로서의 위상 공명}
\label{sec:gates}
% ============================================================================

이제 임계선 위 또는 근방의 ``공명'' 점을 식별하는 2-게이트 기준을 형식화한다.

\subsection{게이트 A: 위상 양자화 (정렬)}

\begin{definition}[게이트 A --- 위상 양자화 조건]
\label{def:gate_a}
복소 $s$-평면에서 기준점 $(1/2, t_{\mathrm{ref}})$으로부터 $(\sigma, t)$까지의 경로를 고려한다. 이 경로를 따른 총 축적 위상 변화는:
\begin{equation}
  \Delta v(\sigma, t) = \int_{t_{\mathrm{ref}}}^{t}
  v_t(\tfrac{1}{2}, \tau)\,d\tau +
  \int_{1/2}^{\sigma}v_\sigma(\sigma', t)\,d\sigma'.
  \label{eq:phase_accumulation}
\end{equation}
게이트~A는 이 위상 축적이 $\pi$의 정수배일 것을 요구한다:
\begin{equation}
\boxed{
  \GateA: \quad \Delta v(\sigma, t) = m\pi, \quad m \in \mathbb{Z}.
}
\label{eq:gate_a}
\end{equation}
\end{definition}

임계선 ($\sigma = 1/2$) 위에서 경로는 $t$-방향으로만 축소되며, 게이트~A는:
\begin{equation}
  \Delta v(\tfrac{1}{2}, t) =
  \vartheta(t) - \vartheta(t_{\mathrm{ref}}) +
  \arg\zeta(\tfrac{1}{2}+it) - \arg\zeta(\tfrac{1}{2}+it_{\mathrm{ref}})
  = m\pi,
  \label{eq:gate_a_critical}
\end{equation}
$t_{\mathrm{ref}} = 0$이고 $\vartheta(0) = 0$이면, 그람 점 조건 $\vartheta(t) = m\pi$로 환원된다 ($\arg\zeta$ 보정까지).

\begin{remark}[QROP-Net 대응: $\sin^2$ 양자화]
\label{rem:gate_a_parallel}
게이트~A는 QROP-Net 양자화 손실 $\mathcal{L}_{\mathrm{quant}}$ (\cref{eq:loss_quant})의 \textit{정확한 해석적 정수론 대응물}이다. QROP-Net에서 $\sin^2$ 벌칙은 연속 위상을 가장 가까운 이산 수준으로 유도하며; 게이트~A는 연속 위상 $v(\sigma,t)$가 $\pi$의 정수배로 양자화된 점을 선택한다. 위상 양자화 연산자 (\cref{eq:pqo})에서 $L=2$로 놓으면 $\mathcal{P}(v; 2) = \sin^2(v) = 0$ ($v = m\pi$에서)를 얻으며, 이것이 정확히 게이트~A이다.
\end{remark}


\subsection{게이트 B: 위상 정류 (곡률)}

\begin{definition}[게이트 B --- 위상 정류 조건]
\label{def:gate_b}
위상이 $\sigma$-방향에서 \textit{정류}이다:
\begin{equation}
\boxed{
  \GateB: \quad v_\sigma(\sigma, t) = 0.
}
\label{eq:gate_b}
\end{equation}
코시--리만 관계식 (\cref{eq:cauchy_riemann})에 의해, 이는 다음과 동치이다:
\begin{equation}
  \mathrm{Im}\!\left(\frac{\xif'(\sigma+it)}{\xif(\sigma+it)}\right) = 0,
  \label{eq:gate_b_cr}
\end{equation}
그리고 진폭 극값 조건 $u_t(\sigma,t) = 0$과도 동치이다, 즉
\begin{equation}
  \frac{\partial}{\partial t}|\xif(\sigma + it)| = 0.
  \label{eq:amplitude_extremum}
\end{equation}
게이트~B는 진폭 $\Ampl(\sigma,t)$가 $t$-방향에서 국소 극대(또는 안장점)를 갖는 점을 선택한다.
\end{definition}

\begin{remark}[물리적 해석: 공명 봉우리]
\label{rem:resonance_peak}
산란 이론~\cite{khare1997}에서 $\xif(s)$는 산란 진폭으로 해석될 수 있으며, 비자명 영점은 공명에 대응한다. 공명에서 산란 진폭은 \textit{국소 극대}에 도달하고 (게이트~B: $u_t = 0$) 위상 이동은 $\pi$-증분을 통과한다 (게이트~A: $\Delta v = m\pi$). 이는 정확히 공진 공동에서 보강 간섭의 조건이다: 파동이 공명에서 ``머물며'', 위상 결맞음을 유지하면서 진폭을 축적한다.
\end{remark}

\begin{remark}[QROP-Net 대응: 충실도 손실]
\label{rem:gate_b_parallel}
게이트~B는 관심 점에서 진폭 $\Ampl(t)$가 정류일 것을 요구한다. QROP-Net에서 충실도 손실 $\mathcal{L}_{\mathrm{fid}}$ (\cref{eq:loss_fid})는 복원된 진폭 패턴을 목표와 비교한다. $\mathcal{L}_{\mathrm{fid}}$를 최소화하면 시스템은 진폭 정류 상태로 향한다: 수렴 시 예측 세기가 목표와 일치하여 $\partial \hat{I}/\partial(\text{params}) = 0$이 된다. 두 시스템 모두 정확한 신호 복원의 필요 조건으로서 진폭 안정성을 강제한다.
\end{remark}


\subsection{결합 판별식: 공명점}

\begin{definition}[공명 소수 선택]
\label{def:resonant_prime}
$s$-평면의 점 $(\sigma^*, t^*)$가 두 게이트를 동시에 만족하면 \textit{공명점}이다:
\begin{equation}
\boxed{
  \pis(\sigma^*, t^*) = 1 \quad\iff\quad
  \begin{cases}
    \GateA: & \Delta v(\sigma^*, t^*) = m\pi \\
    \GateB: & v_\sigma(\sigma^*, t^*) = 0
  \end{cases}
}
\label{eq:discriminant}
\end{equation}
$\sigma^* = \tfrac{1}{2}$ (임계선 위)이면, 공명점은 $\zeta(s)$의 영점 또는 준영점에 대응한다. 정수 $n$이 공명점 $t_n$에서 위상 정렬 조건 $\Phiz(t_n) = \Psip(n)$을 만족하면 \textit{공명 소수}로 선택된다.
\end{definition}

두 조건의 동시 만족이 결정적 특성이다:
\begin{itemize}
  \item 게이트~A만으로는 너무 많은 점을 선택한다 (영점 사이의 모든 그람 점 포함).
  \item 게이트~B만으로는 위상 정렬되지 않은 진폭 극값을 선택한다.
  \item 두 조건의 \textit{교차}가 위상 양자화되고 진폭 극값인 후보---진정한 공명의 특징---를 산출한다.
\end{itemize}


% ============================================================================
\section{$\sigma$-리프트와 오차 허용}
\label{sec:sigma_lift}
% ============================================================================

\subsection{임계선으로부터의 편차}

\begin{definition}[$\sigma$-리프트]
\label{def:sigma_lift}
두 게이트를 만족하는 공명점 $(\sigma^*, t^*)$에 대해 $\sigma$-리프트는:
\begin{equation}
  \ell := \sigma^* - \frac{1}{2}.
  \label{eq:sigma_lift}
\end{equation}
리만 가설은 모든 비자명 공명점에 대해 $\ell = 0$이라는 진술과 동치이다.
\end{definition}

\begin{proposition}[오차 임계치 유비]
\label{prop:threshold}
$\sigma$-리프트 $\ell$은 Kyber의 LWE 오차 $|e|_{\Zq}$와 동일한 구조적 역할을 수행한다:
\begin{center}
\begin{tabular}{lcc}
\toprule
& \textbf{QROP-Net} & \textbf{임계선} \\
\midrule
변수         & $|e|_{\Zq}$            & $\ell = \sigma^* - 1/2$ \\
임계치       & $\ceil{q/4} = 1920$    & $\ell = 0$ \\
성공         & $|e|_{\Zq} < 1920$     & $\ell = 0$ (RH) \\
실패 모드    & $\mathrm{DFR} = 1$     & 임계선 밖의 영점 \\
민감도       & $\Delta d = 0.1\,$mm   & $\ell \neq 0$ \\
\bottomrule
\end{tabular}
\end{center}
두 시스템 모두에서, 신호가 비가역적으로 소실되는 \textbf{경성 경계}가 존재한다. 한 시스템의 ``오차''는 문자 그대로의 복호화 실패이며; 다른 시스템에서는 대칭축으로부터 영점의 이탈이다.
\end{proposition}


\subsection{산란 해석과 유니타리성}

Khare~\cite{khare1997}는 $\zeta(s)$를 산란 진폭으로 해석하고 비자명 영점을 공명으로 보는 것을 제안했다. LeClair와 Mussardo~\cite{leclair2023}는 영점을 \textit{산란 시스템의 양자화된 에너지}로 취급하여 이를 형식화하고, 영점 위치를 파악하는 베테 안자츠/위상 양자화 접근법을 도출했다.

그들의 프레임워크에서, $\sigma^* \neq \tfrac{1}{2}$인 영점 $s = \sigma^* + it^*$는 \textit{비유니타리} 산란 행렬에 대응한다: 시스템에 흡수 또는 누출이 있다. $\sigma$-리프트 $\ell$은 \textit{비유니타리성의 정도}를 정량화한다:
\begin{equation}
  |\ell| = |\sigma^* - \tfrac{1}{2}| \propto
  \text{``흡수 계수''},
  \label{eq:non_unitarity}
\end{equation}
여기서 $\ell = 0$은 유니타리(무손실) S-행렬에, $\ell \neq 0$은 산일 채널에 대응한다.

\begin{remark}[QROP-Net 대응: 채널 손실]
\label{rem:unitarity}
QROP-Net에서 광학 채널은 위상 신호를 열화시키는 잡음 (포아송, 읽기 잡음, ADC 양자화)을 도입한다. ``채널 품질''은 LWE 오차로 직접 측정된다: 더 높은 잡음 $\to$ 더 큰 $|e|_{\Zq}$ $\to$ DFR 임계치에 더 가까움. 유사하게 $|\ell| > 0$은 해석적 연속에서의 ``잡음''을 나타낸다---영점이 정확히 대칭축 위에 있지 않다. 두 시스템 모두 이상적 기준(격자점, 임계선)으로부터의 거리로 신호 열화를 측정한다.
\end{remark}


\subsection{베리 위상과 홀로노미}
\label{sec:berry}

$\xif(s)$의 영점 주위 위상 감김은 위상적 불변량을 제공한다.

\begin{proposition}[위상 감김수]
\label{prop:winding}
$\mathcal{C}$를 $\xif(s)$의 정확히 $N$개 영점을 둘러싸는 $s$-평면의 닫힌 윤곽이라 하자. 그러면:
\begin{equation}
  \oint_{\mathcal{C}} \nabla v \cdot d\mathbf{s} = 2\pi N,
  \label{eq:winding}
\end{equation}
이는 논증 원리에 의한 것이다. 이는 \textit{홀로노미}에 대한 진술이다: 루프 주위의 총 위상 변화는 둘러싸인 영점당 $2\pi$이며, 양자역학의 베리 위상~\cite{berry1984}과 직접 유비된다---순환 경로에 걸쳐 축적된 기하학적 위상이 천 수의 $2\pi$배인 것.
\end{proposition}

게이트~A의 조건 $\Delta v = m\pi$ (기준점에서 $(\sigma, t)$까지의 열린 경로를 따라)는 \textit{부분 홀로노미} 조건이다: 경로가 $m$번의 반회전을 축적하여, 각 공명점에 \textit{위상적 위상 지수} $m$을 고정한다.

\begin{remark}[QROP-Net 대응: 위상 래핑]
\label{rem:phase_wrapping}
QROP-Net에서 2-채널 $(\mathrm{Re}, \mathrm{Im})$ 출력은 $0/2\pi$ 위상 래핑 불연속을 피하기 위해 특별히 설계되었다. 동일한 위상적 장애물이 $\xif$-함수에서도 나타난다: $v(\sigma,t)$는 영점에서 분기 절단과 $2\pi$ 점프를 갖는다. 감김수 (\cref{eq:winding})는 이 점프를 계수하며, 이는 QROP-Net의 원형 통계 (\cref{eq:circular_mean})가 $0/2\pi$ 경계를 올바르게 처리하는 것과 같다.
\end{remark}


% ============================================================================
\section{다원군 영점 탐지: 다원적 프레임}
\label{sec:polyadic}
% ============================================================================

게이트 조건의 수치적 구현에는 임계선 위 로그 도함수 $\xif'/\xif$의 견고한 평가가 필요하며, 여기서 $\xif$는 영점을 가져 $\xif'/\xif$가 발산한다. \textit{다원적 프레임} 접근법은 여러 방향 채널을 사용하여 특이점을 평균화함으로써 이를 해결한다.

\subsection{표적 함수와 평가점}

표적 함수 $f$는 리만 $\xif$와 디리클레 $L$-함수에 걸쳐 통합된다:
\begin{align}
  \xif(s) &= \frac{1}{2}s(s-1)\pi^{-s/2}\Gamma(s/2)\zeta(s),
  \label{eq:xi_target}\\
  \Lambda(s, \chi) &= \left(\frac{q}{\pi}\right)^{(s+\epsilon)/2}
  \Gamma\!\left(\frac{s+\epsilon}{2}\right)L(s,\chi),
  \label{eq:lambda_target}
\end{align}
여기서 $q$는 모듈러스, $\epsilon \in \{0,1\}$은 디리클레 지표 $\chi$의 패리티 지시자이며, 모든 평가는 임계선 $s(t) = \tfrac{1}{2} + it$ 위에서 수행된다.

로그 도함수는 오일러 곱을 통해 소수와 직접 연결된다:
\begin{equation}
  \frac{f'(s)}{f(s)} = -\sum_{\rho}\frac{1}{s-\rho} +
  (\text{매끄러운 부분}),
  \label{eq:log_deriv_poles}
\end{equation}
여기서 합은 비자명 영점 $\rho$에 대해 수행되며, 지배적 기여는 근방 영점에서 온다. 이는 오일러 곱 표현과 직접 관련된다:
\begin{equation}
  \zeta(s) = \prod_{p\;\text{소수}}(1-p^{-s})^{-1}, \qquad
  L(s,\chi) = \prod_p(1-\chi(p)p^{-s})^{-1},
  \label{eq:euler_products}
\end{equation}
여기서:
\begin{equation}
  -\frac{f'(s)}{f(s)} \approx \sum_{n \geq 1}
  \Lambda(n)\chi(n)\,n^{-s},
  \label{eq:euler_log_deriv}
\end{equation}
$\Lambda(n)$은 폰 만골트 함수이다 ($n = p^k$이면 $\Lambda(n) = \log p$, 아니면 $0$). 소수에 대한 진동 합이 게이트~A가 양자화하는 진동을 직접 생성한다.


\subsection{방향 채널}
\label{sec:channels}

$f'/f$를 단일 방향으로만 평가하는 대신 (잡음과 앨리어싱에 민감), $m$개의 방향 채널을 도입한다. 각 채널 $k$는 기저 벡터를 갖는다:
\begin{equation}
  \mathbf{e}_k = (\cos\theta_k, \;\lambda_k\sin\theta_k),
  \label{eq:channel_vector}
\end{equation}
여기서 $\theta_k$는 채널 각도, $\lambda_k$는 비등방 축척 인자이다. 권장 3-채널 표준 집합은:
\begin{equation}
  (\theta, \lambda, w) = (-30^\circ, 1.10, 0.40), \;
  (15^\circ, 0.85, 0.35), \;
  (60^\circ, 1.00, 0.25),
  \label{eq:channel_set}
\end{equation}
채널 가중치 $w_k$는 $\sum w_k = 1$로 정규화된다.

비등방 축척 행렬이 국소 기하에 적응한다:
\begin{equation}
  M(t) = \mathrm{diag}\!\left(\frac{1}{\mu(t)},\;
  \frac{1}{\nu(t)}\right),
  \label{eq:aniso_scale}
\end{equation}
여기서 $\mu(t), \nu(t) > 0$은 매끄러운 프로필 함수이다. 단계 $h > 0$에 대한 실제 방향 변위는:
\begin{equation}
  \delta_k(t) = h\,(v_{k,x} + iv_{k,y}), \quad
  \mathbf{v}_k = R(\theta(t))\,M(t)\,\mathbf{e}_k,
  \label{eq:displacement}
\end{equation}
여기서 $R(\theta)$는 회전 행렬이고 $\theta(t)$는 $t_0$ 주위에서 매끄럽게 변한다.


\subsection{중심 차분 로그 도함수}
\label{sec:central_diff}

채널 $k$를 따른 유효 로그 도함수는 중심 차분으로 계산된다:
\begin{equation}
  D_k\xif(s) \approx
  \frac{\xif(s + \delta_k) - \xif(s - \delta_k)}{2h\,\xif(s)}.
  \label{eq:central_diff}
\end{equation}
\textit{가중 다채널 로그 도함수}는:
\begin{equation}
  \Lhat(t) = \sum_k w_k \cdot
  \frac{f(s + \delta_k) - f(s - \delta_k)}{2h\,f(s)},
  \label{eq:lhat}
\end{equation}
채널 평균화가 방향 편향을 줄이고 영점 근방의 안정성을 향상시킨다. $f$가 해석적이므로 중심 차분의 오차는 $O(h^2)$이다. $\Lhat(t) = v_t(t) + i\,v_\sigma(t)$로 실부와 허부로 분해한다.


\subsection{게이지 ODE에 의한 위상 보정}
\label{sec:gauge_ode}

실부 $v_t(t)$에는 최종 판별식 적용 전에 제거해야 하는 느린 표류가 포함된다. 게이지 ODE를 정의한다:
\begin{equation}
  \tau\dot{\psi} + \psi = v_t, \quad \dot{\phi} = \psi,
  \label{eq:gauge_ode}
\end{equation}
초기 조건 $\phi(t_0) = 0$, $\psi(t_0) = v_t(t_0)$. 여기서 $\tau > 0$은 게이지 감쇠 상수이다 (기본값: $\tau = 0.04$). ODE는 $v_t$를 적분하면서 고주파 잡음을 평활화한다. 수치 기법은 안정적 사다리꼴 이산화를 사용한다:
\begin{align}
  \alpha_j &= \frac{\Delta t_j}{\tau + \Delta t_j},
  \label{eq:ode_alpha}\\
  \psi_j &= \psi_{j-1} + \alpha_j(v_{t,j} - \psi_{j-1}),
  \label{eq:ode_psi}\\
  \phi_j &= \phi_{j-1} + \tfrac{1}{2}(\psi_j + \psi_{j-1})\,\Delta t_j.
  \label{eq:ode_phi}
\end{align}


\subsection{최종 판별 함수}
\label{sec:discriminant}

게이지 보정된 최종 판별식은 두 게이트를 결합한다:
\begin{equation}
\boxed{
  F_2(t) = \mathrm{Im}\!\Big\{
    e^{-i\phi(t)}\Big(\Lhat(t) - \dot{\phi}(t)\Big)
  \Big\} = 0.
}
\label{eq:F2}
\end{equation}
$F_2$의 영점 $t^*$에서:
\begin{itemize}
  \item 허부가 사라져, 위상 보정된 로그 도함수가 축적된 게이지 위상과 정렬됨 $\to$ \textbf{게이트~A} (위상 양자화)가 만족됨.
  \item $F_2$의 영점 통과 부호 변화가 위상 교차를 나타냄 $\to$ \textbf{게이트~B} (정류)가 국소적으로 만족됨.
\end{itemize}

\begin{remark}[QROP-Net 대응: 블라인드 자가 교정]
\label{rem:tto_parallel}
게이지 ODE (\cref{eq:gauge_ode})는 \textit{진정한 영점 위치를 모른 채} 위상 표류를 자기 보정한다. 이는 \textit{진정한 암호문을 모른 채} 렌즈 오정렬을 보정하는 QROP-Net의 블라인드 TTO (\cref{eq:blind_loss})와 구조적으로 동일하다. 두 시스템 모두:
\begin{enumerate}
  \item 정답이 아닌 물리/수학 기반 제약을 사용하며,
  \item 소수의 단계로 수렴하고 (TTO는 15회, $F_2$는 $O(N)$ 격자점),
  \item 신호 구조의 내적 일관성을 통해 자가 교정을 달성한다.
\end{enumerate}
\end{remark}


\subsection{수치 검증과 수렴}
\label{sec:polyadic_results}

판별식 $F_2(t) = 0$은 균일 격자 $[T-W, T+W]$에서 $N$개 점으로 수치 탐색된다 ($W = 0.002$--$0.005$). 알고리즘:

\begin{algorithm}[ht]
\DontPrintSemicolon
\caption{다원적 영점 탐지}
\label{alg:polyadic}
\KwIn{표적 $f$, 대역 중심 $T$, 창 $W$, 채널 $\{(\theta_k, \lambda_k, w_k)\}$, 단계 $h$, 게이지 $\tau$}
\KwOut{영점 후보 $\{t^*\}$}
$s_j \gets \tfrac{1}{2} + it_j$ (균일 격자 $t_j \in [T-W, T+W]$)\;
\ForEach{채널 $k$ 및 격자점 $j$}{
  $\delta_k(t_j) \gets h\,(v_{k,x}+iv_{k,y})$ (\cref{eq:displacement} 경유)\;
  $D_k \gets [f(s_j+\delta_k) - f(s_j-\delta_k)] / (2h\,f(s_j))$\;
}
$\Lhat(t_j) \gets \sum_k w_k\,D_k$\;
추출: $v_t \gets \mathrm{Re}(\Lhat)$,
$v_\sigma \gets \mathrm{Im}(\Lhat)$\;
게이지 ODE (\cref{eq:ode_alpha}--\ref{eq:ode_phi}) 적분으로 $\phi(t_j)$ 획득\;
$F_2(t_j) \gets \mathrm{Im}\{e^{-i\phi_j}(\Lhat_j - \dot{\phi}_j)\}$\;
$F_2$의 부호 변화 탐지; 선형 보간으로 정밀화\;
불변성 검증: 프레임 $M_0$ vs $M_2$, 채널 교환, 창 이동\;
\end{algorithm}

\begin{property}[수렴 속도]
\label{prop:convergence}
$h_{\min}$을 시험된 최소 단계라 하자. 수렴 지표는:
\begin{align}
  p_{\mathrm{LD}} &= \frac{\log\sup_t|\mathrm{LD}(h) - \mathrm{LD}(h_{\min})|}
  {\log h} \approx 2.0, \label{eq:conv_ld}\\
  p_{F_2} &= \frac{\log\sup_t|F_2(h) - F_2(h_{\min})|}
  {\log h} \approx 2.0, \label{eq:conv_f2}
\end{align}
이는 중심 차분에 대해 기대되는 $O(h^2)$ 수렴을 확인한다. 영점 위치 $t^*$는 다음에 대해 불변이다:
\begin{itemize}
  \item 프레임 프로필 ($M_0$ 평탄 vs $M_2$ 강): $|t^*_{M_0} - t^*_{M_2}| / \Delta t_{\mathrm{grid}} < 10^{-3}$.
  \item 채널 구성 (3-채널 vs 2-채널): $r_{\mathrm{channel}} < 0.1$.
  \item 창 이동 ($\pm\delta$ 섭동): $r_{\mathrm{window}} < 0.2$.
\end{itemize}
\end{property}

검증된 표적은 다음을 포함한다:
\begin{center}
\begin{tabular}{lcc}
\toprule
\textbf{함수} & \textbf{대역} & \textbf{$t^*$} \\
\midrule
$\xif(s)$ & $t \approx 141.12$ & 141.123703 \\
$L(s,\chi_3^{\mathrm{odd}})$ & $t \approx 140.61$ & 140.611723 \\
$L(s,\chi_8^{\mathrm{even}})$ & $t \approx 141.08$ & 141.083569 \\
$L(s,\chi_5^{\mathrm{cmplx}})$ & $t \approx 140.97$ & 140.966280 \\
\bottomrule
\end{tabular}
\end{center}
모든 결과는 단계 스위프 $h \in \{2\times10^{-7}, 10^{-6}, 5\times10^{-6}, 10^{-5}\}$ 하에서 불변이다.


\subsection{교차 검증: 세 독립 축}
\label{sec:cross_validation}

$F_2 = 0$으로부터의 영점 후보는 근본적으로 다른 오차 프로필을 가진 세 독립 데이터 소스에 대해 검증된다:

\subsubsection{축 A: 부분 오일러 곱 매칭}
부분합 $S_N(t;\chi) = \sum_{n \leq N}\Lambda(n)\chi(n)\, n^{-1/2-it}$는 충분히 큰 $N$에 대해 $s = \tfrac{1}{2}+it$ 근방에서 $-f'(s)/f(s)$를 근사한다. $\Lhat(t^*)$와 $S_N(t^*;\chi)$ 사이의 허용 범위 내 일치가, 영점 후보가 오일러 곱과 일치함을 확인한다.

\subsubsection{축 B: 명시적 체비셰프 $\psi(x)$ 복원}
소수 계수 함수 $\psi(x) = \sum_{n \leq x}\Lambda(n)$은 명시적 공식을 통해 영점으로부터 재구성될 수 있다:
\begin{equation}
  \psi(x) = x - \sum_{\rho}\frac{x^\rho}{\rho} - \log(2\pi) -
  \frac{1}{2}\log(1-x^{-2}),
  \label{eq:explicit_formula}
\end{equation}
여기서 $\rho$는 $\zeta$의 비자명 영점에 대해 수행된다. $F_2 = 0$에서 얻은 영점 높이 $t^*$를 대입하고 재구성된 $\psi(x)$를 정확한 계단 함수와 비교하면 독립적 검증이 된다.

\subsubsection{축 C: 등차수열 편향 신호}
디리클레 $L$-함수에 대해, 등차수열 $p \equiv a \pmod{q}$에서의 소수 분포가 제3의 신호를 제공한다. $\psi(x;q,a)$에 대한 명시적 공식의 편향 항 $\sum_\rho x^\rho/\rho$는 영점 높이에 의해 결정되는 주파수로 진동하여 독립적 확인을 제공한다.


% ============================================================================
\section{\textcolor{ForestGreen}{\rule{0.9ex}{0.9ex}}\ Hardy-$Z$ 부호변화 탐지기로서의 $F_2$-판별식}
\label{sec:f2_hardy}
% ============================================================================

\noindent\textbf{계층 1.} 본 섹션은 본 프로그램의 핵심 검증 경험 결과를 동일시한다: 제~\ref{sec:polyadic}절의 게이지 ODE에서 도입된 다채널 판별식 $F_2$는, 해석 신호 구성까지의 차이를 제외하면, 정확히 Hardy $Z$-함수의 부호 변화 탐지기이다. 동일시는 두 단계로 확립된다: 실함수의 해석 신호가 영점에서 갖는 인수에 대한 해석적 진술 (\cref{prop:phase_zero}), 그리고 $Z(t)$를 근사하는 매끄러운 네트워크가 단순 영점에서 $\arg(Z_{\mathrm{out}})=\pi/2$로 포화되는 이유를 설명하는 신경 계 corollary (\cref{cor:neural_reg}).

\subsection{Hardy $Z$와 해석 신호 동일시}
\label{sec:hardy_analytic}

\cref{eq:hardy_z}에서 Hardy $Z$-함수
\begin{equation}
  Z(t) = e^{i\vartheta(t)}\,\zeta\bigl(\tfrac{1}{2}+it\bigr)
\end{equation}
는 실수 $t$에 대해 실수이며, 그 부호 변화는 임계선 위 $\zeta$의 비자명 영점과 일대일 대응한다. 이로써 $Z$는 리만 영점 위치를 찾도록 설계된 판별식을 시험할 자연스러운 실수 신호이다.

\begin{proposition}[실함수 영점에서의 위상]
\label{prop:phase_zero}
$f : \mathbb{R} \to \mathbb{R}$를 $t_0$에서 단순 영점을 갖는 실함수라 하고, $z(t) = f(t) + i\,H[f](t)$를 그 해석 신호라 하자 (여기서 $H$는 힐베르트 변환). 그러면
\[
  \arg\bigl(z(t_0)\bigr) = \pm\tfrac{\pi}{2},
\]
이고 부호는 $f$가 음에서 양으로 교차하는지 그 반대인지에 의해 결정된다.
\end{proposition}
\begin{proof}
$t_0$에서 $f(t_0)=0$이므로 $\mathrm{Re}\,z(t_0)=0$이며, Titchmarsh 정리에 의해 $H[f](t_0)\neq 0$이 일반적으로 성립한다. 따라서 $z(t_0)$는 순허수이고 그 부호는 교차 방향으로 결정된다.
\end{proof}

\begin{corollary}[신경망의 부호 변화 정규화]
\label{cor:neural_reg}
$Z(t)$를 연속 복소 함수 $Z_{\mathrm{out}}(t)$로 근사하는 신경망은 단순 영점 $t_0$에서 $\arg(Z_{\mathrm{out}}(t_0)) \to \pm\pi/2$를 만족해야 한다. 왜냐하면 \cref{eq:phase_crit}의 이상적 불연속 $\arg\xif : 0 \to \pi$가 그 중간값으로 정규화되기 때문이다.
\end{corollary}

\cref{prop:phase_zero,cor:neural_reg}을 제~\ref{sec:polyadic}절의 판별식 정의 $F_2(t) = \operatorname{Im}\{e^{-i\varphi}(\hat L - \dot\varphi)\}$와 결합하면 다음 동일시를 얻는다.

\begin{theorem}[$F_2 = $ Hardy $Z$ 부호변화 탐지기]
\label{thm:f2_hardy}
$Z_{\mathrm{out}}(t)$를 \cref{eq:gauge_ode}의 게이지 ODE $\dot\varphi = \hat L(t)$를 만족하는 $Z(t)$의 매끄러운 신경 근사라 하자. 그러면 근사가 충실한 임의의 구간 $I\subset\mathbb{R}$에서
\[
  \{\,t\in I : F_2(t)=0\,\}
  \;=\;
  \{\,t\in I : Z(t)=0\,\}
  \;=\;
  \{\,t\in I : \zeta(\tfrac{1}{2}+it)=0\,\}.
\]
\end{theorem}
\begin{proof}
$F_2(t)=\operatorname{Im}\{e^{-i\varphi(t)}(\hat L(t)-\dot\varphi(t))\}=0$은 $\hat L(t)-\dot\varphi(t)$가 $e^{i\varphi(t)}$와 공선일 때와 동치이며, 이는 게이지 ODE에 의해 해석 신호의 인수가 중간값 $\pm\pi/2$에 도달하는 것과 동치이다. 따라서 \cref{prop:phase_zero}에 의해 $Z$의 영점과 동치이다.
\end{proof}

\subsection{경험 검증: $t\in[10,50]$에서 10/10 영점 탐지}
\label{sec:f2_low}

\begin{observation}[$F_2$ 탐지, 저구간]
\label{obs:f2_low}
Adam ($\eta=10^{-3}$, 200 epoch)으로 $t\in[10,50]$에서 $Z(t)$의 매끄러운 신경 근사는 10개 LMFDB 영점 모두에서 $|F_2(t_k)| < 0.021$로 수렴하며, 15 epoch 이후 잔차 손실 $\mathcal{L}_{\mathrm{res}} < 10^{-4}$이다. 출력 위상 $\arg(Z_{\mathrm{out}}) \in [1.47,1.62]$는 \cref{cor:neural_reg}과 정확히 일치하게 $\pi/2\approx 1.5708$ 주위 $\pm 0.06$ 창 안에 들어간다.
\end{observation}

\begin{table}[h]
\centering
\caption{Adam~$+$~$\cos^2$ PQO의 $t\in[10,50]$ 탐지 결과, 10 LMFDB 영점, 200 epoch.}
\label{tab:f2_low}
\begin{tabular}{c|cccc}
\toprule
$t_k$ & $|F_2|$ & $\sin^2$ PQO & $\cos^2$ PQO & $\arg(Z_{\mathrm{out}})$ \\
\midrule
14.13 & 0.021 & 1.000 & \textbf{0.000} & 1.581 \\
21.02 & 0.001 & 1.000 & \textbf{0.000} & 1.563 \\
25.01 & 0.003 & 1.000 & \textbf{0.000} & 1.557 \\
30.42 & 0.019 & 1.000 & \textbf{0.000} & 1.562 \\
32.94 & 0.000 & 1.000 & \textbf{0.000} & 1.562 \\
37.59 & 0.009 & 1.000 & \textbf{0.000} & 1.571 \\
40.92 & 0.000 & 1.000 & \textbf{0.000} & 1.565 \\
43.33 & 0.004 & 1.000 & \textbf{0.000} & 1.573 \\
48.01 & 0.003 & 1.000 & \textbf{0.000} & 1.584 \\
49.77 & 0.005 & 1.000 & \textbf{0.000} & 1.582 \\
\bottomrule
\end{tabular}
\end{table}

두 관찰이 강조되어야 한다. 첫째, $|F_2|<0.021$과 $\cos^2\mathrm{PQO}<0.0002$의 동시 만족은 자명하지 않은 결합 조건이다: 네트워크가 (i) 해석 신호의 허수 성분이 소멸하고 (ii) 시클로토믹 격자의 양자화된 위상 어트랙터가 만족되는 점을 찾았다는 뜻이다. 같은 10개 점에서 두 조건이 동시에 만족될 수 있다는 사실이 바로 $\cos^2$ Maslov 보정이 작동하고 있음을 보여준다. 둘째, 잔존 격차 $\arg(Z_{\mathrm{out}})\neq\pi/2$ (최대 $\approx 0.01$)는 유한 학습의 아티팩트이며 구조적 장애물이 아니다; 학습 시간이 길어질수록 격차는 단조적으로 좁아진다.

\subsection{영점 난이도 위계와 Lehmer 현상}
\label{sec:lehmer}

모든 영점이 동등하게 쉽게 탐지되는 것은 아니다. \cref{tab:f2_low}의 10개 영점은 $|F_2|$에서 거의 두 자릿수에 이르는 범위를 가진다: 가장 쉬운 영점은 $|F_2|\lesssim 10^{-3}$으로 수렴하고, 가장 어려운 영점 ($t_1=14.13$과 $t_6=37.59$)은 $2\times 10^{-2}$ 근방에서 포화된다.

\begin{remark}[Lehmer--Odlyzko 연결]
\label{rem:lehmer}
$Z(t)$가 실수축을 거의 스치지만 교차하지 않는 영점 쌍 --- 고전적 Lehmer 현상 --- 은 Montgomery--Odlyzko의 GUE 간격 통계에서 가까운 영점 쌍에 대응한다. 우리 설정에서 이들은 해석 신호의 위상이 $\pm\pi/2$ 근처에서 더 오래 머무는 영점이다. 수치 실험(78개 영점, $t\in[0,200]$)에서 근접 쌍의 다발 곡률 프로파일 $\kappa$가 정성적으로 더 넓고(FWHM 비율 ${\approx}\,3.4$, $p{=}0.01$) 다중 피크 구조를 보인다.

고해상도 후속 실험에서 $\delta$-강건성 검증을 확인하였다: 정규화 간격 $g_{\mathrm{norm}}$과 피크 곡률 $\kappa(t_0+\delta)$의 Spearman 순위 상관이 $\rho = 0.835$ ($p < 10^{-21}$, $n=78$)이며, 세 오프셋 값 $\delta\in\{0.01, 0.03, 0.05\}$ 모두에서 $|\rho|>0.83$을 유지한다. 다만 FWHM--간격 상관은 약하며($\rho = -0.276$, $p=0.014$), 간격이 피크 \emph{높이}는 조절하지만 피크 \emph{폭}에는 영향이 적음을 시사한다.
\textcolor{ForestGreen}{[확인 2026-04-14, $\delta$-강건 3/3]}
\end{remark}

% ============================================================================
\section{\textcolor{ForestGreen}{\rule{0.9ex}{0.9ex}}\ 고위영점으로의 확장: $t\in[100,200]$}
\label{sec:high_zeros}
% ============================================================================

\noindent\textbf{계층 1.} \cref{thm:f2_hardy}의 동일시 $F_2 = \text{Hardy-}Z\text{ 부호변화 탐지기}$는 $Z(t)$의 해석 신호에 대한 진술이며 본질적 높이 의존성을 갖지 않는다. 자연스러운 건전성 검사는 제~\ref{sec:f2_low}절의 수치 파이프라인이 타겟 구간이 한 자릿수 올라가도 완전 탐지를 계속 생산하는지 보는 것이다. 결과적으로 그렇다.

\begin{observation}[고위영점 탐지]
\label{obs:high_zeros}
제~\ref{sec:f2_low}절의 실험을 $t\in[100,200]$에서 50개 LMFDB 영점 ($k=26,\ldots,79$에 대해 \texttt{mpmath.zetazero}$(k)$로 계산), 500 격자점, 그리고 확대된 네트워크 (은닉폭 64, $33{,}540$ 매개변수)로 반복하였다. Adam~$+$~$\cos^2$ PQO, $\eta=10^{-3}$, 200 epoch으로:
\end{observation}

\begin{table}[h]
\centering
\caption{저구간 대 고구간 탐지.}
\label{tab:high_zeros}
\begin{tabular}{l|c|c}
\toprule
\textbf{지표} & $t\in[10,50]$, 10 영점 & $t\in[100,200]$, 50 영점 \\
\midrule
탐지율                          & 10/10          & \textbf{50/50} \\
$\max|F_2|$                 & 0.021          & 0.060          \\
$\max \cos^2$ PQO           & 0.0002         & 0.0011         \\
$\arg(Z_{\mathrm{out}})$ 범위 & $[1.55,1.58]$ & $[1.54,1.60]$ \\
학습 벽시계 시간            & $2$\,분       & $19$\,분      \\
\bottomrule
\end{tabular}
\end{table}

50개 고위영점 모두 $|F_2|<0.1$을 만족하고, 48개가 $0.04$ 이하이다. 가장 어려운 두 영점 --- $t=118.79$ ($|F_2|=0.060$)와 $t=114.32$ ($|F_2|=0.058$) --- 조차도 탐지 임계치 내에 편하게 들어있다. \cref{rem:kuramoto_transition}의 임계 동기화 전이가 $\sim 10$ epoch에서 일어나고, $\arg(Z_{\mathrm{out}})\approx\pi/2$ 패턴도 저구간과 동일하며, Maslov 보정과 위상 잠금 메커니즘이 특정 높이 창의 아티팩트가 아님을 확인시켜준다. 이것이 비자명한 스케일에서 \cref{thm:f2_hardy}의 직접적 수치 확증이다.

\begin{remark}[확장이 확립하지 \emph{않는} 것]
\label{rem:high_zeros_limits}
실험은 $F_2$-탐지기가 그 정답 --- 영점 목록 --- 이 공급될 때 계속 작동함을 검증한다. 탐지기가 미지의 영점을 \emph{찾아낼 수} 있는지는 검증하지 \emph{않는다}: 학습 루프가 $F_2$를 평가할 격자점을 결정하기 위해 LMFDB 영점 목록을 사용하기 때문이다. 자가 교정, 블라인드 영점 발견의 문제는 열려 있으며, 제~\ref{part:failed}부와 \cref{sec:open}의 열린 문제 논의의 동기 중 하나이다.
\end{remark}

% ############################################################################
% 제2-b부 — 이론적 기초: 엄밀한 결과
% ############################################################################

% [Paper A: ξ-다발 기하학 — 증명]
\part{$\xi$-다발의 이론적 기초}
\label{part:theory}

\noindent
앞선 부들은 $\xi$-다발 프레임워크를 수치 실험을 통해 발전시켰다.
이제 \emph{엄밀한} 이론적 기반을 확립한다. 아래 결과들은 고전
복소해석학 --- 주로 편각 원리와 아다마르 곱 --- 으로부터 따르지만,
이를 U(1)-다발의 성질로 \emph{기하학적으로 재서술}한 것은, 우리가
아는 한, 새로운 것이다.

본 부 전체에서 $\Lambda$는 차수~$n$의 보형 형식에 대응하는
완비 $L$-함수를 나타낸다: GL(1)의 경우 $\xif(s)$ 또는 $\Lambda(s,\chi)$,
GL(2)의 경우 $(2\pi)^{-s}\Gamma(s)L(s,f)$ (가중치~2 모듈러 형식)
또는 적절한 $\Gamma$-보정 완비. 핵심 요구사항은 $\Lambda$가 유한 위수의
유리형 함수이고 고립된 영점 $\{\rho_k\}$를 가지며 아다마르형
곱 표현을 허용한다는 것이다.


% ============================================================================
\section{영점의 홀로노미 특성화}
\label{sec:holonomy_thm}
% ============================================================================

\begin{definition}[다발 접속과 홀로노미]
\label{def:connection_holonomy}
$\Lambda(s)$를 영역 $\Omega \subset \mathbb{C}$에서 정칙이고
고립 영점 $\{\rho_k\}$를 갖는 완비 $L$-함수라 하자.
\textbf{접속 1-형식}(대수적 도함수)은
\begin{equation}
  \omega := d\log\Lambda = \frac{\Lambda'(s)}{\Lambda(s)}\,ds,
  \label{eq:connection_form}
\end{equation}
이고, 닫힌 곡선 $\mathcal{C} \subset \Omega\setminus\{\rho_k\}$에
대한 \textbf{홀로노미}는
\begin{equation}
  \mathrm{Hol}(\mathcal{C}) := \frac{1}{2\pi i}
  \oint_{\mathcal{C}} \omega
  = \frac{1}{2\pi i}\oint_{\mathcal{C}}
  \frac{\Lambda'(s)}{\Lambda(s)}\,ds
  \label{eq:holonomy_def}
\end{equation}
이다. \textbf{모노드로미 위상}은
$\mu(\mathcal{C}) := 2\pi\,\mathrm{Hol}(\mathcal{C})$로,
$\Lambda$의 $\mathcal{C}$를 따른 총 위상 감김을 측정한다.
\end{definition}

\begin{theorem}[홀로노미 특성화]
\label{thm:holonomy}
$\Lambda$를 위와 같이 두고, $\mathcal{C}_r(s_0)$를 $s_0$ 중심
반지름 $r$의 양방향 원이라 하자. $r$을 $s_0$ 자체(가능한 경우)
외에 다른 영점을 포함하지 않을 만큼 충분히 작게 잡는다.
\begin{enumerate}
  \item[\textup{(i)}] $s_0 = \rho$가 $\Lambda$의 $m$차 영점이면,
    \begin{equation}
      \mathrm{Hol}(\mathcal{C}_r(\rho)) = m,
      \qquad
      \mu(\mathcal{C}_r(\rho)) = 2\pi m.
      \label{eq:hol_zero}
    \end{equation}
  \item[\textup{(ii)}] $\Lambda(s_0) \neq 0$이고 $\mathcal{C}_r(s_0)$가
    영점을 포함하지 않으면,
    \begin{equation}
      \mathrm{Hol}(\mathcal{C}_r(s_0)) = 0,
      \qquad
      \mu(\mathcal{C}_r(s_0)) = 0.
      \label{eq:hol_nonzero}
    \end{equation}
  \item[\textup{(iii)}] 일반적으로, 중복도 $m_1, \ldots, m_k$의
    영점 $\rho_1, \ldots, \rho_k$를 둘러싸는 임의의 곡선 $\mathcal{C}$에
    대해:
    \begin{equation}
      \mathrm{Hol}(\mathcal{C}) = \sum_{j=1}^{k} m_j.
      \label{eq:hol_general}
    \end{equation}
\end{enumerate}
\end{theorem}

\begin{proof}
이것은 \textbf{편각 원리}(argument principle)의 직접적 적용이다
(Ahlfors \cite{ahlfors1979}, 5장 참조). $\Lambda$는 $\Omega$에서
영점 $\{\rho_k\}$를 갖고 극점이 없으므로(완비 $L$-함수는 전해석적이거나
완비 인자에 의해 알려진 극점만 제거), 대수적 도함수
$\Lambda'/\Lambda$는 $\rho_k$에서 잔류값 $m_k$(영점의 위수)인
단순극을 갖는 유리형 함수이다. 따라서:
\[
  \frac{1}{2\pi i}\oint_{\mathcal{C}}
  \frac{\Lambda'(s)}{\Lambda(s)}\,ds
  = \sum_{\rho_k \,\in\, \mathrm{int}(\mathcal{C})} m_k,
\]
이로부터 (i)--(iii)이 즉시 따른다. (ii)의 경우
$\mathcal{C}_r(s_0)$ 내부에 영점이 없으므로 합이 공집합이고
적분이 소멸한다.
\end{proof}

\begin{remark}[기하학적 내용]
\label{rem:hol_geometric}
정리~\ref{thm:holonomy}는 \textbf{고전적 편각 원리}를 U(1)-다발
$P = \{(s, e^{i\arg\Lambda(s)}) : s \in \Omega\setminus\{\rho_k\}\}$
의 접속 $\omega$에 관한 명제로 재서술한 것이다. 이 언어에서: 영점은
\emph{비자명한 홀로노미}(위상적 전하)의 점이고, 비영점은
\emph{평탄}(자명한 홀로노미)이다. 이것이 관측~\ref{obs:midpoint_nonlocal}과
\S\ref{sec:summary}의 모노드로미 측정의 이론적 기초이다.
\end{remark}

\begin{remark}[$\zeta$의 모든 영점은 단순 --- 조건부]
\label{rem:simple_zeros}
$\zeta(s)$의 모든 비자명 영점이 단순($m = 1$)이라는 널리 믿어지는
추측 하에, 정리~\ref{thm:holonomy}(i)은 모든 영점에서 $\mu = 2\pi$를
준다. 수치 측정(mono$/\pi = 2.000 \pm 0.000$, 10/10 영점,
\S\ref{sec:summary})은 이와 일치하며, 단순성에 대한
독립적 수치 증거로 기능한다.
\end{remark}


% ============================================================================
\section{영점에서의 곡률 집중}
\label{sec:curvature_thm}
% ============================================================================

\begin{definition}[다발 곡률 스칼라]
\label{def:curvature_scalar}
$s \in \Omega\setminus\{\rho_k\}$에서 $\xif$-다발의
\textbf{곡률 스칼라}는
\begin{equation}
  \kappa(s) := \left|\frac{\Lambda'(s)}{\Lambda(s)}\right|^2 = |\omega(s)|^2
  \label{eq:kappa_def}
\end{equation}
이다.
\end{definition}

\begin{theorem}[곡률 집중]
\label{thm:curvature}
$\rho$를 $\Lambda$의 $m$차 영점이라 하자.
\begin{enumerate}
  \item[\textup{(i)}] \textbf{영점에서 발산.} $s \to \rho$일 때,
    \begin{equation}
      \kappa(s) = \frac{m^2}{|s - \rho|^2} + O\!\left(\frac{1}{|s-\rho|}\right).
      \label{eq:kappa_diverge}
    \end{equation}
    특히 $s \to \rho$이면 $\kappa(s) \to \infty$.

  \item[\textup{(ii)}] \textbf{영점으로부터 먼 곳에서 유계.}
    $\mathrm{dist}(K, \{\rho_k\}) \geq \delta > 0$인 임의의
    콤팩트 집합 $K \subset \Omega$에서 $\kappa$는 유계:
    \begin{equation}
      \sup_{s \in K} \kappa(s) \leq C(K, \delta) < \infty.
      \label{eq:kappa_bounded}
    \end{equation}

  \item[\textup{(iii)}] \textbf{집중비.} 단순 영점 $\rho$에서
    오프셋 $\delta$에서의 곡률은
    \begin{equation}
      \frac{\kappa(\rho + \delta)}{\kappa(s_{\mathrm{far}})}
      \;\sim\; \frac{|s_{\mathrm{far}} - \rho_{\mathrm{nearest}}|^2}{\delta^2}
      \label{eq:concentration_ratio}
    \end{equation}
    여기서 $s_{\mathrm{far}}$는 최근접 영점이 거리
    $|s_{\mathrm{far}} - \rho_{\mathrm{nearest}}|$에 있는 임의의 점이다.
\end{enumerate}
\end{theorem}

\begin{proof}
\textbf{(i)} $m$차 영점 근방에서
$\Lambda(s) = (s-\rho)^m h(s)$로 쓰면 $h$는 정칙이고
$h(\rho) \neq 0$이다. 그러면:
\[
  \frac{\Lambda'(s)}{\Lambda(s)}
  = \frac{m}{s - \rho} + \frac{h'(s)}{h(s)}.
\]
둘째 항 $h'/h$는 $\rho$에서 정칙($h(\rho)\neq 0$이므로)이다. 따라서:
\[
  \kappa(s) = \left|\frac{m}{s-\rho} + \frac{h'(s)}{h(s)}\right|^2
  = \frac{m^2}{|s-\rho|^2} + O(|s-\rho|^{-1}).
\]

\textbf{(ii)} $K$에서 $\Lambda'/\Lambda$는 정칙(모든 영점을
피하므로)이고 따라서 콤팩트 집합 위에서 연속이므로 유계이다.

\textbf{(iii)} (i)에 의해 $\kappa(\rho+\delta) \approx 1/\delta^2$.
$s_{\mathrm{far}}$에서 아다마르 합 $\sum_k 1/(s-\rho_k)$의
지배항은 거리 $d$의 최근접 영점에서 오므로
$\kappa(s_{\mathrm{far}}) \approx 1/d^2$. 비가 따른다.
\end{proof}

\begin{remark}[수치 검증]
\label{rem:curvature_numerical}
참양성($\delta \approx 0.03$)과 거짓양성(전형적 $d \approx 0.5$)
사이에 관측된 $\sim 300\times$ 곡률비는 \eqref{eq:concentration_ratio}로
예측된다: $d^2/\delta^2 \approx (0.5/0.03)^2 \approx 278$.
\end{remark}

\begin{proposition}[임계선 영점에서의 정확한 로랑 전개]
\label{prop:kappa_exact_expansion}
$\xi(s) = \xi(1-s)$ (함수방정식)과 슈바르츠 반사
$\xi(\bar{s}) = \overline{\xi(s)}$를 만족하는 $\xi$-함수의 임계선 위 단순 영점
$s_0 = \tfrac{1}{2} + it_0$에서, 실수 오프셋 $\delta > 0$에 대해
\begin{equation}
  \frac{\xi'}{\xi}(s_0 + \delta)
  \;=\; \frac{1}{\delta} + iB(t_0) + H_1(t_0)\,\delta + O(\delta^2),
  \label{eq:conn_expansion}
\end{equation}
여기서 $B(t_0) \in \mathbb{R}$, $H_1(t_0) \in \mathbb{R}$이다.
따라서,
\begin{equation}
  \kappa(s_0 + \delta)
  \;=\; \frac{1}{\delta^2} + A(t_0) + O(\delta^2),
  \qquad
  A(t_0) := B(t_0)^2 + 2H_1(t_0) \in \mathbb{R}.
  \label{eq:kappa_exact}
\end{equation}
특히, $\kappa$에서 $O(1/\delta)$ 항의 계수는 항등적으로 0이다.
\end{proposition}

\begin{proof}
$f(\delta) := (\xi'/\xi)(s_0+\delta) = 1/\delta + c_0 + c_1\delta + O(\delta^2)$
로 쓴다. 함수방정식 $\xi(s) = \xi(1-s)$에서 $(\xi'/\xi)(1-s) = -(\xi'/\xi)(s)$이고,
$s = s_0 + \delta$, $1-s = s_0 - \delta$ ($1-s_0 = s_0$ 이용)에서
$f(-\delta) = -f(\delta)$, 즉 $c_0 = -c_0$이므로
$c_0 \in i\mathbb{R}$: $c_0 = iB$. 슈바르츠 반사
$\xi(\bar{s}) = \overline{\xi(s)}$에서 실수 $\delta$에 대해 $c_1 = \bar{c}_1$이므로
$c_1 =: H_1 \in \mathbb{R}$. 따라서
$\kappa = |f|^2 = (1/\delta + H_1\delta)^2 + B^2 = 1/\delta^2 + (B^2 + 2H_1) + O(\delta^2)$.
\end{proof}

\begin{remark}[명제~\ref{prop:kappa_exact_expansion}의 수치 검증 — 결과 \#60]
\label{rem:kappa_scaling_numerical}
결과~\#60은 $\zeta(s)$에 대해 식~\eqref{eq:kappa_exact}를 수치적으로 검증한다:
13개 영점 ($t_0 \in [10,60]$), 5개 오프셋 $\delta \in \{0.01, 0.02, 0.03, 0.05, 0.10\}$,
총 65개 데이터점에서 $\kappa\cdot\delta^2 \in [1.00012, 1.01206]$
(모든 $\delta$에서 $[0.95, 1.10]$ 이내; CV${<}1\%$).
$O(1)$ 보정항 $A(t_0) = \kappa - 1/\delta^2$은 $O(\delta^2)$ 오차 이내에서
$\delta$에 무관함이 확인되었다.

\smallskip
$\delta$-스케일링 요약 (13개 영점 중앙값):

\smallskip
\begin{center}
\small
\begin{tabular}{@{}rrrrl@{}}
\toprule
$\delta$ & $1/\delta^2$ & $\kappa_{\mathrm{near}}$ (중앙값) & CV & $\kappa\cdot\delta^2$ \\
\midrule
0.01 & 10000 & 10001.21 & $0.01\%$ & 1.00012 \\
0.02 &  2500 &  2501.21 & $0.03\%$ & 1.00048 \\
0.03 &  1111 &  1112.32 & $0.06\%$ & 1.00109 \\
0.05 &   400 &   401.21 & $0.16\%$ & 1.00301 \\
0.10 &   100 &   101.21 & $0.62\%$ & 1.01206 \\
\bottomrule
\end{tabular}
\end{center}
\smallskip\noindent
\textcolor{ForestGreen}{[결과 \#60, 확립됨, 2026-04-18; 13영점, 5오프셋, 65 데이터점;
$\kappa{\cdot}\delta^2{\in}[1.00012,1.01206]$; $O(1/\delta)$ 계수 $= 0$ 확인; 수치 검증, 증명 아님]}
\end{remark}

\begin{remark}[$A(t_0)$의 차수 비교: 3-차수 스케일링 (결과~\#61)]
\label{rem:kappa_degree_comparison}
스케일링 법칙 $\kappa = 1/\delta^2 + A(t_0) + O(\delta^2)$
(명제~\ref{prop:kappa_exact_expansion})을 차수 1--3에 걸쳐 120개의 추가 데이터점으로 검증한다:
GL(2) Maass 형식 ($R \approx 13.78$): 12영점 $\times$ 5오프셋 $= 60$점,
GL(3) 대칭 제곱 $L(s, \mathrm{sym}^2(11\mathrm{a}1))$: 12영점 $\times$ 5오프셋 $= 60$점.
보정항 $A(t_0) = \kappa - 1/\delta^2$은 $\delta$에 무관함이 확인됨
(GL(2) spread $< 0.21\%$, GL(3) spread $< 0.70\%$).

차수별 평균$A(t_0)$ 정량 비교:
\smallskip
\begin{center}
\small
\begin{tabular}{@{}llrrl@{}}
\toprule
차수 & $L$-함수 & 평균$(A)$ & 범위$(A)$ & $n$ \\
\midrule
GL(1) & $\zeta(s)$ & 1.30 & $[0.47,\, 2.73]$ & 13 \\
GL(2) & Maass $R{\approx}13.78$ & 3.93 & $[0.50,\, 6.67]$ & 12 \\
GL(3) & $\mathrm{sym}^2(11\mathrm{a}1)$ & 12.79 & $[6.64,\, 23.35]$ & 12 \\
\bottomrule
\end{tabular}
\end{center}
\smallskip\noindent
평균$(A)$는 차수와 함께 단조 증가: $1.30 \to 3.93 \to 12.79$
(비율: GL(2)/GL(1) $\approx 3.0\times$, GL(3)/GL(2) $\approx 3.3\times$).
이 증가는 GL($d$)의 함수방정식 $\Lambda$에 포함된 $d$개 감마 인자의 누적 기여에서 기인한다:
각 감마 인자는 영점 $\rho$에서의 $H_0 = (\Lambda'/\Lambda)$에 $\psi$-함수(디감마) 항을 기여하므로
$|\mathrm{Im}(H_0)|^2$가 $d$와 함께 증가한다.
개별 영점 수준에서는 범위가 약간 겹침 (GL(2) 최대값 $6.67 >$ GL(3) 최솟값 $6.64$)에 유의.
3개의 데이터점만으로는 $A$를 $d$의 함수로 결정할 수 없음.

주의: $\kappa \cdot \delta^2 \approx 1$은 $\xi'/\xi$의 단순 극 구조에서 수학적으로 필연적인 귀결이며
독립적 발견이 아니다. 결과의 핵심 가치는 $O(1)$ 보정항 $A(t_0)$의 \emph{정량적 차수 의존성}:
$1/\delta^2$가 담지 못하는 차수- 및 영점-특유의 산술적 정보를 $A$가 인코딩함.
$\delta = 0.1$에서 GL(3) 보정이 $A \cdot \delta^2 \approx 0.12$ (12\%)에 달하므로,
고차 차수에서는 더 작은 $\delta$가 권장됨.
해석적 분해 $A(t_0) = A_\Gamma + A_L$ (감마 인자 기여와 유한 오일러 곱 기여의 분리)는
미해결 문제로 남아 있다 (경계 B-18; \S\ref{sec:open} 참조).

\textcolor{ForestGreen}{[결과 \#61, 확립됨, 2026-04-18;
GL(2): 12영점, 5오프셋, 60점, $\kappa{\cdot}\delta^2{\in}[1.00,1.04]$;
GL(3): 12영점, 5오프셋, 60점, $\kappa{\cdot}\delta^2{\in}[1.00,1.12]$;
평균$(A)$ 단조 증가 GL(1--3): 1.30, 3.93, 12.79;
$A(d)$ 함수형 3점으로 미결정; 수치 검증, 증명 아님]}
\end{remark}

\begin{remark}[차수~1 내 $A(t_0)$의 도체 의존성: 디리클레 확장 (결과~\#62)]
\label{rem:dirichlet_conductor_A}
스케일링 법칙 $\kappa = 1/\delta^2 + A(t_0) + O(\delta^2)$는
디리클레 $L$-함수 $L(s,\chi)$의 $\chi\bmod 3$, $\chi\bmod 4$, $\chi\bmod 5$
(모두 홀수 문자, $\mu = 1$)로 확장된다.
총 330개 데이터점(66영점 $\times$ 5오프셋
$\delta \in \{0.001, 0.005, 0.01, 0.05, 0.1\}$)이 확인:
\begin{itemize}
  \item 모든 문자·모든 $\delta$에서 $\kappa\cdot\delta^2 \in [1.00000, 1.028]$ (통과);
  \item $A(t_0)$는 $\delta$에 무관 (spread $< 0.1\%$, 세 $\chi$ 모두).
\end{itemize}
보정항 $A(t_0)$는 차수~1 내에서 도체(conductor)에 대한 명확한 의존성을 보인다:
\smallskip
\begin{center}
\small
\begin{tabular}{@{}llcrrl@{}}
\toprule
$L$-함수 & $\mu$ & 도체 $q$ & 평균$(A)$ & $n$ & spread \\
\midrule
$\zeta(s)$ & 0 & 1 & 1.27 & 12 & 0.00\% \\
$L(s,\chi_3)$ & 1 & 3 & 2.25 & 20 & 0.08\% \\
$L(s,\chi_4)$ & 1 & 4 & 2.77 & 22 & 0.02\% \\
$L(s,\chi_5)$ & 1 & 5 & 3.14 & 24 & 0.06\% \\
\bottomrule
\end{tabular}
\end{center}
\smallskip\noindent
동일 $\delta$에서 $\kappa_{\mathrm{near}}$는 모든 차수-1 $L$-함수 사이에서
$< 0.2\%$ 차이로, $\kappa_{\mathrm{near}}$가 실질적으로 차수만의 함수임을 확인한다
(경계~B-09).
도체 의존적 변동은 전적으로 $O(1)$ 보정항 $A(t_0)$에 존재하며,
$q$에 대해 단조 증가한다:
$A(\zeta) = 1.27 < A(\chi_3) = 2.25 < A(\chi_4) = 2.77 < A(\chi_5) = 3.14$.

단순 추정 $\Delta A \approx \log(q/\pi)/2$는 관측된 증가의 ${\sim}13\%$만을 설명하며,
주요 기여는 영점에서의 $L'/L$ 항에서 나온다
(도체 의존적 영점 밀도 반영).
이 세 문자에서 $\mu$와 $q$가 동시에 변하므로 (모두 $\mu = 1$, $q > 1$),
$\mu$-시프트 효과와 도체 효과의 분리는 짝수 문자($\mu = 0$, $q > 1$)가 필요하며
실험~\#76에서 다룬다.

\textcolor{ForestGreen}{[결과 \#62, 확립됨, 2026-04-18;
66영점, 5오프셋, 330데이터점;
$\kappa{\cdot}\delta^2$ 전점 통과; $A$ $\delta$-독립 (spread ${<}0.1\%$);
$A$ 도체-단조: 1.27, 2.25, 2.77, 3.14;
동일-$\delta$ $\kappa_{\mathrm{near}}$ 차이 ${<}0.2\%$ (B-09 확인);
수치 증거, 증명 아님]}
\end{remark}


% ============================================================================
\section{가우스-보네 양자화}
\label{sec:gauss_bonnet_thm}
% ============================================================================

\begin{definition}[곡률 2-형식]
\label{def:curvature_2form}
$\theta(\sigma,t) := \mathrm{Im}\,\log\Lambda(\sigma + it)$에 대해
$\xif$-다발의 곡률 2-형식은
\begin{equation}
  F := d\omega = \left(\frac{\partial^2\theta}{\partial\sigma\,\partial t}
  - \frac{\partial^2\theta}{\partial t\,\partial\sigma}\right)
  d\sigma \wedge dt
  \label{eq:curvature_2form}
\end{equation}
이다. 영점으로부터 먼 곳에서 $\theta$는 매끈하고 $F = 0$(평탄 접속).
영점에서 $\theta$는 대수적 특이점을 가져
$\delta$-함수 플럭스에 기여한다.
\end{definition}

\begin{theorem}[$\xif$-다발에 대한 가우스-보네]
\label{thm:gauss_bonnet}
$R = [\sigma_1,\sigma_2] \times [t_1, t_2]$를 경계가 $\Lambda$의
영점을 피하는 임계 띠 내 직사각형이라 하자. 그러면:
\begin{equation}
  \frac{1}{2\pi}\oint_{\partial R} d\,\mathrm{Im}\,\log\Lambda
  = \frac{1}{2\pi}\oint_{\partial R} \mathrm{Im}\!\left(
  \frac{\Lambda'(s)}{\Lambda(s)}\,ds\right)
  = N(R),
  \label{eq:gauss_bonnet}
\end{equation}
여기서 $N(R)$은 $R$ 내부 $\Lambda$의 영점 수(중복도 포함).
동등하게, 격자 정규화 형태로:
\begin{equation}
  \frac{1}{2\pi}\sum_{\mathrm{plaquettes}\;\square\,\subset\,R}
  \Delta\theta_{\square} = N(R).
  \label{eq:gauss_bonnet_lattice}
\end{equation}
\end{theorem}

\begin{proof}
편각 원리(정리~\ref{thm:holonomy}(iii))에 의해:
\[
  \frac{1}{2\pi i}\oint_{\partial R}
  \frac{\Lambda'(s)}{\Lambda(s)}\,ds = N(R).
\]
$\partial R$을 따른 $\arg\Lambda$의 총 변화량은 $2\pi N(R)$이다.
이것이 $R$에 제한된 U(1)-다발의 \textbf{제1 체른 수}이다:
$c_1(\xif\text{-다발}|_R) = N(R)$.

격자 형태: $R$을 작은 직사각형(플라켓)으로 분할한다.
영점을 포함하지 않는 플라켓에서 감긴 위상 순환
$\Delta\theta_\square = 0$. $m$차 영점을 포함하는 플라켓에서
편각 원리에 의해 $\Delta\theta_\square = 2\pi m$.
모든 플라켓에 걸쳐 합산하고 $2\pi$로 나누면 $N(R)$을 얻는다.
격자 합은 \emph{위상적}이다: 격자 세분에 대해 불변이며,
오직 포함된 영점의 총 수에만 의존한다.
\end{proof}

\begin{corollary}[정수 양자화]
\label{cor:integer_quant}
$\frac{1}{2\pi}\oint_{\partial R} d\arg\Lambda$는 경계가 영점을
피하는 \emph{임의의} 직사각형 $R$에 대해 정확히 정수이다.
\end{corollary}


% ============================================================================
\section{교차 차수 보편성}
\label{sec:universality_thm}
% ============================================================================

\begin{theorem}[교차 차수 보편성]
\label{thm:universality}
정리~\ref{thm:holonomy}, \ref{thm:curvature},
\ref{thm:gauss_bonnet}은 $\mathbb{Q}$ 위 $\mathrm{GL}(n)$의
\emph{임의의} 보형 표현 $\pi$에 대응하는 완비 $L$-함수
$\Lambda(s, \pi)$에 대해 성립한다. 단, $\Lambda(s,\pi)$가:
\begin{enumerate}
  \item[\textup{(a)}] 고립 영점을 갖는 유한 위수 유리형 함수이고,
  \item[\textup{(b)}] 영점에 대한 아다마르형 곱을 허용하면 된다.
\end{enumerate}
특히 홀로노미, 곡률 집중, 가우스-보네 양자화는
차수 $n$, 도체 $N$, 가중치 $k$, 근수 $\varepsilon$에
\textbf{무관}하다.
\end{theorem}

\begin{proof}
정리~\ref{thm:holonomy}--\ref{thm:gauss_bonnet}의 증명은
오직 (1) 편각 원리($\Lambda'/\Lambda$가 영점에서 단순극을 갖는
유리형)와 (2) 로랑 전개 $\Lambda'/\Lambda = m/(s-\rho) + O(1)$
(조건 (b)로부터)만을 사용한다. 두 성질 모두 (a)--(b)를 만족하는
모든 보형 $L$-함수에 성립한다.
\end{proof}

\begin{remark}[수치 확인]
\label{rem:universality_numerical}
교차 차수 보편성은 GL(1)(리만 $\zeta$ + 디리클레 5종)과
GL(2)(11a1, 37a1, 라마누잔 $\Delta$)에서 모두 확인된다.
\emph{FP 모노드로미 교차 순위 검증:} GL(2) 3개 $L$-함수에서
34개 TP mono$/\pi = 2.000 \pm 0.000$,
60개 FP mono$/\pi = 0.000 \pm 0.000$
(Mann--Whitney $p = 9.7 \times 10^{-17}$)으로
GL(1)과 동일한 분리도를 확인.
차수에 따라 변하는 \emph{유일한} 진단은 $\sigma$-유일성으로,
이는 편각 원리의 귀결이 아니라 구체적 $\Gamma$-인자 구조에
의존하기 때문이다(부록~D 참조).
\end{remark}

\begin{remark}[GL(2) 블라인드 영점 예측 --- 교차 순위 보편성]
\label{rem:gl2_blind_prediction}
교차 순위 보편성은 세 GL(2) $L$-함수 및 하나의 GL(1) 디리클레 $L$-함수에 대한
\emph{블라인드} 영점 예측 실험을 통해 추가로 뒷받침된다.
여기서 \emph{블라인드}란 LMFDB 영점 위치가 최종 평가 단계에서만
사용됨을 의미한다: 곡률 스윕 및 $\kappa$-피크 탐지 단계에서
영점 위치를 사전에 참조하지 않는다.

\textbf{실험 \#52 --- 타원곡선 11a1}
(도체 $N=11$, rank~0, 근호 $\varepsilon=+1$):
$t\in[5,30]$의 곡률 스윕 ($\Delta t = 0.1$, \texttt{dps}$=45$)에서
17개의 $\kappa$ 피크를 후보 영점으로 식별한다.
17개의 LMFDB 영점과 매칭: 16개 TP, 1개 FP, 1개 미탐
($\gamma_{17}=29.975$, 스윕 범위 외).
$P = R = F_1 = 0.941$. 위치 오차 최대 $= 0.051$.
유일한 FP ($t \approx 28.4$)는 이중 피크 잔상이다:
윤곽 반경 $r = 0.135$에 LMFDB 영점 미포함 (최근접 거리 $0.284>r$),
물리적 권선수는 0이어야 하며 mono$/\pi = 2.000$은 \texttt{dps}$=45$의 수치 잔상.

\textbf{실험 \#54 --- 타원곡선 37a1}
(도체 $N=37$, rank~1, 근호 $\varepsilon=-1$):
$t\in[5,30]$의 곡률 스윕 ($\Delta t = 0.2$, \texttt{dps}$=60$)에서
22개의 $\kappa$ 피크를 식별한다.
23개의 LMFDB 영점과 매칭: 22개 TP, 0개 FP, 1개 미탐
($\gamma_{24}=29.282$, $\gamma_{23}=29.030$와 간격 $0.252 < 1.5\,\Delta t$).
$P = 1.000$, $R = 0.957$, $F_1 = 0.978$. 위치 오차 최대 $= 0.077$.
FP 0개는 정밀도 향상(\texttt{dps}$=60$)에 의한 결과로,
\#52의 수치 잔상을 제거한다.

\textbf{$t\approx 29.0$에서의 정량적 모노드로미.}
$t = 29.0$에서 반경 $r = 0.36$의 모노드로미 윤곽은
\emph{두 개}의 영점($\gamma_{23}=29.030$, $\gamma_{24}=29.282$)을 포함하여
권선수~2, mono$/\pi = 4.0$을 산출한다.
이는 모노드로미 지표가 영점 유무의 이진 정보를 넘어
윤곽 내 영점 \emph{개수}를 정확히 보고함을 실증한다:
$\Delta t = 0.2$ 해상도 한계로 분리 불가한 인접 영점 쌍을
모노드로미가 보완한다.

\textbf{실험 \#56 --- 라마누잔 $\Delta$ (weight~12)}
(level $N=1$, 근호 $\varepsilon=+1$, 임계선 $\sigma_c = k/2 = 6$):
$t\in[5,30]$의 곡률 스윕 ($\Delta t = 0.2$, \texttt{dps}$=60$)에서
8개의 $\kappa$ 피크를 후보 영점으로 식별한다.
8개의 LMFDB 영점과 매칭: 8개 TP, 0개 FP, 0개 FN --- \emph{완전 탐지}.
$P = R = F_1 = 1.000$. 위치 오차 최대 $= 0.108$.
8개 후보 전부 mono$/\pi = 2.000$ (권선수~1, 단순 영점).
$\kappa$ near/far 비율 $322.8\times$ (near 중앙값 $= 343.5$, far 중앙값 $= 1.06$).
GL(2) 블라인드 3종 중 유일한 완벽 결과:
임계선이 $\sigma_c = 1$ (타원곡선)에서 $\sigma_c = 6$ (weight~12)으로
이동함에도 $\xi$-다발 프레임워크는 오탐·미탐 모두 0을 달성한다.

\textbf{실험 \#57 --- 디리클레 $\boldsymbol{\chi\bmod 7}$}
(위수~6 복소 지표, 도체 $q=7$, $a=1$, $\sigma_c = 1/2$):
$t\in[5,40]$의 곡률 스윕 ($\Delta t = 0.2$, \texttt{dps}$=80$)에서
18개의 $\kappa$ 피크를 후보 영점으로 식별한다.
13개의 LMFDB 영점과 매칭: 13개 TP, 5개 FP, 0개 FN.
$R = 1.000$, $P = 0.722$, $F_1 = 0.839$. 위치 오차 최대 $= 0.084$.
5개 FP는 모두 $t < 17.16$ (첫 LMFDB 영점 $\gamma_1 = 17.161$ 이하)에 발생한다;
독립 평가에서 이 위치에서 $|L(\tfrac{1}{2}+it, \chi)| \gg 0$이 확인되어
임계선 위 영점이 아님이 확정되었다.
FP 원인: 저~$t$ 구간에서 $\Gamma$-배경이 증가하여, 반경 $r = 0.4$
모노드로미 디스크 내의 임계선 밖 영점 포획.
$t \geq 17$ 제한 시: $P = R = F_1 = 1.000$ (13/13 TP, 0 FP).
$\kappa$ near/far 비율 $24.0\times$는 GL(2) 곡선보다 낮으나,
도체에 의한 배경 증가로 비-블라인드 결과(\#37; $36.6\times$)와 일관적이다.

표~\ref{tab:gl1_gl2_blind}는 다섯 $L$-함수에 걸친 블라인드 예측을 비교한다
(GL(1) 및 GL(2), weight~1, 2, 12, rank~0 및~1,
네 도체 ($1, 7, 11, 37$), 양쪽 근호에 걸쳐 있음).

\begin{table}[ht]
\centering
\caption{다섯 $L$-함수에 걸친 블라인드 영점 예측.
  \emph{블라인드}: 평가 단계에서만 LMFDB 영점 사용.
  $\varepsilon$: 근호; $k$: weight; $\sigma_c$: 임계선.}
\label{tab:gl1_gl2_blind}
\small
\begin{tabular}{lrrrrrrrrc}
\toprule
실험 & $N$ & $k$ & rank & $\varepsilon$ & 영점 수 & TP & FP & 정밀도 & 재현율 \\
\midrule
\#2\ \ GL(1) $\zeta$ (블라인드) & 1 & 1 & --- & $+1$ & 7 & 7 & 0 & 1.000 & 1.000 \\
\#52 GL(2) 11a1 (블라인드) & 11 & 2 & 0 & $+1$ & 17 & 16 & 1 & 0.941 & 0.941 \\
\#54 GL(2) 37a1 (블라인드) & 37 & 2 & 1 & $-1$ & 23 & 22 & 0 & 1.000 & 0.957 \\
\#56 GL(2) $\Delta$ (블라인드) & 1 & 12 & --- & $+1$ & 8 & 8 & 0 & 1.000 & 1.000 \\
\#57 GL(1) $\chi\bmod 7$ (블라인드) & 7 & 1 & --- & $-1$ & 13 & 13 & 5 & 0.722 & 1.000 \\
\midrule
\multicolumn{5}{l}{\textbf{합산 (5곡선)}} & \textbf{68} & \textbf{66} & \textbf{6} & \textbf{0.917} & \textbf{0.971} \\
\bottomrule
\end{tabular}
\end{table}

다섯 곡선 합산: 66/68 TP, 6 FP, $F_1 = 0.943$.
다섯 중 네 곡선이 $P \geq 0.94$를 달성하며, 예외인
$\chi\bmod 7$ ($P = 0.722$)의 5개 FP는 $t < 17$ (첫 LMFDB 영점 이전)에
집중되어 저~$t$ 모노드로미 디스크 내 임계선 밖 영점 포획에서 비롯된다.
$t \geq 17$에서는 다섯 곡선 모두 합산 66개 영점에 대해
$P = R = F_1 = 1.000$이다.
라마누잔~$\Delta$ (\#56)는 $\sigma_c = 6$에서 \emph{완벽한}
$P = R = F_1 = 1.000$을, $\chi\bmod 7$ (\#57)는
복소 디리클레 지표에서 완벽한 재현율 ($R = 1.000$, 13/13 TP)을 달성한다.
다섯 실험은 GL(1) 리만, GL(1) 디리클레(복소 지표),
GL(2) (rank~0/1 타원곡선 + weight~12 모듈러 형식)을 포괄하여
테스트된 모든 표준 자기동형 $L$-함수 패밀리를 커버한다.
\emph{이는 자기동형 패밀리에 걸친 블라인드 영점 예측 보편성에 대한
수치적 증거이며, 증명이 아니다.}
\end{remark}


% ============================================================================
\section{이중 기준: 모노드로미 $+$ 곡률}
\label{sec:dual_criterion_thm}
% ============================================================================

\begin{corollary}[정확한 TP/FP 분리]
\label{cor:dual_criterion}
후보 영점 위치 $\{t_1, \ldots, t_N\}$ 각각에 대해:
\begin{align}
  \mu_j &:= |\mathrm{Hol}(\mathcal{C}_r(\tfrac{1}{2}+it_j))|,
  \label{eq:mono_criterion}\\
  \kappa_j &:= \kappa(\tfrac{1}{2} + \delta + it_j).
  \label{eq:kappa_criterion}
\end{align}
\begin{enumerate}
  \item[\textup{(i)}] $t_j$가 참 영점일 필요충분조건은 $\mu_j = 1$
    (단순 영점 가정, 충분히 작은 $r$).
  \item[\textup{(ii)}] 참 영점이면 $\kappa_j \sim 1/\delta^2 \gg 1$.
    거짓양성이면 $\kappa_j \sim 1/d^2 \ll 1/\delta^2$.
  \item[\textup{(iii)}] 이중 기준 $\{\mu_j = 1\} \cap \{\kappa_j > \tau\}$는
    적절한 $\tau$에 대해 \textbf{완벽한 분리}를 달성한다.
\end{enumerate}
\end{corollary}

\begin{proof}
(i)은 정리~\ref{thm:holonomy}(i)--(ii), (ii)는
정리~\ref{thm:curvature}(i)과 (iii), (iii)은 모노드로미 기준만으로
이미 정확하고 곡률 기준이 독립적 연속값 확인을 제공하므로 따른다.
\end{proof}

\begin{remark}[추측 3과의 관계]
\label{rem:conj3}
이전에 ``추측~3''이라 명명된 것(이중 기준이 TP를 FP로부터 분리)은
이제 \textbf{정리}임이 밝혀진다: 편각 원리(정리~\ref{thm:holonomy})와
$\Lambda'/\Lambda$의 로랑 구조(정리~\ref{thm:curvature})의 직접적
귀결이다. 수치적으로는 GL(1) $\zeta$(10 TP, 19 FP, $p = 6.2 \times 10^{-6}$)와
GL(2) 3개 $L$-함수(34 TP, 60 FP, $p = 9.7 \times 10^{-17}$) 양쪽에서
이중 기준의 완벽 분리가 확인되었다.
\end{remark}

\begin{remark}[여전히 추측인 것]
\label{rem:what_conjectural}
증명된 결과(정리~\ref{thm:holonomy}--\ref{thm:universality},
따름정리~\ref{cor:dual_criterion})는 \emph{임의의 위치}의 영점에서
다발을 특성화한다. 영점이 \emph{어디에} 놓이는지는 제약하지 않는다.
모든 영점이 $\sigma = 1/2$ 위에 있다는 명제(리만 가설)는
여기서 확립된 다발 구조와 독립적인 추측으로 남는다. 다만
$\sigma$-국소화(FWHM $< 0.004$)는 위상적 전하가 임계선에
집중됨을 보여주는 강력한 수치적 증거이다.
\end{remark}


% [Paper A: ξ-다발 기하학]
% ############################################################################
% 제3부 — 실패한 방향과 그 진단
% ############################################################################

\part{실패한 방향과 그 진단}
\label{part:failed}

\noindent 본 부는 성공하지 못한 경험 프로그램을 압축된 형태로 모은다. 각 섹션은 \textcolor{red}{\rule{0.9ex}{0.9ex}}\ (계층 3)로 표시되며, 각각 프로그램의 한 단락 설명으로 시작하여 진단 사슬을 거쳐 추출된 교훈으로 닫힌다. 내용은 의도적으로 조밀하다 --- 실패 방향당 전형적으로 1--2 페이지 --- 이는 검증된 자료와 추측 자료를 밀어내지 않으면서 완전한 진단을 유지하기 위함이다. 서두른 독자는 제~\ref{part:failed}부를 건너뛰어도 되지만, 후속 실험을 계획하는 독자는 주의 깊게 읽어야 한다.

% ============================================================================
\section{\textcolor{red}{\rule{0.9ex}{0.9ex}}\ 시클로토믹 $\cos^2$ PQO 학습 ($L=2,4$): $\Psip$-디커플링과 고립 캐스케이드}
\label{sec:cos2_failure}
% ============================================================================

\textbf{프로그램.} 3층 공명 판별 네트워크를 \cref{rem:cos2_pqo}의 시클로토믹 $L\in\{2,4\}$ 코사인 제곱 양자화 손실 $\Lpqo$와 \cref{eq:psi_n}의 합성 크기 전용 타겟 $\Psi_n$을 함께 사용하여 $t\in[100,200]$에서 리만 영점과 비영점 좌표를 분리하도록 학습시킨다.

\textbf{결과.} $L=2$와 $L=4$ 양쪽에서 음성. $L=2$에서 네트워크는 자명한 어트랙터 $\phi\to\pi/2$로 붕괴하여 영점과 비영점 어느 쪽도 판별하지 못한다. $L=4$에서 학습은 어트랙터 붕괴 없이 수렴하지만, 결과 위상 MSE (대 $\arg Z(t)$)는 영점과 비영점 간에 구분이 불가능하다 (세 시드에서 풀링된 Welch $p>0.12$, 시클로토믹 구간 $k=2,3$ 공란). 완전한 표 결과는 \cref{rem:cos2_pqo} 참조.

\textbf{진단.} 실패를 네 단계로 추적한다 (날짜는 실험 일정이 아니라 진단 사슬을 나타낸다):
\begin{enumerate}
  \item \textbf{타겟 디커플링 (\cref{rem:psi_defect}, 2026-04-08).} \cref{eq:psi_n}의 합성 타겟 $\Psi_n$은 $\|X\|$ (또는 $n$)에만 의존하여 $t$-국소 리만 영점 정보를 전혀 운반하지 않는다. $(\alpha,\beta,\gamma,\delta)$의 어떤 선택으로도 이를 복구할 수 없다; 결함은 구조적이다.
  \item \textbf{감독 교체는 필요하지만 충분하지 않다 (\cref{rem:psi_defect}, 2026-04-08).} 타겟을 해석적 정답 $\Psip(t)=\arg\xif(\tfrac{1}{2}+it)$로 교체하고 $\Lpqo$를 학습에서 제거하면 (측정에서만 유지), $L=4$의 비대칭 분지 붕괴가 해제되지만 헤드라인 분리 지표는 영가설 근방에 남는다 ($p=0.976$, 비율 $1.004$).
  \item \textbf{손실 가중 불균형과 산술 평균 읽기는 추가적인 끊긴 고리이다 (\cref{rem:psi_defect}, 2026-04-08).} 잔차 손실 가중 $\lambda_{\mathrm{res}}=1$, $\lambda_{\mathrm{curv}}=0.1$이 감독 $\mathcal{L}_{\mathrm{tgt}}$를 압도하고, $\hat\Psip$을 읽어내기 위해 사용된 채널별 산술 평균이 임계선 위 위상 점프 구조가 예측하는 $\pm\pi$ 랩 신호를 파괴한다.
  \item \textbf{고립이 사슬을 확인한다 (\cref{rem:isolation}, 2026-04-09).} $\lambda_{\mathrm{res}}=\lambda_{\mathrm{curv}}=\lambda_{\mathrm{pqo}}=0$으로 설정하고 원형 평균으로 읽어내면, 비영점 집합에서 $p=1.18\times10^{-15}$로 감독 학습이 회복되며, 영점 잔존 MSE $\approx(\pi/2)^2$은 매끄러운 네트워크가 각 영점에서 $\arg\xif(\tfrac{1}{2}+it)$의 $\pm\pi$ 불연속을 표현할 수 없음으로 완전히 설명된다.
\end{enumerate}

\textbf{교훈.} \emph{네 개의 독립적 실패 모드가 쌓여 하나의 음성 결과를 만들어낼 수 있다.} 해석적 동일시가 실험적으로 반증되었다고 선언하기 전에, 감독 사슬의 모든 고리가 개별적으로 격리되어야 한다: 타겟 주입, 손실 가중, 측정 연산자, 출력 공간 표현력. 프로그램의 해석적 내용 $\Psip(t)=\arg\xif(\tfrac{1}{2}+it)$은 캐스케이드를 살아남는다; 매끄러운 장 출력 표현이 마지막 장애물이었으며, 2026-04-10에 $\{0,\pi\}$ 계단 타겟을 정규화된 매끄러운 대리 $\mathrm{Re}\,\xif/\max|\mathrm{Re}\,\xif|$로 교체하여 \emph{해소}되었다: 영점/비영점 MSE 비율이 $18.2\times$에서 $0.03\times$로 감소한다 (\cref{rem:isolation}). 복소 벡터 읽기 헤드는 미시험 대안으로 남는다 (제~\ref{sec:open}절).

% ============================================================================
\section{\textcolor{red}{\rule{0.9ex}{0.9ex}}\ 위상-게이지 경사하강 (PGGD)}
\label{sec:pggd_failure}
% ============================================================================

\textbf{프로그램.} PGGD는 판별 네트워크의 중간 상태의 $S^1$-값 위상 구조를 활용하기 위해 도입된 사용자 정의 옵티마이저이다. 각 매개변수 업데이트는 주변 원 다발의 접공간으로 투영되어 최적화 궤적이 게이지 대칭 $\phi\mapsto\phi+\mathrm{const}$를 존중하도록 한다. 가설은 PGGD가 허위의 전역 위상 자유도를 제거함으로써 제~\ref{sec:f2_hardy}절의 $F_2$ 기반 영점 탐지 과제에서 Adam을 능가할 것이라는 것이었다.

\textbf{결과.} 시험된 모든 시드에서 Adam 대비 음성. PGGD는 (i) 동일 epoch 예산 내에 수렴하지 못하거나, (ii) Adam 기준선보다 작은 판별 격차로 더 높은 학습 손실에 수렴한다. $t\in[10,50]$ 벤치마크에서 Adam의 최종 $F_2$-MSE는 일관되게 PGGD 최상 실행의 $2$--$4\times$ 작다.

\textbf{진단.} 관련된 두 메커니즘:
\begin{enumerate}
  \item \textbf{$S^1$ 위의 에르고딕 확산.} PGGD가 사용하는 접 투영은 각 업데이트의 반지름 성분을 제거하지만 각도 성분 전체를 보존한다. 다수의 스텝에 걸쳐 각도 업데이트는 실효 스텝 크기 $\eta$와 epoch 수에 선형적으로 커지는 분산으로 대략 확산적으로 누적된다. 따라서 투영된 궤적은 본질적으로 균일하게 $S^1$을 탐색하며 (각도 mod $2\pi$의 Weyl 등분포), 옵티마이저가 구현하도록 설계된 바로 그 게이지 제거를 물리친다. PGGD가 반지름 자유도를 제거해 얻는 것을 각도 자유도의 임의 보행으로 잃는다.
  \item \textbf{$\Lpqo$와의 그래디언트 스케일 불일치.} 코사인 양자화 손실은 최소값에서 멀리 있는 영역에서 $O(1)$ 그래디언트를 갖고 근방에서 $O(|\phi-\phi_*|)$ 그래디언트를 갖는다. PGGD의 접 투영은 Adam의 매개변수별 적응적 스케일링에 비해 첫 번째 영역을 증폭시켜 어트랙터 근방에서 지나치게 큰 실효 학습률을 생성한다.
\end{enumerate}

\textbf{교훈.} \emph{게이지 존중 옵티마이저가 적응적 1차 방법의 표본 효율을 자동으로 물려받지는 않는다.} 올바른 접 투영은 결맞음 탐지와 호환되지만 이를 함의하지는 않는다. PGGD 프로그램은 적응적 각도별 스텝 크기에 의한 가능한 재활 아이디어로 유지된다 (제~\ref{sec:open}절 참조). 그러나 현재 형태로는 $F_2$ 탐지기에서 Adam의 대체재로 사용되어서는 안 된다.


% [Paper A: ξ-다발 기하학]
% ############################################################################
% 제4부 — ξ-다발 수치 증거
% ############################################################################

\part{ξ-다발 수치 증거}
\label{part:synthesis}


% [Paper A: ξ-다발 기하학]
% ============================================================================
\section{\textcolor{ForestGreen}{\rule{0.9ex}{0.9ex}}\ 중간점 곡률 비국소성}
\label{sec:midpoint_nonlocal}
% ============================================================================

\noindent\textbf{Tier 1.}
연속 영점 사이의 중간점 $m = (\gamma_n + \gamma_{n+1})/2$에서,
Hadamard 곱 표현은 두 최근접(NN) 기여가 정확히 상쇄됨을 보장한다:
$1/(s-\rho_n)$ 항과 $1/(s-\rho_{n+1})$ 항이 대칭에 의해 합 0이 된다.
잔여 곡률 $\kappa(m) = |\Lambda'/\Lambda(1/2+im)|^2$는
따라서 먼 영점들의 \emph{비국소} 정보만을 인코딩한다.

\begin{observation}[중간점 곡률이 다음 간격을 인코딩]
\label{obs:midpoint_nonlocal}
$\zeta(s)$의 처음 234개 연속 영점 쌍
($t\in[0,450]$, dps\,$=$\,$100$, 해석적 $\log\xi'$ 공식,
cap률 $0.0\%$)에서 Spearman 순위 상관은 다음과 같다:
\begin{center}
\begin{tabular}{lccl}
\toprule
\textbf{쌍} & $\rho$ & $p$값 & \textbf{해석} \\
\midrule
$\kappa_{\mathrm{mid}}$ vs.\ gap ($\gamma_{n+1}-\gamma_n$)
  & $-0.031$ & $0.64$ & NN 상쇄 확인 \\
$\kappa_{\mathrm{mid}}$ vs.\ gap$_{\mathrm{next}}$ ($\gamma_{n+2}-\gamma_{n+1}$)
  & $-0.654$ & $7.3\times 10^{-30}$ & \textbf{강한 비국소 신호} \\
$\kappa_{\mathrm{mid}}$ vs.\ 국소 밀도
  & $+0.022$ & $0.74$ & 밀도 효과 없음 \\
\bottomrule
\end{tabular}
\end{center}
\noindent
편상관 $\rho(\kappa_{\mathrm{mid}},
\mathrm{gap}_{\mathrm{next}} \mid \mathrm{gap}) = -0.666$은
본질적으로 변하지 않아, 간격 자기상관이 교란 요인이 아님을 확인한다.
\textcolor{ForestGreen}{[확인 2026-04-14]}
\end{observation}

\begin{observation}[디리클레 $L$-함수에서의 보편성]
\label{obs:dirichlet_nonlocal}
중간점 비국소성은 세 디리클레 $L$-함수 ($\chi_3, \chi_4, \chi_5$)에서도
지속된다. 해석적 공식
$\Lambda'/\Lambda = \tfrac{1}{2}\log(q/\pi)
  + \tfrac{1}{2}\psi((s+a)/2) + L'/L(s,\chi)$를 사용하여,
테스트한 모든 $L$-함수에서 $|\rho(\kappa_{\mathrm{mid}},
\mathrm{gap}_{\mathrm{next}})| > 0.55$이며 $p < 10^{-10}$을 확인하였다.
\textcolor{ForestGreen}{[확인 2026-04-14]}
\end{observation}

\begin{table}[h]
\centering
\caption{중간점 곡률 비국소성: $L$-함수 vs.\ 무작위 점 과정.
  ``NN 상쇄'' 열: $|\rho(a)|<0.3$이면 최근접 상쇄가 작동함을 확인.}
\label{tab:midpoint_nonlocal}
\small
\begin{tabular}{@{} l c c c c c @{}}
\toprule
\textbf{점 과정} & \textbf{$n$ 쌍} &
  $\rho(a)$ gap & $\rho(b)$ gap$_{\text{next}}$ &
  \textbf{편상관} & \textbf{판정} \\
\midrule
$\zeta(s)$  & 234 & $-0.031$ & $\mathbf{-0.654}$ & $-0.666$ & \checkmark\ 비국소 \\
$L(s,\chi_3)$  & 112 & $-0.045$ & $\mathbf{-0.567}$ & $-0.576$ & \checkmark\ 비국소 \\
$L(s,\chi_4)$  & 120 & $-0.016$ & $\mathbf{-0.608}$ & $-0.629$ & \checkmark\ 비국소 \\
$L(s,\chi_5)$  & 127 & $-0.013$ & $\mathbf{-0.607}$ & $-0.613$ & \checkmark\ 비국소 \\
\midrule
GUE ($500\times 500$, 20회) & $20\times 399$
  & $-0.037$ & $-0.008\pm 0.034$ & $-0.020$ & $\times$\ 부재 \\
Poisson (234점, 20회) & $20\times 232$
  & $-0.290$ & $-0.090\pm 0.045$ & $-0.098$ & $\times$\ 기하학적 \\
\bottomrule
\end{tabular}
\end{table}

\begin{observation}[GUE/Poisson 대조]
\label{obs:gue_poisson}
비국소 신호가 반발 점 과정의 일반적 특성인지 검증하기 위해,
GUE 고유값($500{\times}500$ Hermitian 행렬, bulk 80\%, 슬라이딩 윈도우
unfolding, 20회)과 Poisson 점(234개 균등분포 점, 20회)에 대해
$R(m)^2 = (\sum_{\text{non-NN}} 1/(m - \lambda_k))^2$를 계산하였다.

\smallskip\noindent
GUE: $\rho(b) = -0.008 \pm 0.034$ ($t$검정 $p = 0.29$) ---
GUE의 spectral rigidity에도 불구하고 통계적으로 0과 구분 불가.
Poisson: $\rho(b) = -0.090 \pm 0.045$ ($p = 4.5\times 10^{-8}$) ---
통계적으로 유의하나 $|{\rho}|$가 $L$-함수 값의 ${\sim}1/7$ 수준;
이 잔여는 순수 기하학적 효과에서 기인한다: 지배 비-NN 항
$1/(m - \lambda_{i+1}) \approx -2/(\mathrm{gap} + \mathrm{gap}_{\mathrm{next}})$가
$\mathrm{gap}_{\mathrm{next}}$에 약하게 결합한다.

\smallskip\noindent
결론: 중간점 곡률 비국소성은 테스트한 4개 $L$-함수
($\zeta$, $\chi_3$, $\chi_4$, $\chi_5$)의 산술적 영점 배치에 고유하다.
GUE 레벨 반발이나 Poisson 독립성으로는 재현되지 않는다.
메커니즘은 미확인 상태이며, 열린 문제로 기록한다.
\textcolor{ForestGreen}{[확인 2026-04-14, 20회 앙상블]}
\end{observation}

\begin{observation}[비국소 도달 범위 감쇠 및 전방--후방 부호 반전]
\label{obs:nonlocal_reach}
중간점 곡률이 ``얼마나 멀리 보는지'' 정량화하기 위해,
관측~\ref{obs:midpoint_nonlocal}을 확장하여
$\rho(\kappa_{\mathrm{mid}},\, \mathrm{gap}_{n+k})$ ($k = 0, 1, \dots, 10$, 전방) 및
$\rho(\kappa_{\mathrm{mid}},\, \mathrm{gap}_{n-k})$ ($k = 1, \dots, 10$, 후방)을
4개 $L$-함수 전체에 대해 측정하였다.

\smallskip\noindent
\textbf{전방 감쇠 ($\zeta$, $n=234$):}
\begin{center}
\small
\begin{tabular}{@{} r r@{\;}c@{\;}l r r @{}}
\toprule
$k$ & \multicolumn{3}{c}{$\rho_{\mathrm{fwd}}$}
  & $p$-값 & 편상관 $\rho$ \\
\midrule
0 & $-$&$0.031$& & $6.3\times 10^{-1}$ & --- \\
\textbf{1} & $\mathbf{-}$&$\mathbf{0.658}$& & $\mathbf{3.7\times 10^{-30}}$ & $\mathbf{-0.668}$ \\
2 & $+$&$0.068$& & $3.0\times 10^{-1}$ & $+0.071$ \\
3 & $+$&$0.114$& & $8.6\times 10^{-2}$ & $+0.121$ \\
4 & $+$&$0.057$& & $3.9\times 10^{-1}$ & $+0.065$ \\
5 & $+$&$0.003$& & $9.7\times 10^{-1}$ & $+0.009$ \\
\bottomrule
\end{tabular}
\end{center}
신호는 $k{=}1$ ($|\rho|=0.658$)에서 $k{=}2$ ($|\rho|=0.068$)로 한 자릿수 감소하여,
비국소 결합이 가장 가까운 비인접 간격에 집중됨을 확인한다.

\smallskip\noindent
\textbf{$k{=}1$에서의 전방--후방 부호 반전 (4개 $L$-함수 전부):}
\begin{center}
\small
\begin{tabular}{@{} l c c c @{}}
\toprule
$L$-함수 & $\rho_{\mathrm{fwd}}(k{=}1)$ & $\rho_{\mathrm{bwd}}(k{=}1)$
  & 부호 반전? \\
\midrule
$\zeta(s)$     & $-0.658$ & $+0.576$ & \checkmark \\
$L(s,\chi_3)$  & $-0.555$ & $+0.504$ & \checkmark \\
$L(s,\chi_4)$  & $-0.608$ & $+0.588$ & \checkmark \\
$L(s,\chi_5)$  & $-0.606$ & $+0.602$ & \checkmark \\
\bottomrule
\end{tabular}
\end{center}
\noindent
모든 테스트된 $L$-함수에서 전방 $k{=}1$ 상관은 \emph{음}이고
후방 상관은 \emph{양}이며, 크기는 유사하다
($||\rho_{\mathrm{fwd}}| - |\rho_{\mathrm{bwd}}|| < 0.08$).
이 부호 반전은 Hadamard 구조의 직접적 귀결이다:
중간점 $m = (\gamma_n + \gamma_{n+1})/2$에서 오른쪽 지배 비-NN 기여($\gamma_{n+2}$)와
왼쪽 기여($\gamma_{n-1}$)는 $\Lambda'/\Lambda$에 반대 부호로 진입한다.

\smallskip\noindent
\textbf{감쇠 모델.}
$|\rho_{\mathrm{fwd}}|$ 수열 ($k=1,\dots,10$)에 대해
멱법칙 피팅 $|\rho| \sim A\,k^{-\alpha}$이 일관되게 지수 피팅보다 우수하다
($R^2$: $\zeta$\,0.71, $\chi_4$\,0.76, $\chi_5$\,0.76;
$\chi_3$은 3시드 예비 결과 $R^2 = 0.38$).
급격한 감쇠로 단일 특성 길이 $\xi$ 정의는 어렵고,
비국소 신호는 $k{=}1$에 거의 전적으로 집중된다.
\textcolor{ForestGreen}{[확인 2026-04-14, 4개 $L$-함수, 3시드 예비]}
\end{observation}


\subsection{Chern 수 해석: 수치 검증}
\label{sec:chern_verification}

\cref{sec:berry}의 이론적 프레임워크는 영점 계수 함수 $N(T)$가 위상적 해석을 가짐을 예측한다:
\begin{equation}
  N(T) \;=\; c_1\bigl(\text{$\xif$-다발 위 직사각형}\bigr)
        \;=\; \frac{1}{2\pi}\oint_{\mathcal{C}}
              d\!\operatorname{arg}\Lambda(s),
  \label{eq:chern_counting}
\end{equation}
여기서 $\mathcal{C}$는 직사각형 $[\sigma_{\min}, \sigma_{\max}]
\times [t_0, T]$이고 $\Lambda$는 완비된 $L$-함수이다.
이를 \emph{위상 추적법}으로 수치 검증한다: 윤곽을 따라
$\operatorname{Im}\!\log\Lambda(s)$를 계산하고 $2\pi$-되감기로
위상차를 누적하며, $t > 300$에서 언더플로우를 일으키는
$\xif(s)$ 직접 계산을 회피한다.

\begin{observation}[$\zeta(s)$에 대한 Chern 수 검증]
\label{obs:chern_zeta}
$\sigma \in [0.3, 0.7]$, dps${}=100$, 64분할 사다리꼴 적분으로
세 수준의 검증을 수행하였다:

\smallskip\noindent
\textbf{(A) 개별 홀로노미 (첫 20개 영점).}
각 영점 $\gamma_n$ 주위에 반지름 $r = 0.3$의 원을 추적.
20개 잔류값 모두 $\operatorname{Re}(\oint/2\pi i) = 0.9984$
(64점 적분의 상수 체계적 오프셋; 128분할에서 $0.9996$),
단순 영점 잔류값${}= 1$ 확인.

\smallskip\noindent
\textbf{(B) 직사각형 위상 추적 ($n = 10$--$200$).}
\begin{center}
\small
\begin{tabular}{@{} r r r r r @{}}
\toprule
$n$ & $t_2$ & $N_{\mathrm{actual}}$ & $N_{\mathrm{ctr}}$ & $|N_{\mathrm{ctr}} - N_{\mathrm{actual}}|$ \\
\midrule
  10 &  51.37 &  10 &   9.9882 & 0.0118 \\
  50 & 144.56 &  50 &  49.9882 & 0.0118 \\
 100 & 237.15 & 100 &  99.9882 & 0.0118 \\
 200 & 397.15 & 200 & 199.9882 & 0.0118 \\
\bottomrule
\end{tabular}
\end{center}
\noindent
6개 테스트($n \in \{10, 20, 50, 100, 150, 200\}$) 모두
상수 오프셋 $|N_{\mathrm{ctr}} - N_{\mathrm{actual}}| = 0.0118$로 통과.
수평변 64분할의 $O(1/N^2)$ 사다리꼴 오차에 기인
(128분할에서 $0.003$으로 확인).

\smallskip\noindent
\textbf{(C) 연속 계단함수 ($t \in [14, 430]$, 61개 점).}
누적 위상 함수 $C(t)$는 $N(t)$를 완벽한 계단함수로 추적하며,
각 영점에서 $+1$ 점프, 상수 오프셋 $C(t) = N(t) - 0.022$.
최대 편차: $|C - N| = 0.0223$.
\textcolor{ForestGreen}{[확인 2026-04-14, dps${}=100$, $t \le 430$ ($\approx 219$개 영점)]}
\end{observation}

\begin{observation}[Chern 수 공식의 디리클레 보편성]
\label{obs:chern_dirichlet}
위상 추적 검증을 홀수($\chi(-1) = -1$, $a = 1$) 및 짝수($\chi(-1) = +1$,
$a = 0$) 원시 디리클레 지표 모두로 확장하여
공식~\eqref{eq:chern_counting}의 보편성을 확인한다.

홀수 지표의 경우 $L(\tfrac{1}{2} + it, \chi)$가 복소수이므로
$|L|$ 최소화 + \texttt{findroot}로 영점을 탐색한다.
짝수 실수 지표의 경우 완전화 함수 $\Lambda(\tfrac{1}{2} + it, \chi)$는
실수이나, $L$ 자체는 $\Gamma(s/2)$ 인자의 복소 위상을 포함하므로
$|L|$ 최소화 방식이 보편적으로 올바른 접근이다.

\begin{center}
\small
\begin{tabular}{@{} l l r r r r r @{}}
\toprule
지표 & 패리티 & $n$ & $t_2$ & $N_{\mathrm{actual}}$ & $N_{\mathrm{ctr}}$
  & $|N_{\mathrm{ctr}} - N_{\mathrm{actual}}|$ \\
\midrule
$\chi_3$ (mod 3) & 홀 &  5 & 22.26 &  5 &  4.9988 & 0.0012 \\
                  &    & 10 & 34.75 & 10 &  9.9988 & 0.0012 \\
                  &    & 15 & 45.20 & 15 & 14.9988 & 0.0012 \\
                  &    & 20 & 54.92 & 20 & 19.9988 & 0.0012 \\
\midrule
$\chi_4$ (mod 4) & 홀 &  5 & 19.87 &  5 &  4.9985 & 0.0015 \\
                  &    & 10 & 31.12 & 10 &  9.9983 & 0.0017 \\
                  &    & 15 & 41.06 & 15 & 14.9983 & 0.0017 \\
                  &    & 20 & 50.70 & 20 & 19.9983 & 0.0017 \\
\midrule
$\bigl(\frac{\cdot}{5}\bigr)$ (mod 5) & 짝 &  5 & 18.55 &  5 &  4.9986 & 0.0014 \\
                  &    & 10 & 29.08 & 10 &  9.9983 & 0.0017 \\
                  &    & 15 & 38.84 & 15 & 14.9983 & 0.0017 \\
                  &    & 20 & 47.42 & 20 & 19.9983 & 0.0017 \\
\midrule
$\chi_8$ (mod 8) & 짝 &  5 & 16.11 &  5 &  4.9981 & 0.0019 \\
                  &    & 10 & 25.58 & 10 &  9.9974 & 0.0026 \\
\bottomrule
\end{tabular}
\end{center}
\noindent
14/14 테스트(홀수 8 + 짝수 6) 모두 $|N_{\mathrm{ctr}} - N_{\mathrm{actual}}| < 0.003$으로 통과하여,
Chern 수 해석~\eqref{eq:chern_counting}이 $\zeta(s)$를 넘어
양쪽 패리티의 디리클레 $L$-함수로 확장됨을 확인한다.
오프셋(0.0012--0.0026 vs.\ $\zeta$의 0.0118)은
낮은 높이 범위($t \le 55$ vs.\ $t \le 430$)에 기인한다.

\smallskip\noindent
\textit{방법론 참고.}
짝수 지표 $(\frac{\cdot}{5})$에서 두 영점 탐색 방식---$\operatorname{Re}L$
부호변화 vs.\ $|L|$ 최소화---을 교차 검증한 결과,
전자가 27개 중 2개의 접선 영점($\operatorname{Re}L$이 부호를 바꾸지 않는 영점)을
놓치는 것을 확인하여, $|L|$ 최소화가 엄밀하게 더 신뢰할 수 있음을 확인했다.
\textcolor{ForestGreen}{[확인 2026-04-14; 홀: $\chi_3$~mod~3 + $\chi_4$~mod~4;
짝: $(\frac{\cdot}{5})$~mod~5 + $\chi_8$~mod~8]}
\end{observation}


\begin{observation}[$\xif$-다발의 Gauss--Bonnet 면적분]
\label{obs:gauss_bonnet}
윤곽 적분 Chern 수(관찰~\ref{obs:chern_zeta}--\ref{obs:chern_dirichlet})의
면적분 등가물로서, 위상장
$\theta(\sigma,t) = \operatorname{Im}\!\log\xif(\sigma+it)$ 위에서
\textbf{격자 플라켓 곡률}---격자 게이지 이론의 표준 방법---을 계산한다.
각 격자 셀의 감긴 위상 순환이 국소 와동 전하를 주며,
전체 합은 이산 Stokes 정리에 의해 $2\pi N$을 준다:
\begin{equation}
\label{eq:gauss_bonnet_vortex}
\iint_R F_2\,d\sigma\,dt
  \;=\; \sum_{\text{플라켓}} \Delta\theta_{\square}
  \;=\; 2\pi\,N(R),
\end{equation}
여기서 $N(R)$은 영역 $R$ 내부의 $\xif$ 영점 수이다.

\smallskip\noindent
\textbf{결과} ($\mathrm{dps}=100$, $N_\sigma=64$,
$\sigma\in[0.3,0.7]$):

\begin{center}
\begin{tabular}{llrrrc}
\toprule
테스트 & 영역 & $\iint F_2/(2\pi)$ & 기대 $N$ & $|\text{오차}|$ & \\
\midrule
\S A: 단일 영점 ($\gamma_1$) & $[0.3,0.7]\!\times\![13,15.5]$ & 1.0000 & 1 & 0.000000 & \checkmark \\
\S B: $k\!=\!1$ & $[0.3,0.7]\!\times\![0.5,17.58]$ & 1.0000 & 1 & 0.000000 & \checkmark \\
\S B: $k\!=\!5$ & $[0.3,0.7]\!\times\![0.5,35.26]$ & 5.0000 & 5 & 0.000000 & \checkmark \\
\S B: $k\!=\!10$ & $[0.3,0.7]\!\times\![0.5,51.37]$ & 10.0000 & 10 & 0.000000 & \checkmark \\
\S B: $k\!=\!20$ & $[0.3,0.7]\!\times\![0.5,78.24]$ & 20.0000 & 20 & 0.000000 & \checkmark \\
\S C: 빈 영역 & $[0.3,0.7]\!\times\![\gamma_1\!+\!0.5,\gamma_2\!-\!0.5]$ & 0.0000 & 0 & 0.000000 & \checkmark \\
\S D-1: $\chi_5$ 단일 & $(\frac{\cdot}{5})$ 첫 영점 & 1.0000 & 1 & 0.000000 & \checkmark \\
\S D-2: $\chi_5$ 5영점 & $(\frac{\cdot}{5})$ 5개 영점 & 5.0000 & 5 & 0.000000 & \checkmark \\
\bottomrule
\end{tabular}
\end{center}

\noindent
8/8 테스트 모두 $2\pi$의 \emph{정확한} 정수배를 준다 (표시 정밀도 내).
이는 위상적 불변량의 본질적 성질이다:
플라켓 자속은 격자 해상도에 무관하게 양자화된다
($16\!\times\!16$에서 독립 검증).
빈 영역 테스트(\S C)는 $F_2$가 영점에 국소화됨을 확인한다---
각 영점이 정확히 $2\pi$의 자속을 기여하며,
인접 영역으로의 누출이 없다.

\smallskip\noindent
윤곽 적분 결과(관찰~\ref{obs:chern_zeta}--\ref{obs:chern_dirichlet})와
결합하면, Stokes 쌍대성이 확립된다:
$\oint_{\partial R}\!\operatorname{Im}(L\,ds)
= \iint_R F_2\,d\sigma\,dt = 2\pi N$.
이는 $\xif$-다발의 곡률 구조를 경계와 벌크 양쪽 관점에서 확인한다.
면적분은 플라켓 감김이 정확한 $2\pi$ 양자화를 강제하므로
윤곽 방법($|\text{오차}|=0.012$)보다 높은 수치 정밀도($|\text{오차}|=0$)를 달성한다.
\textcolor{ForestGreen}{[확인 2026-04-14; $\zeta$ + $(\frac{\cdot}{5})$~mod~5]}
\end{observation}


\begin{observation}[위상적 전하의 $\sigma$-국소화]
\label{obs:sigma_localisation}
Gauss--Bonnet 적분(관찰~\ref{obs:gauss_bonnet})은 직사각형 영역 $R$ 내의 영점 수를 센다.
이제 곡률 $F_2$가 $\sigma$-방향 어디에 집중되는지 묻는다.
동일한 $t$-범위 $[13,\,(\gamma_{20}+\gamma_{21})/2]$ (처음 20개 $\zeta$-영점 포함)를
유지하면서 $\sigma$-strip을 좁히는 일련의 \textbf{strip 축소 테스트}를 수행한다.

\smallskip\noindent
\textbf{결과} ($\mathrm{dps}=100$, $N_\sigma=64$--$128$):

\begin{center}
\begin{tabular}{llrrrc}
\toprule
테스트 & $\sigma$-strip & $\iint F_2/(2\pi)$ & 기대값 & $|\text{오차}|$ & \\
\midrule
\S A1: 기준 & $[0.30,\,0.70]$ & 20.0000 & 20 & 0.000000 & \checkmark \\
\S A2: 좁은 & $[0.49,\,0.51]$ & 20.0000 & 20 & 0.000000 & \checkmark \\
\S A3: 극좁은 & $[0.499,\,0.501]$ & 20.0000 & 20 & 0.000000 & \checkmark \\
\S A4: 오른쪽 off-critical & $[0.60,\,0.90]$ & 0.0000 & 0 & 0.000000 & \checkmark \\
\S A5: 왼쪽 off-critical & $[0.10,\,0.40]$ & 0.0000 & 0 & 0.000000 & \checkmark \\
\S A6: 왼쪽 반-strip & $[0.30,\,0.499]$ & 0.0000 & 0 & 0.000000 & \checkmark \\
\S A7: 오른쪽 반-strip & $[0.501,\,0.70]$ & 0.0000 & 0 & 0.000000 & \checkmark \\
\bottomrule
\end{tabular}
\end{center}

\noindent
\textbf{$\sigma$-프로파일.}\;
$\sigma\in[0.3,0.7]$의 200개 bin에서 곡률 밀도
$\rho(\sigma)=\sum_j \Delta\theta_{\square}(\sigma,t_j)$를 측정:
피크값 $\rho(0.5000)=20\times 2\pi$,
$\mathrm{FWHM}<\delta\sigma=0.004$ (해상도 한계);
나머지 199개 bin의 잔류값 $<1.4\times10^{-15}$ (머신 엡실론).

\smallskip\noindent
\textbf{디리클레 확장:}\;
$\chi_5=(\frac{\cdot}{5})$에 대해 좁은 strip $[0.49,0.51]$에서
$N=5.000000$ (기대 5); off-critical $[0.6,0.9]$에서
$N=0.000000$---$\sigma=1/2$에 동일한 $\delta$-집중 확인.

\smallskip\noindent
이 결과는 처음 20개 $\zeta$-영점과 5개 $L(\chi_5)$-영점에 대해
RH와 일관적인 기하학적 증거를 제공한다:
플라켓 곡률 $F_2$는 $\sigma=1/2$에 $\delta$-함수로 집중되며,
외부에는 잔류 자속이 없다.
$\sigma=0.5$를 0.001만큼만 배제해도(\S A6--A7 테스트) 즉시 $N=0$이 된다.
이는 이 영점들이 임계선 위에 있다는 알려진 사실의 기하학적 재표현이며---
이 실험의 기여는 집중의 \emph{정량적 선예도}이다.
\textcolor{ForestGreen}{[확인 2026-04-14; $\zeta$ 7/7 + 프로파일 + $\chi_5$ 2/2;
$16\!\times\!16$, $\mathrm{dps}=50$에서 독립 재현]}
\end{observation}


\begin{observation}[三重 위상적 일관성 표]
\label{obs:triple_consistency}
세 독립 방법---직접 영점 수 세기~(A),
Chern 윤곽 적분~(B, 관찰~\ref{obs:chern_zeta}),
Gauss--Bonnet 면적분~(C, 관찰~\ref{obs:gauss_bonnet})---을
결합하여, 영점 수 $N=10,\,20,\,50,\,100$ 범위에서
$\xif$-다발 프레임워크의 자기 정합성을 검증한다.
각 $N$에서 상한 $t_2 = (\gamma_N+\gamma_{N+1})/2$로 설정하여
경로가 영점을 지나지 않도록 한다.

\smallskip\noindent
\textbf{결과} ($\mathrm{dps}=100$, $\sigma\in[0.3,0.7]$,
$N_\sigma=64$, Chern 수평변 64분할):

\begin{center}
\begin{tabular}{rr rrr rr}
\toprule
$N$ & $t_2$ & A (직접) & B (Chern) & C (GB)
    & $|B{-}A|$ & $|C{-}A|$ \\
\midrule
10  & 51.37 & 10  & 9.9882  & 10.0000 & 0.0118 & 0.0000 \\
20  & 78.24 & 20  & 19.9881 & 20.0000 & 0.0119 & 0.0000 \\
50  & 144.56 & 50 & 49.9882 & 50.0000 & 0.0118 & 0.0000 \\
100 & 237.15 & 100 & 99.9882 & 100.0000 & 0.0118 & 0.0000 \\
\bottomrule
\end{tabular}
\end{center}

\noindent
\textbf{정밀도 계층.}\;
Gauss--Bonnet 방법은 모든 $N$에서 정확한 정수 양자화($|C{-}A|=0$)를 달성하고,
Chern 윤곽 적분은 $N$에 무관한 일정한 체계적 오프셋 $|B{-}A|\approx 0.012$를 보인다.
이 오프셋은 위상 추적 윤곽의 수평변 이산화에서 기인하며,
영점당 오차 누적이 \emph{아니다}---$N$-무관성 자체가 정합성 검증이다.

\smallskip\noindent
\textbf{고높이 $\sigma$-국소화 ($N\!=\!100$).}\;
관찰~\ref{obs:sigma_localisation}을 처음 100개 $\zeta$-영점($t_2=237.15$)으로 확장:
좁은 strip $\sigma\in[0.49,0.51]$에서 $N=100.000000$;
off-critical $\sigma\in[0.6,0.9]$에서 $N=0.000000$.
$\sigma=1/2$의 $\delta$-집중은 관찰~\ref{obs:sigma_localisation}에서
테스트한 높이의 5배까지 유지된다.

\smallskip\noindent
$N$의 4개 값에 걸친 三重 일치는 미분기하학적 프레임워크의 내적 정합성을 확립한다:
곡률(국소), Chern 수(반국소), Gauss--Bonnet(전역)이 모두 동일한 위상적 내용을 부호화한다.
\textcolor{ForestGreen}{[확인 2026-04-14; 4/4 三重 + 2/2 $\sigma$-국소화;
$16\!\times\!128$, $\mathrm{dps}=50$에서 독립 재현]}
\end{observation}


\begin{observation}[합성 off-critical 음성 대조군 (특이성)]
\label{obs:synthetic_control}
플라켓 방법이 $\sigma=1/2$를 \emph{내재적으로} 선호하는지 배제하기 위해,
영점이 알려진 위치에 놓인 세 합성 다항식 함수에 적용한다:
\begin{itemize}
\item $f_1(s)=\prod_{k}(s-z_k)$: 모든 영점이 $\sigma=1/2$ 위
      (``RH-유사'': $z_k=\tfrac12+14i,\;\tfrac12+21i,\;\tfrac12+25i$);
\item $f_2(s)$: 혼합 영점
      ($\sigma=0.3,\;0.7,\;0.5$);
\item $f_3(s)$: $\sigma=1/2$ 위 영점이 \emph{없음}
      ($\sigma=0.3,\;0.7,\;0.6,\;0.4$).
\end{itemize}

\smallskip\noindent
\textbf{결과} ($\mathrm{dps}=30$, $N_\sigma=64$, $N_t=200$;
3개 함수에 걸쳐 11종 테스트):

\begin{center}
\begin{tabular}{llcccl}
\toprule
테스트 & 함수 & $\sigma$-범위 & $N_{\mathrm{plaq}}$ & 기대값 & 통과 \\
\midrule
\S A-1 & $f_1$ & $[0.3,0.7]$    & 3.0000 & 3 & \checkmark \\
\S A-2 & $f_1$ & 프로파일        & 1피크@0.500 & 1피크 & \checkmark \\
\S A-3 & $f_1$ & $[0.6,0.9]$    & 0.0000 & 0 & \checkmark \\
\S B-1 & $f_2$ & $[0.2,0.8]$    & 3.0000 & 3 & \checkmark \\
\S B-2 & $f_2$ & 프로파일        & 3피크 & 3피크 & \checkmark \\
\S B-3 & $f_2$ & $[0.49,0.51]$  & 1.0000 & 1 & \checkmark \\
\S B-4 & $f_2$ & $[0.6,0.8]$    & 1.0000 & 1 & \checkmark \\
\S B-5 & $f_2$ & $[0.2,0.4]$    & 1.0000 & 1 & \checkmark \\
\S C-1 & $f_3$ & $[0.2,0.8]$    & 4.0000 & 4 & \checkmark \\
\S C-2 & $f_3$ & 프로파일        & 4피크, $\rho(\tfrac12)=0$ & 4피크 & \checkmark \\
$\star$\,\S C-3 & $f_3$ & $[0.49,0.51]$ & \textbf{0.0000} & \textbf{0} & \checkmark \\
\bottomrule
\end{tabular}
\end{center}

\noindent
\textbf{핵심 결과 (\S C-3).}\;
임계선 위에 영점이 \emph{없는} $f_3$에서,
$\sigma\in[0.49,0.51]$ strip의 플라켓 자속은 정확히 0이다.
프레임워크가 인위적으로 $\sigma=1/2$에 전하를 집중시킨다면,
영점 위치와 무관하게 이 strip에서 비영 자속이 관측될 것이다.
이 영 결과는 \textbf{특이성}을 확립한다:
$\xif$에서 관측된 $\delta$-집중(관찰~\ref{obs:sigma_localisation}--\ref{obs:triple_consistency})은
$\xif$-영점 자체의 성질이지, 방법론적 인공물이 아니다.

\smallskip\noindent
\textbf{\S B 개별 $\sigma$-분리.}\;
$f_2$의 세 영점 $\sigma=0.3,\,0.5,\,0.7$이 각각
자신의 좁은 strip에만 $N=1$로 기여(\S B-3/B-4/B-5)하여,
플라켓이 각 영점의 실수부를 해상함을 확인한다.

\smallskip\noindent
11종 테스트 모두 정확한 정수 양자화(오차 $=0$)를 보이는데,
다항식이므로 당연하다; 의의는 정밀도가 아니라
\emph{$\sigma=1/2$ 편향의 부재}에 있다.
\textcolor{ForestGreen}{[확인 2026-04-14; 11/11;
$32\!\times\!100$ 및 $48\!\times\!100$ 격자에서
(\S C-3, \S B-3, \S C-1) 독립 재현]}
\end{observation}

\begin{observation}[수 분산: Berry 산술 보정 레짐에서의 sub-GUE]
\label{obs:number_variance}
플라켓 곡률장이 개별 홀로노미뿐만 아니라 리만 영점의
\emph{집합적} 통계를 올바르게 인코딩하는지 검증하기 위해,
수 분산 $\Sigma^2(L)$을 두 가지 독립적 방법으로 계산한다:
(i)~전개(unfolded) 창 크기 $L$에서의 직접 영점 계수,
(ii)~슬라이딩 창에 걸친 플라켓 총 플럭스 $N_{\mathrm{plaq}}$.

\smallskip\noindent
\textbf{설정.} 격자: $N_\sigma=32$, $N_t=2000$,
$\sigma\in[0.3,0.7]$, $t\in[50,600]$, dps${}=80$.
전개 창 $L\in\{5,10,20,50\}$, 슬라이딩 step $=1$ 영점,
경계 $t\in[70,580]$ 제외.
$[45,605]$ 구간에서 337개 영점; 평균 전개 간격 $=1.0008$
(이론$=1.000$, 오차$<0.1\%$).
GUE 예측: $\Sigma^2_{\mathrm{GUE}}(L)\approx
\tfrac{2}{\pi^2}\!\ln L+0.44$ (Berry--Keating 산술 보정 포함~\cite{berry1988}).

\smallskip\noindent
\textbf{결과.}

\begin{center}
\small
\begin{tabular}{@{} r c c c c c c @{}}
\toprule
$L$ & $N_{\mathrm{win}}$ & $n_{\mathrm{eff}}$ &
  $\Sigma^2_{\mathrm{dir}}$ & $\Sigma^2_{\mathrm{plaq}}$ &
  $\Sigma^2_{\mathrm{GUE}}$ & $\Sigma^2/\Sigma^2_{\mathrm{GUE}}$ \\
\midrule
 5 & 304 & 60 & 0.3960 & 0.4257 & 0.7682 & 0.515 \\
10 & 299 & 29 & 0.5134 & 0.5604 & 0.9086 & 0.565 \\
20 & 289 & 14 & 0.4826 & 0.4792 & 1.0491 & 0.460 \\
50 & 259 &  5 & 0.4494 & 0.4841 & 1.2348 & 0.364 \\
\bottomrule
\end{tabular}
\end{center}

\smallskip\noindent
\textbf{세 가지 발견.}
\begin{enumerate}
  \item \textbf{방법론 검증 (플라켓 $\approx$ 직접 계수).}
        $\Sigma^2_{\mathrm{plaq}}$와 $\Sigma^2_{\mathrm{dir}}$이
        모든 창 크기에서 $10\%$ 이내로 일치한다.
        플라켓 곡률장이 직접 영점 계수와 동일한 집합적 통계 구조를
        포착함을 확인 — 위상적 대응을 개별 와동(홀로노미,
        관측 \ref{obs:chern_zeta}--\ref{obs:chern_dirichlet})에서
        앙상블 분산으로 확장한다.

  \item \textbf{Sub-GUE 레짐: Berry 산술 보정.}
        $\Sigma^2/\Sigma^2_{\mathrm{GUE}}=0.36$--$0.57$로
        일관되게 1 미만이다.
        이는 이상 현상이 아니라, Berry~\cite{berry1988}가 예측하고
        Odlyzko~\cite{odlyzko1987}가 수치적으로 확인한
        \emph{산술 보정}의 기대 결과이다:
        높이 $t\sim 300$에서 GUE 로그 성장은 소수합 산술 진동에 의해
        억제되며, 완전한 GUE 수렴은 $t\gg 10^6$에서만 나타난다.
        $\Sigma^2(L)$의 log--log 기울기는 $0.036$
        (GUE 예측 $\sim\!1$; 완전 평탄 $=0$)으로,
        Berry의 ``산술 포화(arithmetic saturation)'' 레짐을
        정확히 재현한다.

  \item \textbf{Poisson 완전 배제.}
        Poisson 예측은 $\Sigma^2_{\mathrm{Poi}}=L$.
        관측 분산은 Poisson 수준보다 $63$--$156\,\sigma$ 낮아,
        플라켓 플럭스가 영점 배열의 수준 반발(level repulsion)
        구조를 올바르게 반영함을 확인한다.
\end{enumerate}

\noindent
\textbf{해석 주의.}\;
Sub-GUE 결과를 GUE 보편성의 실패로 해석해서는 안 된다.
이는 알려진 유한 높이 효과~\cite{berry1988}이다:
``sub-GUE''와 ``GUE''는 높이 레짐이 다를 뿐 동일한 보편성 부류에 속한다.
여기서의 핵심 의의는 \emph{방법론 검증}에 있다:
플라켓이 개별 권선수를 넘어 집합적 통계를 올바르게 포착한다.
\textcolor{ForestGreen}{[확인 2026-04-14; 창 크기 4종;
플라켓을 직접 영점 계수와 독립적으로 교차 검증]}
\end{observation}


% ----------------------------------------------------------------------------
\begin{observation}[Epstein 제타 확장]
\label{obs:epstein_extension}
다발 프레임워크가 리만 $\zeta$ 및 디리클레 $L$-함수를 넘어 확장되는지 검증하기 위해
이차형식 $Q(m,n) = m^2 + \lambda n^2$의 Epstein 제타 함수
$Z_Q(s) = \sum_{(m,n)\neq(0,0)} Q(m,n)^{-s}$를 검사한다.

\smallskip\noindent
\textbf{설정.} 두 격자 매개변수: $\lambda=1$ (대조군;
$Z_Q = 4\,\zeta(s)\,L(s,\chi_{-4})$, RH 만족) 및 $\lambda=5$
(실험군; 류수 $h(-20)=2$, 개별 Epstein 제타에서 RH 위반 가능).
$\theta$-Mellin 적분으로 해석적 접속; $Z_{Q,\lambda=1}$과
$4\,\zeta(s)\,L(s,\chi_{-4})$의 상대오차 $\sim 10^{-70}$으로 검증.
다발 진단: $\delta=0.03$ 오프셋에서 $\kappa$, 폐곡선 적분 모노드로미
(반경 $=0.5$, 64단계), $\sigma$-국소화 FWHM.
임계선 스캔: $t\in[2,35]$; off-critical 격자:
$\sigma\in[0.25,0.75]$, $t\in[2,35]$, $12\times 100$ 격자.

\smallskip\noindent
\textbf{결과.}

\begin{center}
\small
\begin{tabular}{@{} l c c c c @{}}
\toprule
& $\kappa$ & mono$/\pi$ & FWHM & off-crit.\ 영점 \\
\midrule
$\lambda=1$ ($N=8$) & $1007.9 \pm 72.6$ & $2.00 \pm 0.00$ & $0.044$ & 0 \\
$\lambda=5$ ($N=8$) & $1006.2 \pm 47.8$ & $2.00 \pm 0.00$ & $0.044$ & 0 \\
\bottomrule
\end{tabular}
\end{center}

\smallskip\noindent
\textbf{발견.}
\begin{enumerate}
  \item \textbf{On-critical 서명 보편성.}
        세 다발 지표($\kappa$, 모노드로미, FWHM) 모두 $\lambda=1$과 $\lambda=5$에서
        구별 불가(편차 ${}<1\%$). 프레임워크의 진단 서명이
        이차형식에 무관함을 시사한다.

  \item \textbf{Off-critical 영점 미발견.}
        탐색 범위 내 $\sigma\neq\tfrac{1}{2}$ 영점 없음.
        이는 부재 증거가 아닌 \emph{범위 한계}: $\lambda=5$의
        off-critical 영점은 $t \gg 35$에 존재할 수 있다.
\end{enumerate}

\noindent
\textbf{해석.}\;
Epstein 제타 결과는 다발 프레임워크가 세 번째 $L$-함수 계열($\zeta$ 및
디리클레 $L$ 이후)로 확장 가능함을 확인하지만,
반증 가능성 테스트는 완성하지 \emph{못한다}.
off-critical 영점의 서명 차이를 측정하려면 off-critical 영점이
\emph{보장된} 함수가 필요하다 (예: Davenport--Heilbronn 함수; \S\ref{sec:open} 참조).
\textcolor{ForestGreen}{[확인 2026-04-14; $\lambda\in\{1,5\}$;
dps${}=80$; 78.7분 실행]}
\end{observation}


% ----------------------------------------------------------------------------
\begin{observation}[고 $t$ 견강성 및 $\kappa$ 교정]
\label{obs:high_t_robustness}
다발 프레임워크의 핵심 진단이 큰 허수 부분에서도 안정적으로
유지됨을 검증하기 위해, $t \approx 14$에서 $t \approx 10\,000$까지
6개 리만 영점($N=1$ \textasciitilde{} $N=10\,142$)에 걸쳐
$\kappa$, 모노드로미, $\sigma$-국소화를 측정한다.

\smallskip\noindent
\textbf{설정.}\;
모든 측정은 해석적 로그-도함수 공식을 사용한다:
\[
  \frac{\xi'(s)}{\xi(s)}
  \;=\;
  \frac{1}{s} + \frac{1}{s-1}
  - \frac{\log\pi}{2}
  + \frac{\psi(s/2)}{2}
  + \frac{\zeta'(s)}{\zeta(s)}.
\]
이는 큰 $t$에서 유한 차분 $\xi'$의 수치 불안정성을 회피한다.
$\kappa = |\xi'/\xi|^2$를 오프셋 $\delta=0.03$ ($\sigma=0.53$)에서 평가.
$\sigma$-프로파일은 \emph{고정된} 201점 격자
($\sigma\in[0.2,0.8]$, 간격 $0.003$)를 모든 $t$에서 사용하여,
예비 실행에서 발견된 해상도 아티팩트를 제거.
모노드로미: 윤곽 적분, 반경 $0.2$, 64단계.
정밀도: $t<2000$일 때 $\mathrm{dps}=60$, $t\ge 2000$일 때 $\mathrm{dps}=80$.

\smallskip\noindent
\textbf{파트~A: $\kappa$ 교정.}\;
먼저 동일한 해석적 공식으로 저 $t$ 기준 $\kappa$를 재교정한다
(이전에 \texttt{bundle\_utils}의 $h=10^{-20}$ 유한 차분으로 $\kappa\approx 1030$ 측정).

\begin{center}
\small
\begin{tabular}{@{} r r r r @{}}
\toprule
$N$ & $t$ & $\kappa$ (해석적) & $\kappa/(1/\delta^2)$ \\
\midrule
1   & 14.135  & 1111.58 & 1.0004 \\
2   & 21.022  & 1111.96 & 1.0008 \\
3   & 25.011  & 1111.78 & 1.0006 \\
4   & 30.425  & 1112.58 & 1.0013 \\
5   & 32.935  & 1111.92 & 1.0007 \\
\midrule
\multicolumn{2}{l}{평균 ($t<100$)} & $1111.96 \pm 0.33$ & --- \\
\bottomrule
\end{tabular}
\end{center}

해석적 공식의 저 $t$ $\kappa$는 $1112.0$으로, 고 $t$ 값 $1118.6$과
$0.6\%$ 이내로 일치한다. 이전 기준값 $\kappa\approx 1030$은
$h=10^{-20}$ 유한 차분 방식의 과잉 상쇄(over-cancellation)로 인한
$8\%$ 과소평가였다.
교정 기준값: $\kappa \approx 1112$--$1120$ (해석적).

\smallskip\noindent
\textbf{파트~B: $t$-스케일링 비교 (동일 201점 격자).}

\begin{center}
\small
\begin{tabular}{@{} l c c c @{}}
\toprule
지표 & $t<100$ (3 영점) & $t>1000$ (3 영점) & 차이 \\
\midrule
$\kappa$ (해석적) & $1112.0 \pm 0.3$ & $1118.6 \pm 1.9$ & $0.6\%$ \\
$\mathrm{peak}_\sigma$ & $0.4970$ & $0.4970$ & $0.0000$ \\
FWHM & $0.0060$ & $0.0060$ & $1.000\times$ \\
mono$/\pi$ & $2.0000$ & $2.0000$ & 정확 \\
\bottomrule
\end{tabular}
\end{center}

\smallskip\noindent
\textbf{발견.}
\begin{enumerate}
  \item \textbf{$\kappa$ $t$-안정성.}\;
        곡률 스칼라 $\kappa$는 $t$의 $700\times$ 범위에 걸쳐 $0.6\%$만
        변동한다. 선도 차수 값 $\kappa \approx 1/\delta^2 = 1111.1$은
        $\xi'/\xi$의 극 구조에서 해석적으로 예상되며;
        차선도 $\zeta'/\zeta$ 보정은 ${}<0.5\%$ 기여하고 $t$에 무관하다.

  \item \textbf{모노드로미 안정성.}\;
        6개 영점 모두 mono $= 2.0000\pi$. 반경 $0.2$에서의 단일 영점
        감싸기가 $t \approx 10\,000$에서도 확인됨
        (이 높이에서 평균 영점 간격 $\approx 0.68$).

  \item \textbf{$\sigma$-국소화 불변성.}\;
        동일 격자 해상도에서 $\mathrm{peak}_\sigma$와 FWHM은
        저 $t$와 고 $t$ 간 구별 불가. 예비 실행에서 가변 격자 크기로
        나타난 외견상 $t$-의존성은 전적으로 해상도 아티팩트였다.
\end{enumerate}

\noindent
\textbf{해석.}\;
이 결과는 새로운 수학적 발견이 아닌 \emph{견강성 확인}이다:
해석적 $\xi'/\xi$ 인프라와 세 가지 다발 진단($\kappa$, 모노드로미,
$\sigma$-FWHM)이 $t = 10^4$까지 수치적으로 안정함을 검증한다.
$\kappa$ 교정은 또한 유한 차분과 해석적 기준선 사이의 오랜
$8\%$ 불일치를 해결한다.
\textcolor{ForestGreen}{[확인 2026-04-15; 6 영점 ($N\in\{1,2,3,649,4520,10142\}$);
해석적 $\xi'/\xi$ 공식; 201점 격자; 105초 실행]}
\end{observation}


% ----------------------------------------------------------------------------
\begin{observation}[곡률 차선도 구조와 Hadamard 분해]
\label{obs:curvature_subleading}
곡률 스칼라 $\kappa = |\xi'/\xi|^2$를 선도항 $1/\delta^2$ 극 너머로
전개하여, \emph{차선도} 보정 $\Delta\kappa := \kappa - 1/\delta^2$를
분리하고 이것이 Hadamard 곱 분해를 통해 산술적 내용을 인코딩함을
보인다.

\smallskip\noindent
\textbf{설정.}\;
20개 리만 영점($N = 1$ \textasciitilde{} $N = 20\,000$,
$t \approx 14$ \textasciitilde{} $t \approx 18\,046$)을
오프셋 $\delta = 0.03$에서 탐침한다.
해석적 로그 도함수 $\xi'/\xi(s) = G(s) + \zeta'/\zeta(s)$는
감마/극 기여 $G(s) = 1/s + 1/(s{-}1) - \tfrac{1}{2}\log\pi + \tfrac{1}{2}\psi(s/2)$와
디리클레 급수 부분 $\zeta'/\zeta(s)$로 분리된다.

$R(s) = \xi'/\xi(s) - 1/\delta$ (최근접 영점 극 제거 후 나머지)로
쓰면:
\[
  \Delta\kappa
  \;=\;
  \frac{2\,\operatorname{Re} R}{\delta}
  + |R|^2.
\]

\smallskip\noindent
\textbf{파트~A: $\delta$-독립성.}\;
보정항 $\Delta\kappa$는 $\delta$에 본질적으로 무관하다:
멱법칙 적합 $\Delta\kappa \propto \delta^{\alpha}$에서
$\alpha \approx -0.002$로 영과 일치한다.
이는 $\kappa = 1/\delta^2 + f(t)$로 깔끔히 분해됨을 의미한다.

\smallskip\noindent
\textbf{파트~B: 로그 스케일링.}\;
차선도 보정은 대수적으로 성장한다:
\[
  \Delta\kappa
  \;\approx\;
  1.59 \cdot \log\!\bigl(t/(2\pi)\bigr) + c,
  \qquad
  R^2 = 0.860,\;\; p = 4.1 \times 10^{-9}.
\]
$R^2 = 0.86$은 ``확립'' 기준 $0.9$에 미달하지만,
높은 유의성의 $p$-값이 대수적 추세를 확인한다.

\smallskip\noindent
\textbf{파트~C: 기울기의 Hadamard 분해.}\;
기울기 $a = 1.59$는 두 큰 기여의 근사 상쇄에서 발생한다:
\begin{center}
\small
\begin{tabular}{@{} l r r @{}}
\toprule
성분 & $\log(t/2\pi)$ 대비 기울기 & $R^2$ \\
\midrule
$G(n)$ [감마 + 극]      & $+0.5000$ & $1.0000$ \\
$S_{\mathrm{all}}(n)$ [제타 영점] & $-0.4768$ & $0.9996$ \\
$\operatorname{Re} R = G + S_{\mathrm{all}}$  & $+0.0232$ & $0.854$ \\
$2\operatorname{Re} R/\delta$  & $+1.548$ & $0.854$ \\
$\operatorname{Im}(R)^2$      & $+0.042$ & $0.041$ \\
$\Delta\kappa$ (전체)        & $+1.595$ & $0.860$ \\
\bottomrule
\end{tabular}
\end{center}
감마 기여(+0.500)와 영점합 기여($-0.477$)가 $95.4\%$ 상쇄되고,
잔차 $0.023$이 $2/\delta = 66.7$배로 증폭되어 관측 기울기
$\approx 1.55$를 만든다.
Hadamard 곱 구조에 기반한 이론 추정 $a \approx 1.5$는
실측 $1.59 \pm 0.15$와 $1\sigma$ 이내 일치한다.

\smallskip\noindent
\textbf{파트~D: 이상치 분석 ($N{=}7000$).}\;
최대 잔차($\Delta\kappa = 14.08$)는 근접 영점 쌍에서 기인한다:
영점 6999와 7000의 간격은 $0.475$ (지역 평균 $0.89$)로,
이례적으로 큰 $\operatorname{Re} R = 0.204$를 생성한다.
이는 GUE 레벨 반발 통계와 일관적이다.

\smallskip\noindent
\textbf{주의.}\;
영점 간격과 $\Delta\kappa$의 raw 상관 $r = -0.755$는
공통 교란변수 $t$에 의한 허위 연관이다.
Unfolding 후 $r = +0.256$ ($p = 0.28$):
간격--$\Delta\kappa$ 결합은 통계적으로 유의하지 않다.

\smallskip\noindent
\textbf{해석.}\;
차선도 구조 $\Delta\kappa \approx 1.59\log(t/2\pi)$는
$\xi$-다발 곡률이 자명한 $1/\delta^2$ 극 너머의 산술적 내용을
포착한다는 예비적 수치 증거를 제공한다.
Hadamard 분해는 (오차 $0.00\%$로) 이 기울기가 감마 함수와
제타 영점 기여의 근사 상쇄에서 발생함을 검증한다.
\smallskip\noindent
\textbf{파트~E: NNS 의존 Lorentzian$^2$ 보정 (99 영점).}\;
99개 영점($N = 203$ \textasciitilde{} $20\,000$,
$t \approx 402$--$18\,046$)으로 확장하여,
최근접 이웃 간격(NNS)이 1변수 log-적합의 잔차
$\operatorname{res}_n := \Delta\kappa_n - (a\log(t_n/2\pi) + b)$를
설명하는지 검증한다.

1변수 OLS는 $R^2 = 0.015$ ($p = 0.22$): 이 규모에서 두 극단 이상치
($\Delta\kappa = 568, 259$; 모두 NNS${}_\mathrm{unf} < 0.15$)가
잔차를 지배한다.

\emph{NNS 상관.}\;
Pearson $r(\mathrm{res},\mathrm{NNS}_\mathrm{unf}) = -0.474$
($p = 7.4 \times 10^{-7}$); Spearman $\rho = -0.739$
($p = 2.6 \times 10^{-18}$). 이론적으로 동기 부여된
Lorentzian$^2$ 가설 $f_1 = \delta^2/(\delta^2 + g_{\min}^2)^2$
($g_{\min}$은 최근접 raw 간격)에서
Pearson $r(\mathrm{res}, f_1) = 0.953$ ($p \approx 10^{-52}$).

\emph{2변수 회귀.}\;
\begin{center}
\small
\begin{tabular}{@{} l c c @{}}
\toprule
NNS 함수 $f$ & $R^2_{\text{2-var}}$ & $F$-통계량 \\
\midrule
$f_1 = \delta^2/(\delta^2 + g_{\min}^2)^2$ (Lorentzian$^2$)
  & $\mathbf{0.9185}$ & 1064 \\
$f_2 = 1/\text{NNS}_\mathrm{unf}$ & 0.8806 & 696 \\
$f_3 = \log(\text{NNS}_\mathrm{unf})$ & 0.5182 & 100 \\
\bottomrule
\end{tabular}
\end{center}
Lorentzian$^2$ 모델은 $R^2$를 0.015에서 0.919로 끌어올려,
이전에 설명되지 않은 분산의 $90.3\%$를 설명한다.
이론 상관 $|\operatorname{Im} R| \sim 1/g_{\min}$은
$r = 0.892$ ($p < 10^{-34}$)로 확인된다.

\emph{이상치 해부.}\; $|{\mathrm{res}}| > 2\sigma$인
2개 이상치($N = 11919, 16767$)는 모두
$\text{NNS}_\mathrm{unf} < 0.15$---GUE 레벨 인력 사건을
100\% 포착한다.

\emph{주의.}\; KS 검정은 표본 NNS 분포와 Wigner surmise의 불일치를
기각한다($D = 0.41$, $p < 10^{-15}$). 등간격 인덱스 선택(무작위
표본추출 아님)에 기인하며, 상관 분석의 유효성에는 영향을 미치지
않으나 분포 수준 주장은 제한된다.

\textcolor{ForestGreen}{[확인 2026-04-15; 99 영점 ($N{=}203$ \textasciitilde{} $20000$);
$\delta{=}0.03$; $R^2_\text{1-var}{=}0.015$, $R^2_\text{2-var}{=}0.919$;
Hadamard 분해 오차 ${=}0.00\%$;
3-시드 예비]}

\smallskip\noindent
\textbf{파트~F: $\delta$-sweep 보편성 ($\delta \in \{0.01, 0.02, 0.03, 0.05, 0.10\}$).}\;
Lorentzian$^2$ 보정이 탐침 거리에 무관하게 보편적인지 검증하기 위해,
5개 오프셋으로 확장한다 (각 99 영점, $N = 203$--$20\,000$).

\emph{$\delta$별 2변수 회귀.}\;
\begin{center}
\small
\begin{tabular}{@{} c c c c @{}}
\toprule
$\delta$ & $R^2_\text{1-var}$ & $R^2_\text{2-var}$ & $c(\delta)$ \\
\midrule
0.01 & 0.014 & \textbf{0.925} & 147.8 \\
0.02 & 0.015 & \textbf{0.923} & \ 39.4 \\
0.03 & 0.015 & \textbf{0.919} & \ 19.4 \\
0.05 & 0.017 & \textbf{0.909} & \  9.1 \\
0.10 & 0.024 & \textbf{0.890} & \  4.9 \\
\bottomrule
\end{tabular}
\end{center}
Lorentzian$^2$ \emph{형태}는 5개 오프셋 전부에서 $R^2 > 0.88$을 달성하여,
$\delta$의 10배 범위에 걸쳐 보정 \emph{형태}의 보편성을 확립한다.

\emph{음성 결과: 계수 비보편성.}\;
계수 $c(\delta)$는 $147.8$ ($\delta = 0.01$)에서 $4.9$ ($\delta = 0.10$)로
30배 변동한다 --- 비자명 지수 $\alpha$에 대한 $c(\delta) \sim \delta^{-\alpha}$를 시사한다.
곱 $a(\delta)\cdot\delta$의 변동계수는 $47.8\%$ (이론 예측: $\lesssim 20\%$).
통합 3-파라미터 모델
$\Delta\kappa = (a/\delta)\log(t/2\pi) + b/\delta^2 + c\cdot f_1$은
$R^2 = 0.449$---기준 $0.90$ 미달. 이 실패는 계수의 $\delta$-스케일링이
단순 멱급수가 아님을 나타내며 추가 해석적 규명이 필요하다.

\emph{$\delta = 0.01$에서의 수치 안정성.}\;
$1/\delta^2 = 10\,000$으로, 차감 $\Delta\kappa = \kappa - 10\,000$은
자릿수 상쇄 위험이 있다. $\text{dps} = 80$ 상향 시
$R^2_{2\text{-var}} = 0.925$이고 $\Delta\kappa > 0$ 전원: 수치 아티팩트 없음.

\textcolor{ForestGreen}{[확인 2026-04-15; 5 오프셋 $\delta\in\{0.01,0.02,0.03,0.05,0.10\}$;
각 99 영점 ($N{=}203$--$20000$);
Lorentzian$^2$ 보편 ($5/5$, $R^2{>}0.88$);
통합 3-파라미터 모델 $R^2{=}0.449$ (음성)]}
\end{observation}


% ----------------------------------------------------------------------------
\begin{observation}[RMT 연관의 한계]
\label{obs:rmt_boundaries}
세 결과(\#20 쌍 상관, \#28 수 분산, \#32 NNS--$\Delta\kappa$
Lorentzian$^2$)에서 RMT 양성 연관을 확립한 후, 곡률장 $\kappa(t)$가
개별 분포 및 상관 함수 수준에서 더 세밀한 RMT 예측을 따르는지
검증하는 네 가지 실험을 수행하였다.
이 더 세밀한 예측들이 \emph{성립하지 않음}을 확립함으로써,
RMT 유비의 적용 범위를 정밀하게 획정한다.

\smallskip\noindent
\textbf{결과~\#34 (실험~\#36): 간격비 분포.}\;
비율 $\tilde{r}_n = \min(s_n,s_{n+1})/\max(s_n,s_{n+1})$을
$t > 600$ 범위의 리만 영점 $N = 4657$개 연속 간격으로 구성한다.
이론 평균을 $\mathbb{E}[\tilde{r}]_\text{GUE} = 0.6027$
(Atas \textit{et al.}\ 2013)로 교정한 후, 실측 평균은 $\bar{r} = 0.6161$로
$+2.2\%$ 체계적 초과($4.15\sigma$)를 보인다.
GUE Wigner surmise에 대한 KS 검정: $D = 0.030$,
$p = 4.7 \times 10^{-4}$.
비율 $D_\text{GUE}/D_\text{Poisson} = 0.079$는 간격비 분포가
\emph{GUE에 가장 근접}함을 확인하지만,
체계적 편향은 모든 $t$-절단값($t > 400, 600, 800, 1000$)에서
지속되며 $N$ 증가에 따라 GUE로 수렴하지 않는다.
\textbf{판정: 음성} (모든 절단값에서 $p < 5 \times 10^{-4}$;
GUE로의 수렴 부재).

\smallskip\noindent
\textbf{결과~\#35 (실험~\#37): 곡률 값 분포.}\;
변수 변환 $\kappa = 1/(\delta^2 + s_\text{unf}^2)^2$ 하에서,
$s_\text{unf}$가 GUE 레벨 간격 분포를 따른다면 $\kappa$는
예측 가능한 분포 $P_\text{pred}(\kappa)$를 따라야 한다.
$t \in [1000, 5000]$, $\delta = 0.03$에서 유효 표본 $N = 1870$개로
검증한다.
KS 검정: $D = 0.361$, $p = 6.3 \times 10^{-213}$.
로그 스케일 Q--Q 상관 $r = 0.971$은 분포의 \emph{형태}가 보존되지만
\emph{스케일}이 중앙값에서 $7.6\times$ 불일치함을 나타낸다.
근본 원인은 두 가지이다: (i)~균등 $t$ 샘플링이 GUE 이론 기대값 $1.000$
대신 $\bar{s}_\text{unf} = 0.306$을 산출하며,
(ii)~$\kappa = |G(s) + \zeta'/\zeta(s)|^2$의 배경항 $G(s)$가
$|s - \gamma_n| \gg \delta$ 거리에서 지배적 극 $1/\delta^2$를
부분적으로 상쇄하여, 단순 모델 $1/(\delta^2 + s^2)^2$에서 포착되지
않는 체계적 스케일 오프셋을 도입한다.
\textbf{판정: 음성} (KS $p < 10^{-200}$; 로그 스케일 형태 보존,
스케일 $7.6\times$ 불일치).

\smallskip\noindent
\textbf{결과~\#36 (실험~\#38): $\kappa$ 자기상관.}\;
$N = 8001$개 등간격 점 ($\Delta t = 0.5$, $t \in [1000, 5000]$)에서
경험적 자기상관 함수
$C(\tau) = \operatorname{Corr}(\kappa(t), \kappa(t+\tau))$를 계산하고,
중점 평균 영점 밀도 $\bar{d} = 0.982$를 이용하여
$\tau_\text{unf} = \tau \cdot \bar{d}$로 변환한다.
GUE 두점 군집 함수 $b_2^c(r) = 1 - (\sin\pi r/\pi r)^2$는
첫 번째 영교차가 $\tau_\text{unf} = 1.0$에서 발생함을 예측한다.
실측에서 첫 번째 영교차는 $\tau_\text{unf} = 0.448$
(편차 $0.552$, 허용 기준 $\pm 0.20$ 초과)에서 발생한다.
진폭 피팅 결과 $A = 0.190$: $C_\text{emp}$의 진동 진폭이 GUE 예측의
$19\%$에 불과하다.
전체 지연 범위($\tau_\text{unf} \leq 20$): $R^2 = 0.48$.
결정적으로, $C(1) = -0.081$: 자기상관이 지연~1
($\tau_\text{unf} \approx 0.48$)에서 이미 거의 0에 도달하여,
\emph{$\kappa(t)$가 격자 간격에서 사실상 비상관}임을 확립한다.
\textbf{판정: 음성} (영교차 $\tau_\text{unf} = 0.448 \neq 1.0$;
$R^2 = 0.48$; 지연 $\geq 1$에서 $\kappa$ 사실상 비상관).

\smallskip\noindent
\textbf{해석: 대국적 통계 vs.\ 미시적 분포.}\;
확립된 RMT 양성 결과(\#20 쌍 상관, \#28 수 분산)는 영점 \emph{위치} 자체의
두점 통계로, GUE 보편성이 견고한 스펙트럼의 평균 밀도 구조를 포착한다
(Berry--Tabor, Bohigas--Giannoni--Schmit).
음성 결과(\#34--\#36)는 대신 곡률 $\kappa$의 \emph{전체 분포} 또는
\emph{시간적 자기상관}을 탐침하는데, 이는 GUE 모델에 없는 추가 요소를
포함하는 영점 위치의 비선형 범함수이다:
\begin{itemize}
  \item 감마·극 기여로부터의 배경항 $G(s)$
        ($|G(s)|^2 \approx 6$--$9$)가 $|s - \gamma_n| \gg \delta$
        거리에서 지배적 극 $1/\delta^2$를 부분적으로 상쇄하여,
        단순 GUE 대입에서 보이지 않는 체계적 스케일 오프셋을 도입한다.
  \item 곡률 $\kappa(t)$는 최근접 영점 위치에 의해 \emph{국소적으로 결정}된다.
        공간 상관 길이는 격자 간격 수준($\Delta t = 0.5$,
        $\tau_\text{unf} \approx 0.48$)으로, GUE 상관 길이
        ($\tau_\text{unf} \approx 1$)보다 훨씬 짧다.
\end{itemize}

따라서 다음 경계를 확립한다:
\begin{quotation}\noindent
  \emph{RMT 연관은 스펙트럼 두점 통계(쌍 상관, 수 분산) 수준과
  NNS--곡률 범함수 상관(Lorentzian$^2$ 모형) 수준에서 작동하지만,
  간격비 분포 전체, 곡률 값 분포, 곡률 범함수의 시간적 자기상관으로는
  확장되지 않는다.}
\end{quotation}
이 경계는 정밀하고 유익한 과학적 발견이지 실패가 아니다: $\xi$-다발
기하학의 어떤 측면이 랜덤 행렬 보편성의 지배를 받고, 어떤 측면이
산술적 배경항에 민감한지를 정확히 식별한다.

\textcolor{red}{[음성, 확인 2026-04-15;
\#34: KS $p = 4.7\times10^{-4}$, 체계적 초과 $+2.2\%$ (모든 절단값);
\#35: KS $D = 0.361$, 스케일 $7.6\times$, 로그 Q--Q $r = 0.971$;
\#36: 영교차 $\tau_\text{unf} = 0.448 \neq 1.0$, $R^2 = 0.48$;
독립 검증: 10/10 항목 재현]}
\end{observation}


% [Paper B: 스펙트럼 통계]
% ============================================================================
\section{구조적 대응 표}
\label{sec:correspondence}
% ============================================================================

본 연구 전반에 걸쳐 식별된 7개의 구조적 대응을 형식화한다:

\begin{table}[ht]
\centering
\caption{QROP-Net과 공명 프레임워크 사이의 7개 구조적 대응.}
\label{tab:correspondence}
\rowcolors{2}{lightbg}{white}
\small
\begin{tabular}{@{} c p{3.5cm} p{3.5cm} p{4.5cm} @{}}
\toprule
\textbf{\#} & \textbf{QROP-Net (구체적)} & \textbf{공명 (해석적)}
& \textbf{공유 수학 구조} \\
\midrule
1 & $\sin^2(2^{d-1}\phi)$ 양자화 손실
  & $\Delta v = m\pi$ 위상 양자화 (게이트~A)
  & 주기적 위상 벌칙: $\mathcal{P}(\phi;L) = \sin^2(L\phi/2) = 0$
    (격자 위치에서) \\

2 & NTT (이산) $+$ FrFT (연속 분수)
  & 임계선 위 $\xif(s)$의 푸리에 분석
  & 스펙트럼 분해: 대수 구조 $\leftrightarrow$ 주파수 영역 \\

3 & $X^n + 1 = 0$ 근: 단위원 위 $n$개 등간격 위상
  & 위상 감김 $\oint\nabla v \cdot ds = 2\pi N$
  & 시클로토믹/원형 양자화: $S^1$ 위의 이산 위상 \\

4 & LWE 오차 임계치: $|e|_{\Zq} < q/4 = 1920$
  & $\sigma$-리프트: $|\sigma^* - 1/2| = 0$ (RH)
  & 신호 복원이 실패하는 경성 오차 경계 \\

5 & $U_0 = e^{i\phi}$, $|U_0| \equiv 1$; 세기 $I = |U_d|^2$
  & $\Ampl(t) = |\xif|$ vs $\Phiz(t) = \arg\xif$
  & 진폭--위상 분리: 위상에 정보 인코딩, 진폭으로 측정 \\

6 & 블라인드 TTO: 정답 없는 자가 교정
  & 게이지 ODE: 영점 위치 모른 채 $F_2 = 0$
  & 내적 일관성에 의한 정답 없는 수렴 \\

7 & $\Delta d = 0.1\,$mm $\to$ DFR $= 1$
  & $\sigma$-리프트 $\neq 0$ $\to$ 비유니타리성 (흡수)
  & 민감도 임계치: 작은 매개변수 이동 $\to$ 파국적 실패 \\
\bottomrule
\end{tabular}
\end{table}

\subsection{대응의 수학적 형식화}

각 대응을 구조적 사상으로 표현한다.

\begin{definition}[위상 양자화 함자]
\label{def:functor}
$(\mathcal{A}, \Phi_{\mathcal{A}}, L_{\mathcal{A}})$를 대수 구조 $\mathcal{A}$, 위상 함수 $\Phi_{\mathcal{A}}: \mathcal{A} \to [0, 2\pi)$, 격자 수준 수 $L_{\mathcal{A}} \in \mathbb{N}$으로 이루어진 \textit{위상 양자화 삼중체}라 하자. 위상 양자화 연산자 $\mathcal{P}(\cdot; L)$은 모든 이러한 삼중체에 동일하게 작용한다:

\medskip
\begin{center}
\begin{tabular}{lccc}
\toprule
& $\mathcal{A}$ & $\Phi_{\mathcal{A}}$ & $L_{\mathcal{A}}$ \\
\midrule
QROP-Net   & $\Zq^{1024}$ & $\phi_i = 2\pi c_i/2^{d_u}$
           & $2^{d_u} = 2048$ \\
시클로토믹 & $R_q$ & $\arg(\omega_j) = \pi(2j+1)/n$
           & $2n = 512$ \\
임계선     & $\{t \in \mathbb{R}^+\}$
           & $\vartheta(t) \bmod \pi$
           & $2$ (그람 구간당) \\
\bottomrule
\end{tabular}
\end{center}
\medskip

각 경우에, 조건 $\mathcal{P}(\Phi_{\mathcal{A}}; L_{\mathcal{A}}) = 0$이 ``정렬된'' 원소를 선택하며, $\nabla\mathcal{P}$가 비정렬 원소를 가장 가까운 격자 위치로 유도한다.
\end{definition}

\subsection{\textcolor{orange}{\rule{0.9ex}{0.9ex}}\ 동역학적·스펙트럼적 대응}
\label{sec:extended_corr}

\noindent\textbf{계층 2.} 위의 대응표는 구조적이다. 본 소절은 \cref{sec:f2_hardy}--\cref{sec:high_zeros}의 실증 현상을 고전 수학 틀에 연결하는 네 가지 추가 대응을 기록한다. 각 대응은 유비와 부분적 수치 일치 수준에서 수립되며, 엄밀한 동등성을 주장하지 않는다.

\subsubsection{Maslov 보정 --- 메타플렉틱 시트 이동}
\label{sec:maslov_corr}

\cref{sec:f2_hardy} 전반에서 언급된 ``$\cos^2$ Maslov 보정'' --- 양자화 벌칙을 $\sin^2(L\phi/2)$에서 $\cos^2(L\phi/2)$로 교체하는 것 --- 은 고전적 메타플렉틱 현상의 신경망적 그림자이다.

\begin{theorem}[게이트~A의 메타플렉틱 보정]
\label{thm:maslov}
실함수의 영점에서 $\varphi\to\pi/2$ (\cref{cor:neural_reg})일 때 적용되는 위상 양자화 연산자 $\PQO(\varphi; L)$에 대하여, $\sin^2$ 형식은 $\PQO(\pi/2; 2) = 1$ (최대 벌칙) 을 주는 반면, $\cos^2$ 형식은
\begin{equation}\label{eq:pqo_cos_thm}
  \PQO_{\cos}(\varphi; L) := \cos^2\!\bigl(L\varphi/2\bigr),
  \quad \PQO_{\cos}(\pi/2; 2) = 0
\end{equation}
을 준다. 교체 $\sin^2\to\cos^2$은 지수 $\mu=2$를 가지는 \emph{Maslov 보정}이며, 이는 위상 이동 $\mu\pi/4=\pi/2$에 해당한다.
\end{theorem}

\begin{table}[ht]
\centering
\caption{수학 전반에 걸친 $\sin^2\leftrightarrow\cos^2$ 쌍대성.}
\label{tab:maslov_duality}
\small
\begin{tabular}{@{}p{3cm}|p{3.8cm}|p{3.8cm}|p{3cm}@{}}
\toprule
\textbf{영역} & \textbf{$\sin^2$ 국면} & \textbf{$\cos^2$ 국면} & \textbf{이동의 이름} \\
\midrule
WKB/준고전     & 전환점에서 떨어진 곳
               & 전환점에서
               & Maslov 지수 $\mu$ \\
심플렉틱 기하  & 표준 라그랑지안 차트
               & 교차하는 Maslov 순환
               & $\pi_1(\Lambda(n))$ \\
조화 해석      & 동상 성분 (I)
               & 직교 성분 (Q)
               & 힐베르트 변환 \\
격자 암호      & 격자 복호
               & 잉여류 복호
               & $\Lambda\to\Lambda+v$ \\
Berry--Keating & $\vartheta(t_n)=n\pi$
               & $\vartheta(t_n)=(n-\tfrac12)\pi$
               & 반정수 이동 \\
표현론         & 콤팩트 카르탕
               & 분할 카르탕
               & Cayley 변환 \\
\bottomrule
\end{tabular}
\end{table}

통합 대상은 $\mathrm{Sp}(2n,\mathbb{R})$의 이중 피복인 메타플렉틱 군 $\mathrm{Mp}(2n,\mathbb{R})$이며, 그 두 시트는 정확히 두 양자화 위상에 대응한다. 함수방정식 $\xif(s)=\xif(1-s)$는 $\tau\to-1/\tau$을 통해 세타 함수 위에서 작용하는 이 군에서 유래하며, Berry--Keating 양자화 조건
\begin{equation}\label{eq:berry_keating}
  \vartheta(t_n)=\bigl(n-\tfrac12\bigr)\pi + O(1/t_n)
\end{equation}
는 해석적 측면에서 동일한 $-1/2$ 이동을 부호화한다.

\subsubsection{게이지 ODE 임계 전이 --- Kuramoto 동기화로서}
\label{sec:kuramoto}

\begin{observation}[게이지 임계 전이]
\label{obs:gauge_transition}
\cref{sec:f2_hardy}의 $F_2$-탐지기에서 게이지 버퍼 상태 $(\varphi,\psi)$는 ${\sim}10$ epoch에서 급격한 전이를 보인다: 표준편차 $\mathrm{std}(\varphi)$가 $\approx 0.10$에서 $\approx 0.03$으로 떨어지며, 이는 $F_2$가 모든 훈련 영점에서 수렴에 도달하는 epoch과 일치한다. 이 전이는 시드에 불변이며 $t\in[10,50]$과 $t\in[100,200]$ 모두에서 재현된다.
\end{observation}

\begin{remark}[Kuramoto 동기화]
\label{rem:kuramoto_transition}
\cref{obs:gauge_transition}은 유한 모집단 Kuramoto형 동기화 전이이다. 게이지 위상 $\{\varphi_i\}$은 배치 지역 장이며, 공유되는 잔차 손실 $\mathcal{L}_{\mathrm{res}}$가 이들 사이의 결합 범함수로 작용한다. 유효 결합 강도 (학습률과 타겟의 날카로움에 의해 제어됨) 가 임계값을 초과할 때, 위상들은 타겟 영점 위치에 의해 결정되는 공통 주파수 프로파일로 잠긴다. epoch $\sim 10$에서의 $\mathrm{std}(\varphi)$의 $\approx 3\times$ 하락은 이 잠김의 신경망적 서명이다.

\textcolor{ForestGreen}{[confirmed 2026-04-12]}
Kuramoto 순서 매개변수 $r(e)=|N^{-1}\sum_j e^{i\varphi_j}|$를
$3\text{ seeds}\times 2\text{ }t$-범위에서 측정하였다
(5회 완료 + 1회 epoch 90/100에서 OOM 불완전 종료; 불완전 run도 동일 단조 패턴).
$\mathrm{std}(\varphi)$ 비율은 $3.43\pm0.09\times$
(초기 $0.108\pm0.008$, 최종 $0.032\pm0.002$)로,
위의 ${\sim}3\times$ 관찰과 정량적으로 일치한다.
순서 매개변수는 epoch~$0$부터 $r>0.99$
($r_{\mathrm{init}}=0.994\pm0.001$, $r_{\mathrm{final}}=0.9994\pm0.0001$)를
만족하며, 네트워크 초기화가 이미 게이지 위상을 초임계 영역 $K\gg K_c$에
배치함을 나타낸다; 학습은 무질서-질서 전이가 아닌 동기화의 정련이다.
$t\in[10,50]$와 $t\in[100,200]$ 모두 동일 패턴을 보여 $t$-범위 불변성을 확인하였다.
\end{remark}

\begin{remark}[von Mises 자기일관성]
\label{rem:von_mises}
최종 위상 분산 $\sigma_\varphi = 0.032 \pm 0.002$는
집중 매개변수 $\kappa = 1/\sigma_\varphi^2 \approx 980$인
von~Mises 분포와 정확히 일치한다: 자기일관성 관계
\[
  r \;=\; e^{-\sigma^2/2} \;=\; e^{-0.032^2/2} \;\approx\; 0.9995
\]
는 측정된 $r_{\mathrm{final}} = 0.9994 \pm 0.0001$과
실험 오차 내에서 일치한다. 이는 학습된 게이지 위상 앙상블이
\emph{거의 축퇴된} von~Mises 분포임을 확인하며,
Kuramoto 동기화가 단순한 비유가 아님을 확립한다:
위상들은 $\kappa \approx 10^3$인 von~Mises 분포를
통계적으로 따른다.
\textcolor{ForestGreen}{[confirmed 2026-04-13]}
\end{remark}

\begin{proposition}[자발적 게이지 고정]
\label{prop:spontaneous_gauge}
$H$를 은닉 채널 수라 하자.
$H$개의 게이지 위상 $\{\varphi_j(x)\}_{j=1}^H$은
초기화 시 $r_{\mathrm{init}} \approx 1 - c\,H^{-\alpha}$를 만족하며,
$\alpha \approx 0.38$, $c \approx 0.033$이다 (5-시드 피팅, $R^2=0.94$).
학습은 집중도를 $r_{\mathrm{final}} > 0.998$로 균일하게 정련하여,
자발적 대칭 파괴 $U(1)^H \to U(1)$을 달성한다.

\smallskip\noindent\textit{$H$-스케일링 데이터
(5 시드 $\times$ 5 너비, $t\in[100,200]$, 30 에포크):}

\smallskip
\begin{center}
\begin{tabular}{rcccc}
\toprule
$H$ & $r_{\mathrm{init}}$ & $1{-}r_{\mathrm{init}}$ &
      $r_{\mathrm{final}}$ & $\Delta r/\mathrm{ep}\;(\times 10^{-3})$ \\
\midrule
 16 & $0.9878\pm0.0016$ & $0.0122$ & $0.9982\pm0.0003$ & $0.348\pm0.043$ \\
 32 & $0.9915\pm0.0011$ & $0.0085$ & $0.9985\pm0.0001$ & $0.236\pm0.038$ \\
 64 & $0.9938\pm0.0010$ & $0.0062$ & $0.9988\pm0.0002$ & $0.166\pm0.027$ \\
128 & $0.9954\pm0.0003$ & $0.0046$ & $0.9987\pm0.0001$ & $0.110\pm0.010$ \\
256 & $0.9955\pm0.0005$ & $0.0045$ & $0.9989\pm0.0001$ & $0.112\pm0.014$ \\
\bottomrule
\end{tabular}
\end{center}

\smallskip\noindent
스케일링은 $H\ge128$에서 포화한다 ($1{-}r_{\mathrm{init}}\approx 0.0045$).
단순 $O(1/H)$ 법칙은 기각됨. $H=16$에서조차
$r_{\mathrm{init}}=0.988 \gg 1/\sqrt{16}=0.25$:
이 근접-단위 동기화는 유한 크기 효과가 아니라 \emph{구조적}이다
(Xavier 초기화와 ReLU 비선형성이 출력 위상 다양성을 억제).
정련 속도 $\Delta r/\mathrm{ep}$는 $H$에 대해 단조 감소하며,
이는 $r_{\mathrm{init}}\to 1$일 때 한계 이득이 줄어드는 것과 일치한다.
\textcolor{ForestGreen}{[confirmed $H{\in}\{16,32,64,128,256\}$,
2026-04-13]}
\end{proposition}

\begin{remark}[von Mises--Kuramoto--$L_{\mathrm{geo}}$ 삼각 연결]
\label{rem:triangle}
측지 손실 $L_{\mathrm{geo}} = 1 - \cos(\Delta\varphi)$는
가산 상수를 제외하면 von~Mises 분포의 음의 로그우도이다:
$-\log p(\varphi \mid \mu,\kappa)
  = -\kappa\cos(\varphi-\mu) + \mathrm{const}$.
$L_{\mathrm{geo}}$를 최소화하면 von~Mises MLE를 수행하게 되며,
이는 $\kappa$를 최대화하고, 이는 다시 Kuramoto
순서 매개변수 $r \to 1$을 구동한다. 세 개념---von~Mises MLE,
$S^1$ 측지 손실, Kuramoto 동기화---은 따라서 하나의
최적화 문제의 서로 다른 관점이다. 이것이
$\Ltgt$를 $L_{\mathrm{geo}}$로 교체하면
(\cref{sec:open}, Tier~1) 검출 개선이 나타나는 이유를
설명한다: 손실 함수가 게이지 동기화를 암묵적으로
강화하기 때문이다.
후속 극소값 기반 평가(5시드, 2개 높이 구간)는
\emph{밀도 의존적} 효과를 보여준다:
\begin{itemize}
  \item \textbf{$t\in[500,600]$} (밀도 $0.72$/단위):
    $F_1$ 동등 (BL $0.373\pm0.020$ vs $L_{\mathrm{geo}}$
    $0.361\pm0.012$, $p{=}0.45$);
    재현율 분산 축소 유의
    ($F{=}16.2$, $p{=}0.010$).
  \item \textbf{$t\in[1000,1100]$} (밀도 $0.81$/단위):
    $F_1$ 개선 (BL $0.395\pm0.009$ vs $L_{\mathrm{geo}}$
    $0.411\pm0.005$, 대응 $t{=}2.84$, $p{=}0.047$,
    5/5시드 방향 일관);
    재현율 분산도 유의 ($F{=}6.5$, $p{=}0.049$).
\end{itemize}
두 구간 모두에서 $L_{\mathrm{geo}}$는 재현율 하한을
안정화한다 ($t\in[500,600]$에서 $R\ge 0.972$;
$t\in[1000,1100]$에서 $R\ge 0.975$; 10개
시드--구간 조합 전체).
더 높은 영점 밀도에서 측지 손실은 추가로 평균 $F_1$을
끌어올리며, 이는 $\Ltgt$의 $\pm\pi$ 불연속 부담이
밀도에 따라 증가함을 시사한다.
효과 크기는 소($\Delta F_1 \approx 0.016$)이고
$p$-값은 경계선이다.
파국적 발산은 1/5시드(시드 특이적, best-checkpoint
선택으로 구제).
\textcolor{ForestGreen}{[2026-04-13 확인, 5시드, 2구간]}
\end{remark}

\subsubsection{Lehmer 현상과 GUE 간격 통계}
\label{sec:lehmer_gue}

\cref{tab:f2_low}의 난이도 위계 --- 일부 영점이 이웃보다 한두 자릿수 느리게 수렴하는 것 --- 는 Lehmer 쌍의 경험 분포와 일치한다. Montgomery의 쌍 상관 추측과 Odlyzko의 수치 연구는 고위 리만 영점의 정규화된 간격이 GUE 분포
\begin{equation}\label{eq:gue}
  p_{\mathrm{GUE}}(s) = \frac{32}{\pi^2}s^2 e^{-4s^2/\pi}
\end{equation}
를 따름을 확립하였고, 가까운 영점 쌍 --- Lehmer 쌍 --- 은 예측 가능한 빈도로 나타난다. 우리 탐지기에서 이들은 $Z(t)$가 실수축을 거의 스치지만 교차하지 않는 영점이며, 해석 신호가 $\pm\pi/2$ 근방에 더 오래 머물고, 다발 곡률 프로파일이 정성적으로 더 복잡해진다(넓은 FWHM, 다중 피크). 그러나 Hadamard 부분합 수렴 속도는 영점 간격에 둔감하다($|r|<0.03$). 이로부터 검증 가능한 예측이 나온다: 고정 epoch 예산에서 영점별 수렴 시간 분포가 GUE 밀도와 해석 신호 체류 시간의 컨볼루션과 일치해야 한다. 이 양은 닫힌 형식으로 계산될 수 있다. 아직 이 적합을 수행하지 않았으며, \cref{sec:open}의 열린 문제로 기록한다.

\subsubsection{Weyl 등분포와 PGGD 각도 보행}
\label{sec:weyl}

\cref{sec:pggd_failure}는 PGGD의 실패 원인을 $2\pi$ 모듈로 각도 업데이트의 Weyl형 등분포 탓으로 돌린다. 기저 사실은 무리수 $\alpha$에 대해 수열 $\{n\alpha\bmod 1\}$이 $[0,1)$에서 등분포한다는 고전 정리이다. PGGD 맥락에서 스텝당 각도 업데이트는 근사적으로 $\eta\,\hat L(t)$이며, 일반적 타겟 함수에 대해 유한 샘플 위에서 유리 독립이다; 접 투영은 전체 업데이트를 이 각도 성분에 집중시키며, epoch 수가 증가함에 따라 경험 분포가 $S^1$ 위의 Haar 측도로 수렴하는 등분포 보행을 생성한다. 이는 \cref{rem:kuramoto_transition}의 잠김의 역이다: 여기서는 결합이 구성에 의해 제거되었고, 등분포가 승리한다.


% [Paper A: ξ-다발 기하학]
% ============================================================================
\section{\textcolor{orange}{\rule{0.9ex}{0.9ex}}\ $|\Ftwo|$ 잔차: 잠정적 중시각 관측량}
\label{sec:f2_residual}
% ============================================================================

\noindent\textbf{계층 2 (2026-04-08 부분 재해석 포함).}
제~\ref{sec:high_zeros}절의 고위영점 확장을 실행하면서, 개별 영점에서의 수렴 후 크기 $|F_2(t_k)|$가 무특징적 임의 산포가 아니라 체계적 높이 및 공간 패턴을 보인다는 점을 발견하였다. 본 섹션은 원래 관찰 (2026년 4월, $14{,}441$개 영점)과 이후 2026-04-08 부분 재해석을 기록한다. 재해석은 개별 영점 수준의 주장은 약화시키지만 총괄적 스케일 주장은 건드리지 않는다. 독자가 주장의 범위가 어떻게 좁혀졌는지 볼 수 있도록 양쪽 모두 보존한다.

\subsection{원래 관찰 (2026년 4월)}
\label{sec:f2_residual_original}

$|F_2|$ 잔차를 $[10,50]$에서 $[5000,10000]$까지의 7개 광역 구간과 $[10000,13000]$을 덮는 23개 세부 구간을 합쳐 총 $14{,}441$개 영점에서 계산하였다. 누적 탐지율 (균일 임계치 이하인 $|F_2|$를 갖는 영점)은 $99.99\%$였다.

\begin{table}[ht]
\centering
\caption{원래의 $14{,}441$ 영점 스윕에서 얻은 높이 구간별 $|F_2|$ 평균 잔차.}
\label{tab:f2_scale}
\small
\begin{tabular}{lrrr}
\toprule
\textbf{구간} & \textbf{영점 수} & \textbf{탐지율}
  & $\overline{|F_2|}$ \\
\midrule
$[10, 50]$       & 10     & 10/10     & 0.0209 \\
$[100, 200]$     & 50     & 49/50     & 0.0208 \\
$[200, 500]$     & 190    & 190/190   & 0.0119 \\
$[500, 1000]$    & 380    & 380/380   & 0.0083 \\
$[1000, 2000]$   & 868    & 868/868   & 0.0058 \\
$[2000, 5000]$   & 3{,}003  & 3{,}003/3{,}003 & 0.0020 \\
$[5000, 10000]$  & 5{,}622  & 5{,}622/5{,}622 & 0.0014 \\
\bottomrule
\end{tabular}
\end{table}

이 데이터셋에서 원래 세 가지 개별 영점 수준 주장이 제기되었다:
\begin{enumerate}
  \item[(O1)] \textbf{고전 통계로부터의 독립성.} 개별 영점 잔차 $|F_2(t_k)|$는 $|Z'(t_k)|$, $|Z''(t_k)|$, 정규화 간격 $s_k$, 인수 함수 $S(t_k)$, 리만--지겔 도함수 $\vartheta'(t_k)$, 국소 밀도, 간격 분산, 간격 자기상관, 그리고 국소 $Z$-면적 $\int_{t_k-0.3}^{t_k+0.3}|Z|$ 모두에 대해 Pearson $|r|<0.09$와 Spearman $|\rho|<0.09$를 보였다.
  \item[(O2)] \textbf{공간 클러스터링.} $|F_2(t_k)|$는 가장 가까운 6개 영점 이웃의 $|F_2|$ 평균과 Spearman $\rho=+0.895$ ($p\approx 10^{-141}$)로 상관되었으며, 이는 임계선 위에 서로 다른 성격의 ``중시각 도메인''이 존재한다는 증거로 해석되었다.
  \item[(O3)] \textbf{장거리 자기상관.} 시차-$\Delta$ 자기상관 $C(\Delta)$는 $\Delta\sim 3$ 내에서 $\approx 0.4$로 급격히 떨어지고 $\Delta=200$까지 $\approx 0.33$에서 거의 평탄한 상태로 유지되며, 두 성분 구조 (단거리 지수 + 장거리 평탄)를 시사한다.
\end{enumerate}

종합하여 이것들은 게이지 ODE가 ``임계선의 새 렌즈''로서 작용하여 무작위 행렬 이론과 고전적 $S(t)$-기반 도구로는 보이지 않는 구조를 드러낸다는 증거로 읽혔다.

\subsection{2026-04-08 부분 재해석}
\label{sec:f2_residual_reinterp}

\begin{remark}[$|F_2|$ 개별 영점 주장의 부분 재해석]
\label{rem:f2_reinterp}
2026-04-08에 수행된 모델 독립성 감사는 (O1)--(O3)을 \emph{개별 영점} 수준에서 상당히 약화시키는 반면, \cref{tab:f2_scale}의 총괄적 스케일 구조는 건드리지 않았다.

구체적으로, 동일 구조를 다른 시드, 다른 게이지 버퍼 $(\varphi_0,\psi_0)$ 초기화, 다른 학습 영점 목록 셔플로 재학습하면 개별 영점의 $|F_2(t_k)|$ 값의 순위 정렬이 실행 간에 재현되지 \emph{않는다}. 동일 구간에 대한 독립적 두 실행의 개별 영점 Spearman 상관은 대부분의 구간에서 $+0.2$ 이하로 떨어지며, 이는 $|F_2(t_k)|$를 영점 $t_k$의 고유 성질로 해석하는 것과 불일치한다. 특히, 공간 클러스터링 통계 $\rho_{\mathrm{cluster}}=+0.895$는 이웃 영점 잔차를 결합시키는 학습 시간 매끄러움 사전 분포에 의해 부분적으로 부풀려져 있고, $C(\Delta)$의 장거리 평탄치 $\approx 0.33$은 모든 실행에서 나타나지만 학습 의존적 수준으로 나타난다.

재해석에서 살아남는 것은 \cref{tab:f2_scale}의 구간들에서 나타나는 $\overline{|F_2|}$의 높이 의존적 \textbf{평균 스케일}이다: $t\sim 50$에서의 $\overline{|F_2|}\approx 0.021$로부터 $t\sim 10{,}000$에서의 $\overline{|F_2|}\approx 0.0014$에 이르는 단조 감소가 모든 시드에서 상대 변동 ${\sim}10\%$ 이내로 재현되며, 그 형태는 $\alpha\in[0.4,0.5]$인 멱법칙 $\overline{|F_2|}\sim t^{-\alpha}$와 일치한다.

따라서 \cref{tab:f2_scale}과 평균 스케일 주장은 탐지기 잔차의 진정한 높이 의존성에 대한 계층 2 증거로 \emph{유지}하며, 개별 영점 주장 (O1)--(O3)은 ``시사적이지만 모델 오염된'' 상태로 \emph{강등}한다: 패턴은 학습된 모델의 진정한 특징이지만, 모델로부터 임계선의 근저 산술에 이르는 추론은 개별 영점 수준에서 관찰된 비재현성에 의해 차단된다. 본 섹션은 바로 이 축소 자체가 정보 가치를 가진 진단이기 때문에 이 논문에 유지된다; 그러한 철회 기록의 역할에 관해서는 \cref{sec:open}도 참조.
\end{remark}

\subsection{2026-04-10 후속: ``어려움 군집''은 기저 대역폭 잡음}
\label{sec:f2_residual_cluster_artifact}

2026-04-09에서 2026-04-10에 걸쳐, $|F_2|$ 데이터 내부의 잔차 어려움 군집에 대해 열 개의 실험 연쇄 (Lehmer 회귀, 5-seed 앙상블, hidden 폭 스케일링, 표본 밀도 스케일링, 잔차 스펙트럼, 군집 물리, 입력 기저 스케일링, 수렴 확장)를 수행하였다. 구체적으로, 학습된 모델(\texttt{in\_features}$=64$, \texttt{hidden}$=64$, seed$=42$)에서 밴드 $t\in[100,200]$의 $t\in[165,173.5]$ 구간 다섯 영점이 평균 $|F_2|\approx 0.036$을 보였으며, 이는 밴드 평균 $0.020$의 약 $1.8\times$에 해당한다. 연쇄 실험은 네 가지 후보 해석을 기각하고 다섯 번째를 식별하였다.

\smallskip
\noindent\textbf{기각된 가설.}
\begin{enumerate}
  \item[(E1)] \emph{Lehmer형 근접쌍.} $51$개 학습 밴드에 걸친 $15{,}422$개 영점 전체에 대한 $|F_2(t_k)|$와 Riemann--von Mangoldt 정규화 간격 $g_{\min}^{\mathrm{norm}}(t_k)$의 다변량 회귀, 그리고 별도의 \texttt{mpmath} 기반 고위 $t$ 스캔 모두 $|r|<0.03$을 주었다. 군집 밴드에서 가장 조밀한 정규화 간격을 가진 영점 $t=121.37$ ($g_{\min}^{\mathrm{norm}}=0.027$)은 오히려 밴드 내 \emph{가장 쉬운} 영점이다.
  \item[(E2)] \emph{다중 basin 시드 의존성.} 동일 하이퍼파라미터로 5-seed 앙상블 학습을 수행한 결과 영점 전체에 걸쳐 $\mathrm{corr}(\mathrm{mean},\mathrm{std})=+0.70$을 얻었으나, 어려운 영점에서 앙상블 이득은 없었다 ($|F_2|^{\mathrm{ensemble}}/|F_2|^{\mathrm{seed-mean}}\approx 1.00$). 이는 군집을 다수-국소최소 현상으로 해석하는 것을 배제한다.
  \item[(E3)] \emph{은닉층 용량.} \texttt{in\_features}$=64$ 고정에서 $\texttt{hidden}\in\{32,64,128\}$ 스케일링은 \emph{단조 역방향} 경향을 주었다: 전역 $\overline{|F_2|}$가 $0.019\to 0.022\to 0.030$으로 증가하여, 더 넓은 은닉층이 군집에 도움이 되지 않았다.
  \item[(E4)] \emph{학습 표본 밀도.} \texttt{num\_points}를 $500$에서 $1500$으로 증가시키면 전역 $\overline{|F_2|}$가 $32\%$, 군집 최댓값 $|F_2|$가 $57\%$ 개선되었다; $t=173.41$에서는 $65\%$ 개선. 이는 처음에는 ``샘플링 해상도'' 치료제로 읽혔으나, 결국 아래에서 식별되는 기저 대역폭 병목의 대리지표임이 드러났다.
\end{enumerate}

\smallskip
\noindent\textbf{식별된 원인: 입력 기저 대역폭.} 입력 임베딩은 스칼라 $t$를 푸리에 특성 벡터로 승격시킨다:
\[
  t\;\longmapsto\;\Bigl(\tfrac{t-\mu}{\sigma},\;
  \sin(\omega_1 t),\cos(\omega_1 t),\dots,
  \sin(\omega_K t),\cos(\omega_K t)\Bigr),
\]
여기서 $\omega_k=2\pi k/(t_{\max}-t_{\min})$이고 $K=(\texttt{in\_features}-1)/2$이다. 밴드 $[100,200]$에서 \texttt{in\_features}$=64$이면 $K=31$이고 대역폭 천장이 고정된다:
\[
  \omega_{\max}^{(64)}
  \;=\; \frac{2\pi\cdot 31}{100}\;\approx\; 1.948\;
  \mathrm{rad}/\mathrm{unit}\,t,
  \qquad \text{최소 주기 }\approx 3.23\text{ (단위 }t\text{)}.
\]
군집 내부에서 $t\approx 169.095$와 $t\approx 169.912$의 두 영점은 간격 $\Delta t\approx 0.82$의 쌍을 이룬다. 이 간격 사이에서 두 번의 부호 변화를 표현하려면 기저가 다음을 포함해야 한다:
\begin{equation}
  \omega\;\gtrsim\; \frac{\pi}{\Delta t}\;\approx\; 3.83\;
  \mathrm{rad}/\mathrm{unit}\,t,
  \label{eq:required_omega}
\end{equation}
이는 \texttt{in\_features}$=64$ 천장의 약 $1.97\times$이다. 이 관찰을 반증 가능한 명제로 정식화한다.

\begin{proposition}[입력 기저 대역폭 필요조건]
\label{prop:basis_bandwidth}
폭 $W=t_{\max}-t_{\min}$의 밴드 $[t_{\min},t_{\max}]$ 위에서, $\omega_k=2\pi k/W$인 푸리에 특성 사상 $t\mapsto(\sin\omega_k t,\cos\omega_k t)_{k=1}^{K}$에 대해 \emph{선형}인 함수들의 클래스를 $\mathcal{F}_K^{\mathrm{lin}}$라 하자. 임의의 두 종대 $t_a,t_b\in[t_{\min},t_{\max}]$에 대해 $\Delta t:=|t_a-t_b|>0$이라 하자. 다음 샘플링 이론적 한계가 성립한다: 만약
\begin{equation}
  K \;<\; \frac{W}{2\,\Delta t},
  \label{eq:basis_bandwidth_condition}
\end{equation}
이면, $f(t_a)=f(t_b)=0$이고 $f'(t_a)\cdot f'(t_b)<0$인 비자명한 $f\in\mathcal{F}_K^{\mathrm{lin}}$가 존재할 수 없다. 특히 간격 $\Delta t$에서 $f\in\mathcal{F}_K^{\mathrm{lin}}$의 두 부호 변화 영점은 $K\geq W/(2\Delta t)$를 강제한다.
\end{proposition}

\begin{proof}[증명 개요]
$\mathcal{F}_K^{\mathrm{lin}}$의 함수는 $[0,W]$ 위의 차수 $\leq K$의 실수 삼각다항식이다. Nyquist--Shannon 샘플링 정리를 그 $2K$차원 푸리에 기저에 적용하면, 연속되는 부호 변화 사이의 최소 주기가 $W/(2K)$로 아래로 유계이다. 따라서 간격 $\Delta t<W/(2K)$에서의 두 부호 변화는 $\mathcal{F}_K^{\mathrm{lin}}$에서 실현 불가능하며, 이가 \eqref{eq:basis_bandwidth_condition}이다.
\end{proof}

\begin{remark}[비선형 확장과 intermodulation 간극]
\label{rem:intermodulation_gap}
문제의 네트워크는 $\mathcal{F}_K^{\mathrm{lin}}$에 속하지 \emph{않는다}: 세 겹의 ResonantBlock 층이 푸리에 특성에 곱셈, 리지 비선형성, 게이지 ODE 적분을 합성한다. 원리상 $\sin(\omega_j t)\cdot\sin(\omega_k t)=\tfrac{1}{2}[\cos((\omega_j-\omega_k)t)-\cos((\omega_j+\omega_k)t)]$ 형태의 곱과 고차 intermodulation 캐스케이드는 $\omega^{\mathrm{eff}}_{\max}=L\cdot\omega_{\max}$까지의 성분을 만들 수 있으며, 여기서 $L$은 네트워크 깊이이다. 이는 형식적으로 대역폭 천장을 들어올려야 한다. 그러나 우리 실험 자료(\cref{tab:basis_scaling})는 선형 기저 천장 \eqref{eq:basis_bandwidth_condition}에서 정확히 날카로운 실패를 보이며, 비선형적으로 들어올려진 천장 $L\cdot\omega_{\max}$에서는 실패하지 \emph{않는다}.

이 간극은 신경망의 \emph{spectral bias} 문헌과 일관된다: 유한한 학습 예산에서 MLP는 최저주파 성분을 먼저 학습하는 경향을 가지며, intermodulation 캐스케이드는 주파수에 따라 발산하는 학습 시간에서만 활성화된다 \cite{Rahaman2019SpectralBias,BasriGeifman2020FrequencyBias}. 이는 Neural Tangent Kernel 고유값 감쇠로부터 유도 가능하다 \cite{Jacot2018NTK}. 이러한 유한-예산 체제 안에서, 학습된 네트워크의 \emph{유효} 대역폭은 형식적 곱셈 상승 $L\cdot\omega_{\max}$가 아니라 선형 기저 천장의 상수배에 의해 위로 유계이다. ResonantBlock 아키텍처에 대한 ``$\omega^{\mathrm{eff}}_{\max}(E)\leq c(E)\cdot\omega^{\mathrm{basis}}_{\max}$, $c(E)=O(1)$ ($E\to E_{\max}$ 유한 학습 예산)'' 형태의 엄밀한 진술은 여기서 유도하지 않으며 열린 문제로 남긴다; 실험이 보이는 것은 천장이 \emph{경험적으로} $L\cdot\omega_{\max}$가 아니라 명제 \ref{prop:basis_bandwidth}처럼 거동한다는 것이다.
\end{remark}

$[100,200]$ 밴드와 $t\approx 169$ 쌍의 경우 \eqref{eq:basis_bandwidth_condition}은 $K\geq 61$, 즉 $\texttt{in\_features}\geq 123$을 준다.

\smallskip
\noindent\textbf{Riemann--von Mangoldt 대입.} 조건 \eqref{eq:basis_bandwidth_condition}은 Riemann--von Mangoldt 국소 영점 밀도 $N'(t)\sim (2\pi)^{-1}\log(t/2\pi)$과 결합될 때 해석적 수론적 내용을 얻는다. 이로부터 평균 영점 간격은 다음과 같다:
\begin{equation}
  \langle\Delta t\rangle(t)\;\sim\;\frac{2\pi}{\log(t/2\pi)}.
\end{equation}
Lehmer 압축 인자 $c_{\mathrm{L}}\in(0,1]$을 도입하여 높이 $t$에서의 폭 $W$의 전형적 밴드의 최소 간격이 $\Delta t_{\min}\gtrsim c_{\mathrm{L}}\cdot\langle\Delta t\rangle(t)$를 만족한다고 가정하자. 이를 \eqref{eq:basis_bandwidth_condition}에 대입하면 \emph{높이 의존 대역폭 요구조건}을 얻는다:
\begin{equation}
  K(t,W)\;\geq\;
  \frac{W\,\log(t/2\pi)}{4\pi\,c_{\mathrm{L}}},
  \label{eq:basis_bandwidth_rvm}
\end{equation}
동치로 $\texttt{in\_features}\geq (W\log(t/2\pi))/(2\pi\,c_{\mathrm{L}})+1$. $W=100$, $t\approx 150$, $c_{\mathrm{L}}\approx 1/3$ (Lehmer형 군집)인 경우 $K\gtrsim 82$, $\texttt{in\_features}\gtrsim 165$을 주며, 이는 $\texttt{in\_features}\in[64,128]$의 경험적 천이 구간을 포함한다. 식 \eqref{eq:basis_bandwidth_rvm}은 \emph{아키텍처-해석} 다리이다: $\xi$-영점 탐지를 위해 요구되는 푸리에 특성 용량은 높이에 \emph{로그적}으로 증가하며, 이는 $N'(t)$ 법칙과 일치한다. 특히 \cref{tab:f2_scale}에서 $\overline{|F_2|}$가 높이에 따라 겉보기로 개선되는 현상은 원래 스윕이 $\texttt{num\_points}=\max(200,10\,n_{\mathrm{zeros}})$ 자동 스케일링으로 실행되어 무차원 비율 $K/(W\log t)$를 보존했기 때문에 \eqref{eq:basis_bandwidth_rvm}을 암묵적으로 만족하는 고정 $\texttt{in\_features}$와 \emph{양립}한다; 이는 높이 스케일링 표를 $\xi$의 어떤 해석적 성질이 아니라 아키텍처 조건에 연결한다.

\smallskip
\noindent\textbf{기저 스케일링에 의한 확인.} \texttt{hidden}$=64$, seed$=42$, \texttt{num\_points}$=500$, 학습 예산 $100$ epoch 고정하에 $\texttt{in\_features}\in\{32,64,128,256\}$을 스케일링한 결과:

\begin{table}[ht]
\centering
\caption{고정 학습 예산에서의 군집 vs. 비군집 $\overline{|F_2|}$. 5-시드 앙상블 행 (시드 $42,7,123,2026,314$)이 단일 실행 주장을 대체한다; 군집 개선은 평균 $25\%$이나 시드 분산이 높다.}
\label{tab:basis_scaling}
\small
\begin{tabular}{rrrrrr}
\toprule
\texttt{in\_feat} & $\omega_{\max}$ & 전역 &
  군집 & 대조 & 비율 \\
\midrule
32   & 0.942 & 0.0140 & 0.0140 & 0.0140 & $1.00\times$ \\
64   & 1.948 & 0.0218 & 0.0361 & 0.0202 & $1.79\times$ \\
128 (시드 42)  & 3.958 & 0.0164 &
  0.0150 & 0.0166 & $0.90\times$ \\
\textbf{128 (5-시드)} & \textbf{3.958} & \textbf{0.0207} &
  \textbf{0.0269}$\pm$\textbf{0.012} & \textbf{0.0200} & $\mathbf{1.42\pm0.73\times}$ \\
256  & 7.980 & 0.0214 & 0.0581 & 0.0173 & $3.36\times$ \\
\bottomrule
\end{tabular}
\end{table}

\noindent 두 가지 특징이 두드러진다. 첫째, \texttt{in\_features}$=64$에서의 군집 초과는 \texttt{in\_features}$=128$에서 최선의 시드에서는 \emph{사라지지만}, 5-시드 앙상블 (시드 $42, 7, 123, 2026, 314$)은 상당한 변동성을 드러낸다: 군집 대 비군집 비율의 평균은 $1.42\pm0.73$ (범위 $0.61$--$2.44\times$)이며, $3/5$ 시드가 비율${}<1.2$을 달성하고 $2/5$에서는 군집이 \emph{지속}되거나 악화된다. 군집 영점별 변동계수(CV)는 시드 간 $46\%$에서 $106\%$에 이른다. 따라서 대역폭 조건 \eqref{eq:basis_bandwidth_condition}은 군집 해소의 \emph{필요}조건이지만 \emph{충분}조건은 아니다: Nyquist 천장이 통과된 후에도 학습 역학(시드 의존 초기화)이 상당한 영향을 유지한다. 앙상블 평균에서 군집 $\overline{|F_2|}$는 \texttt{in\_features}$=64$ 기준선 대비 $25\%$ 개선된다 ($0.0269$ vs.\ $0.0361$), 이는 대역폭 메커니즘을 모집단 수준에서 확인하면서 단일 실행의 완전 해소 주장을 배제한다. 둘째, \texttt{in\_features}$=256$은 $100$ epoch에서 수렴에 실패하며 \emph{새로운} 무관한 어려움 군집을 보인다.

\smallskip
\noindent\textbf{과매개변수화 잡음의 분리.} 학습을 $300$ epoch까지 확장하고 $50$-epoch 간격의 best-validation 체크포인트를 취하면 대역폭 효과를 최적화 효과로부터 분리할 수 있다. \texttt{in\_features}$=128$의 경우 best validation이 epoch $143$에서 $0.0550$으로 안정화되고 이후 평탄하며, 군집 비율은 $0.75\times$로 개선되고 $\overline{|F_2|}^{\mathrm{cluster}}$가 $0.0106$까지 낮아진다 ($\overline{|F_2|}^{\mathrm{other}}=0.0142$ 대비). \texttt{in\_features}$=256$의 경우 best validation이 $300$ epoch 내내 $0.224$에 머물러 \texttt{in\_features}$=128$ 최솟값보다 약 $4\times$ 나쁘며, 전역 $\overline{|F_2|}$는 $0.0214$에서 $0.0099$로 개선되지만 군집 비율은 $3.87\times$로 유지된다. 따라서 \texttt{in\_features}$=256$의 실패는 독립적인 과매개변수화 병리 현상 (확장된 파라미터 공간에 대해 Adam+gradient-clip 학습 예산이 불충분)이며 대역폭 문제가 아니다; 이는 $\texttt{in\_features}\in\{32,64,128\}$에서의 대역폭 해석을 훼손하지 않는다.

\smallskip
\noindent\textbf{2026-04-08 재해석에 대한 귀결.} 본 실험은 \cref{sec:f2_residual_open}의 두 번째 열린 질문을 닫는다: $[165,173.5]$의 군집이라는 구체적이고 \emph{모델 의존적}인 개별 영점 어려움이, 영점의 고유 성질이 아니라 임베딩 $\{\omega_k\}_{k=1}^{K}$의 성질임을 드러냈다. 따라서 개별 영점 $|F_2|$ 값을 임계선 위의 관측량으로 읽으려는 시도는, 먼저 해당 밴드의 가장 작은 영점 쌍에 대해 기저 대역폭 조건 \eqref{eq:basis_bandwidth_condition}이 충족되는지를 검증해야 한다. \cref{tab:f2_scale}의 평균 스케일 주장은 영향을 받지 않는다. 왜냐하면 그것은 자동 스케일링 규칙 $\texttt{num\_points}=\max(200,10\cdot n_{\mathrm{zeros}})$을 통해 (암묵적으로) 조건이 충족되거나 대략 충족되는 밴드들에 대한 평균이기 때문이다.

\subsection{여전히 열려 있는 것}
\label{sec:f2_residual_open}

2026-04-10 후속 이후 여전히 열려 있는 질문은 다음과 같다:
\begin{enumerate}
  \item \textbf{높이 스케일링 법칙
        \textcolor{ForestGreen}{[2026-04-12 부분 해결]}.}
        블라인드 외삽 재현율은 $t{\sim}200$에서 $t{\sim}1100$까지
        $97.5{-}98.8\%$로 안정적이며 (3시드, $K\in\{128,256\}$),
        정밀도는 고높이에서 \emph{개선}된다. 이는 대규모 성능 저하를
        배제한다. 세부 질문---$\overline{|F_2|}$이 진정한 멱법칙을
        따르는지, 그 지수가 $N(t)$의 $\log\log t$ 보정과
        일치하는지---은 분리된 시드 가족에 대한 오차 막대가 있는
        연속 스윕을 여전히 요구한다.
  \item \textbf{대역폭-청정 개별 영점 관측량.} \texttt{in\_features}$=128$의 5-시드 앙상블은 \eqref{eq:basis_bandwidth_condition} 충족 후에도 개별 영점 $|F_2|$ 값의 변동계수가 시드 간 $46$--$106\%$에 달함을 확인한다. 따라서 개별 영점 값은 순진한 의미에서 \emph{재현 불가}이다. \emph{앙상블 평균} 군집 $\overline{|F_2|}$는 대역폭 부족 기준선 대비 $25\%$ 개선되어, 모집단 수준의 메커니즘은 확인된다. 더 정교한 불변량 (예: 부트스트랩 앙상블에서 $|F_2(t_k)|$의 \emph{최솟값}을 ``달성 가능 잔차''로 해석)을 추출할 수 있는지는 여전히 열려 있다.
  \item \textbf{과매개변수화 실패 양상.} \texttt{in\_features}$=256$이 $300$ epoch 내에 수렴에 실패한 것 (best validation $0.224$, \texttt{in\_features}$=128$의 $0.055$ 대비)은 동일한 $t\in[165,173.5]$에서 대역폭 기원일 수 없는 ``새로운'' 군집을 보인다. 이것이 (a) Adam+gradient-clip 학습 스케줄 잡음, (b) 다른 기전을 통해 동일 군집을 부활시키는 진정한 과매개변수화 병리, 또는 (c) 시드 우연인지는 미해결이다.
\end{enumerate}


% ============================================================================
\section*{수치 증거 요약}
\label{sec:summary}
\addcontentsline{toc}{section}{수치 증거 요약}
% ============================================================================

표~\ref{tab:summary}는 여덟 가지 검증 축에 걸쳐 얻은 62개의 수치 결과를
종합한다. 상태 코드: \textbf{E} = 확립(Established),
\textbf{C} = 조건부 양성(Conditionally positive),
\textbf{N} = 음성(Negative, 완전 공개).

\begin{center}
\small
\begin{longtable}{@{}r l l p{3.8cm} l@{}}
\caption{여덟 검증 축에 걸쳐 62개 수치 결과 요약.\label{tab:summary}}\\
\toprule
\# & 결과명 & 상태 & 핵심 지표 & 축 \\
\midrule
\endfirsthead
\multicolumn{5}{c}{\small\itshape (표~\ref{tab:summary} 계속)}\\
\toprule
\# & 결과명 & 상태 & 핵심 지표 & 축 \\
\midrule
\endhead
\midrule
\multicolumn{5}{r}{\small\itshape 다음 페이지 계속}\\
\endfoot
\bottomrule
\endlastfoot
1  & 쿠라모토 동기화           & \textbf{E} & $r=0.9994\pm0.0001$              & 국소 \\
2  & $S^1$ 측지 손실           & \textbf{E} & 탐지 이득 $4\times$              & 국소 \\
3  & 고높이 스케일링           & \textbf{E} & recall $97.5\%+$ ($t=1100$까지) & 국소 \\
4  & $H$-스케일링 지수         & \textbf{C} & $\alpha\approx{-0.4}$            & 국소 \\
5  & FP 해부                   & \textbf{E} & $85.1\%$ 유사-영점              & 국소 \\
6  & Conj.\,3 $\kappa$-분리    & \textbf{E} & FP max $<$ TP min, $300\times$  & 국소 \\
7  & $\sigma=\tfrac{1}{2}$ 유일성 & \textbf{E} & 에너지 비율 $385\times$       & 국소 \\
8  & 블라인드 예측             & \textbf{E} & F1$=1.000$, 7/7                  & 국소 \\
9  & $S^1$ 고높이              & \textbf{C} & $\Delta\text{F1}=+0.1$\%p        & 국소 \\
\midrule
10 & 디리클레 다발             & \textbf{E} & 3 지표 $\times$ 4 성질          & 디리클레 \\
11 & 디리클레 교차 비교        & \textbf{E} & 4 함수 $\times$ 5 성질          & 디리클레 \\
12 & 디리클레 블라인드         & \textbf{E} & F1$=1.000$, 21/21               & 디리클레 \\
13 & 고-conductor              & \textbf{E} & 5/5 성질                        & 디리클레 \\
14 & 실수 지표 에너지          & \textbf{E} & $\chi_4{=}132.5\times$, $\chi_8{=}194.3\times$ & 디리클레 \\
15 & off-critical 해부         & \textbf{E} & Conj.\,1 디리클레 일반화        & 디리클레 \\
\midrule
16 & Lehmer 쌍 곡률            & \textbf{C} & $\rho=0.835$                     & 전역 \\
17 & Gram 점 곡률              & \textbf{N} & $\kappa$-향상 없음               & 음성 \\
18 & 중간점 비국소성           & \textbf{E} & $\rho={-0.654}$ ($p<10^{-29}$)  & 전역 \\
19 & 디리클레 중간점 비국소    & \textbf{E} & 3/3 $|\rho|>0.55$               & 전역 \\
20 & GUE/Poisson 대조          & \textbf{E} & GUE $\rho(b)\approx 0$          & 전역 \\
21 & 비국소 도달 감쇠          & \textbf{C} & $k=1$ 집중                      & 전역 \\
\midrule
22 & 홀로노미 / Chern 수       & \textbf{E} & $\zeta$ 87/87, $\chi_3$ 4/4, $\chi_4$ 4/4 & 위상적 \\
23 & 짝수 Chern + 방법론       & \textbf{E} & $\chi_5$ 4/4, $\chi_8$ 2/2      & 위상적 \\
24 & Gauss--Bonnet 면적분      & \textbf{E} & $\zeta$ 6/6, $\chi_5$ 2/2, 오차$=0$ & 위상적 \\
25 & off-critical $\sigma$-국소화 & \textbf{E} & $\zeta$ 7/7, FWHM $<0.004$  & 위상적 \\
26 & 삼중 위상적 정합성        & \textbf{E} & $4{\times}3$ 행렬, $N=100$      & 위상적 \\
\midrule
27 & 합성 함수 음성 대조군     & \textbf{E} & 11/11, 독립 재현 3회            & 음성 \\
28 & 수 분산 (RMT)             & \textbf{C} & Var$_\text{plaq}{\approx}$Var$_\text{dir}$, Berry sub-GUE & 전역 \\
\midrule
29 & Epstein 제타 확장          & \textbf{C} & $\lambda{=}1$ vs $\lambda{=}5$, $\kappa$ 편차 ${<}1\%$  & 확장 \\
\midrule
30 & 고 $t$ 견강성               & \textbf{E} & $\kappa$ CV${=}0.2\%$, $t{=}14$--$10^4$ & 견강성 \\
\midrule
31 & $\Delta\kappa$ 차선도       & \textbf{C} & 기울기${=}1.59$, Hadamard 오차${=}0.00\%$ & 산술적 \\
32 & NNS--$\Delta\kappa$ Lorentzian$^2$ & \textbf{C} & $R^2_\text{2-var}{=}0.919$, $r{=}0.953$ & 산술적 \\
33 & $\delta$-sweep Lorentzian$^2$ 보편성 & \textbf{C} & $5/5$ $R^2{>}0.88$; 통합 $R^2{=}0.449$ (음성) & 산술적 \\
\midrule
34 & 간격비 vs.\ GUE (Wigner surmise)    & \textbf{N} & KS $p{=}4.7{\times}10^{-4}$; $D_\text{GUE}/D_\text{Poisson}{=}0.079$ & RMT 경계 \\
35 & $\kappa$ 분포 vs.\ GUE 예측         & \textbf{N} & KS $D{=}0.361$; 스케일 $7.6{\times}$; Q--Q $r{=}0.97$ & RMT 경계 \\
36 & $\kappa$ 자기상관 vs.\ GUE 두점 상관 & \textbf{N} & $\tau_\text{unf}{=}0.45{\neq}1.0$; $R^2{=}0.48$; 지연-1 $C{=}{-}0.08$ & RMT 경계 \\
\midrule
37 & 디리클레 $\chi\bmod 7$ (고 conductor) & \textbf{E} & 136 영점; 모노 136/136; 블라인드 7/7; conductor 효과 B/E & 디리클레 \\
38 & 교차비교: 5개 GL(1) $L$-함수 & \textbf{E} & $\kappa$ 단조 감소 Spearman $r_s{=}{-}1.000$; $\sigma$-유일성 5/5; 예비 5점 경향 & 디리클레 \\
\midrule
39 & 타원곡선 11a1 (GL(2)) & \textbf{E} & 모노 17/17; $\kappa$ $972\times$; 블라인드 15/15; $\sigma$-유일성 N/A & GL(2) \\
40 & 타원곡선 37a1 (GL(2), rank 1) & \textbf{E} & 모노 23/23; $\kappa$ $2684\times$; 블라인드 20/20; $\sigma$-유일성 N/A & GL(2) \\
41 & 라마누잔 $\Delta$ (GL(2), weight 12) & \textbf{E} & 모노 8/8; $\kappa$ $1396\times$; 블라인드 7/7; $\sigma$-유일성 N/A & GL(2) \\
42 & $\sigma$-유일성 메커니즘 (GL(1) vs GL(2)) & \textbf{C} & $R_{\mathrm{GL}(1)}{\approx}0.4$, $R_{\Delta}{=}1.000$ (항등식); $\chi_3$ 경험 $0.38{\approx}0.37$ & GL(2) \\
\midrule
43 & FP 모노드로미 해부 (추측~3) & \textbf{E} & TP mono$/\pi{=}2.000{\pm}0$; FP mono$/\pi{=}0.000$; $p{=}6.2{\times}10^{-6}$; 이중기준 정밀도 $+44\%$p & 위상학적 \\
44 & GL(2) FP 모노드로미 (교차 순위) & \textbf{E} & 34 TP mono$/\pi{=}2.000{\pm}0$; 60 FP mono$/\pi{=}0.000$; $p{=}9.7{\times}10^{-17}$; 3/3 $L$-함수 일관 & GL(2) \\
45 & GL(2) 블라인드 영점 예측 (11a1) & \textbf{E} & 16/17 TP ($P{=}R{=}F_1{=}0.941$); 위치 오차 최대${=}0.051$; $\kappa$-피크 블라인드; FP $t{\approx}28.4$ (이중피크 잔상) & GL(2) \\
46 & GL(2) 블라인드 영점 예측 (37a1, rank 1) & \textbf{E} & 22/23 TP ($P{=}1.000$, $R{=}0.957$, $F_1{=}0.978$); FP${=}0$; mono$/\pi{=}4.0$ (권선수${=}2$, 윤곽 내 2영점); $\varepsilon{=}{-1}$, rank~1 & GL(2) \\
47 & GL(2) 블라인드 영점 예측 ($\Delta$, weight 12) & \textbf{E} & 8/8 TP ($P{=}R{=}F_1{=}1.000$); FP${=}0$; $\sigma_c{=}6$; 위치 오차 최대${=}0.108$; $\kappa$-비율 $322.8\times$; weight 보편성 & GL(2) \\
48 & 디리클레 블라인드 영점 예측 ($\chi\bmod 7$) & \textbf{C} & 13/13 TP ($R{=}1.000$, $P{=}0.722$, $F_1{=}0.839$); FP 5개 ($t{<}17$, off-line 포획); $t{\geq}17$: $P{=}R{=}F_1{=}1.000$; $\kappa$-비율 $24.0\times$ & 디리클레 \\
49 & GL(2) $t$-스케일링 안정성 (11a1, $t{\approx}51$--$100$) & \textbf{C} & 18개 영점 ($t{\approx}48$--$104$, 전부 $\sigma{=}1.000$); $\kappa$ 안정 $1118$--$1191$ (${\leq}7\%$); peak$_\sigma{=}1.000$; mono$/\pi{=}2.0$ (3구간 중 2개); $t{\approx}147$: 수치 정밀도 장벽 (이론적 한계 아님) & GL(2) \\
\midrule
50 & GL(3) sym$^2$(11a1) 4성질 & \textbf{E} & $\kappa$ $283\times$; mono 12/12; FE${}=0.0$; $\sigma$-유일성 FAIL & GL(3) \\
51 & GL(3) sym$^2$(37a1) 4성질 & \textbf{E} & $\kappa$ $222\times$; mono 15/15; FE${}=0.0$; $\sigma$-유일성 FAIL & GL(3) \\
52 & GL(3) 블라인드 영점 예측 & \textbf{E} & $R{=}1.000$, $P{=}1.000$, $F_1{=}1.000$ (교정); 21개 신규 영점 & GL(3) \\
53 & GL(3) $\sigma$-유일성 5단계 역보정 & \textbf{N} & 5/5 FAIL; 원인${}={}$디리클레 급수 수렴 & GL(3) 경계 \\
54 & GL(3) FP 검증 (AFE) & \textbf{E} & 21/21 영점; 중앙값 $|\Lambda|{=}2.5{\times}10^{-27}$ & GL(3) \\
55 & GL(3) $\kappa$-도체 스케일링 (6곡선) & \textbf{E} & 스케일링: $\alpha{=}0.003$, $R^2{=}0.0002$ (음성); $\kappa_{\mathrm{near}}{\approx}1125$ (CV${=}0.15\%$, 5곡선); 도체 독립 보편상수 & GL(3) \\
\midrule
56 & Maass 형식 $L$-함수 (홀수, $R{=}9.53$) & \textbf{E} & FE${}=0.0$; 18영점; $\kappa_{\mathrm{near}}{=}1114$ (CV${=}0.1\%$, 8 TP); 모노 8/8; $\sigma$-유일성 3/3 PASS; $\kappa_{\mathrm{near}}{\neq}$GL(3) & GL(2) Maass \\
57 & Maass 형식 $L$-함수 (짝수, $R{=}13.78$) & \textbf{E} & FE${}=0.0$; 24영점; $\kappa_{\mathrm{near}}{=}1115$ (CV${=}0.2\%$, 12 TP); 모노 12/12; $\sigma$-유일성 24/24 PASS; $\kappa_{\mathrm{near}}(2) \approx 1115$ 확정 & GL(2) Maass \\
58 & 라마누잔 $\Delta$ $L$-함수 ($N{=}1$, $w{=}12$) 4성질 & \textbf{E} & FE${}=0.0$; 23영점; $\kappa_{\mathrm{near}}{=}1114$ (CV${=}0.1\%$, 12 TP); 모노 12/12; $\sigma$-유일성 23/23 PASS; $N{=}1 \Leftrightarrow$ PASS 확정 (경계 B-10 해결) & GL(2) $\Delta$ \\
59 & GL(1) $\zeta(s)$ 4성질 검증 & \textbf{E} & FE rel${=}1.66\times10^{-50}$; 13영점 ($t{\in}[10,60]$); $\kappa_{\mathrm{near}}{=}1112.32$ (CV${=}0.1\%$, 13 TP); 모노 13/13 ($\mathrm{mono}/\pi{=}2.0000$); $\sigma$-유일성 13/13 PASS ($\kappa(0.5)/\kappa(0.45){\in}[10^{26},10^{27}]$); $\kappa_{\mathrm{near}}(1){=}1112.32$ 앵커; $\kappa$ 비율 $2200.7{\times}$ (최고기록) & GL(1) $\zeta$ \\
60 & $\kappa$-$\delta$ 스케일링 법칙 검증 & \textbf{E} & 13영점, 5오프셋 $\delta{\in}\{0.01,\ldots,0.10\}$, 65점; $\kappa{\cdot}\delta^2{\in}[1.00012,1.01206]$ (전 CV${<}1\%$); $O(1/\delta)$ 계수 ${=}0$ 확인; $A(t_0){=}\kappa{-}1/\delta^2$ $\delta$-독립; B-12 해소: $1/\delta^2$ 보편, 차수 차이는 $A(t_0,\mu)$; 명제~\ref{prop:kappa_exact_expansion} & GL(1) $\zeta$ \\
\midrule
61 & 3-차수 $\kappa$-$\delta$ 스케일링 $+$ $A(t_0)$ 차수 비교 & \textbf{E} &
   GL(2) Maass ($R{\approx}13.78$): 12영점 $\times$ 5$\delta$ $=$ 60점, $\kappa{\cdot}\delta^2{\in}[1.00,1.04]$;
   GL(3) sym$^2$(11a1): 12영점 $\times$ 5$\delta$ $=$ 60점, $\kappa{\cdot}\delta^2{\in}[1.00,1.12]$;
   평균$(A)$: GL(1)${=}1.30$, GL(2)${=}3.93$, GL(3)${=}12.79$ (단조 증가);
   $A$ $\delta$-독립: GL(2) spread ${<}0.21\%$, GL(3) ${<}0.70\%$;
   스케일링 법칙 $\kappa{=}1/\delta^2{+}A(t_0){+}O(\delta^2)$ 차수 1--3 전반 성립
   (비고~\ref{rem:kappa_degree_comparison} 참조);
   $A(d)$ 함수형 미결정 (3점); 개방 문제: B-18
   & GL(1--3) \\
\midrule
62 & 디리클레 $\chi$ $\kappa$-$\delta$ 스케일링 $+$ $A(t_0)$ 도체 의존성 & \textbf{E} &
   $\chi\bmod 3,4,5$ (홀수, $\mu{=}1$): 66영점 $\times$ 5$\delta$ $=$ 330점;
   $\kappa{\cdot}\delta^2$ 전점 통과 (최대 1.028);
   $A(t_0)$ $\delta$-독립 (spread ${<}0.1\%$);
   $A$: $\zeta{=}1.27$, $\chi_3{=}2.25$, $\chi_4{=}2.77$, $\chi_5{=}3.14$ (도체 단조 증가);
   동일-$\delta$ $\kappa_{\mathrm{near}}$ 차이 ${<}0.2\%$ $\Rightarrow$ 차수-전용 확인 (B-09);
   비고~\ref{rem:dirichlet_conductor_A} 참조
   & GL(1) 디리클레 \\
\end{longtable}
\end{center}

\noindent\textit{상태}: \textbf{E} = 확립(시드/축 전반 확인);
\textbf{C} = 조건부 양성(유의하나 단서 있음);
\textbf{N} = 음성(완전 공개; \cref{app:gram_negative,obs:synthetic_control,obs:rmt_boundaries} 참조).

% ============================================================================
\section{토론}
\label{sec:discussion}
% ============================================================================

\subsection{대응이 의미하는 것}

7개의 대응은 수학적 문법 수준에서의 \textit{구조적 동형}을 입증한다: 주기적 위상 벌칙, 경성 오차 경계, 진폭--위상 분리, 자가 교정 수렴의 동일한 패턴이 구체적 공학 시스템(QROP-Net)과 추상적 해석 프레임워크(임계선 위의 공명 이론) 모두에서 나타난다.

이는 다음 사이의 보다 깊은 연결을 시사하지만---\textbf{증명하지는 않는다}:
\begin{enumerate}
  \item 시클로토믹 다항식환의 대수적 구조 ($X^n + 1 = 0$, 격자 암호의 기반),
  \item 리만 $\xif$-함수의 스펙트럼 구조 (복소 해석적 위상에서의 소수 분포 인코딩),
  \item 결맞음 광학 전파의 물리적 구조 (분수 푸리에 변환으로서의 프레넬 회절 핵).
\end{enumerate}

통합의 실마리는 \textit{단위원}이다: $S^1$ 위의 다항식 근, $S^1$ 위의 광학 위상, 그리고 $S^1$ 위의 $\xif$-함수 값 (위상 성분). 위상 양자화---연속 신호를 이산 각도 위치로 강제---가 공유 연산이다.


\subsection{대응이 의미하지 않는 것}

본 연구의 경계에 대해 명시적이어야 한다:

\begin{enumerate}
  \item \textbf{이것은 리만 가설의 증명이 아니다.}
        구조적 병렬은 공유 수학적 문법을 식별하지만, 논리적 증명을 구성하지 않는다. RH는 특정 해석 함수의 영점에 대한 진술이며, 광학 시스템에 대한 어떤 구조적 유비도 엄밀한 복소 해석적 논증을 대체할 수 없다.

  \item \textbf{대응은 구조적이지 인과적이 아니다.}
        QROP-Net의 $\sin^2$ 손실과 게이트~A의 $\Delta v = m\pi$가 같은 함수 형태를 가진다는 사실이, 하나가 다른 것을 ``야기''하거나 ``설명''함을 의미하지 않는다. 이들은 같은 수학적 패턴의 독립적 실현이다.

  \item \textbf{다원적 프레임은 수치 도구이지 증명 기법이 아니다.}
        $F_2(t) = 0$이 알려진 영점 위치로 수렴하는 것은 수치 방법을 검증하지만, 영점을 찾는 것은 그것들이 모두 $\sigma = 1/2$ 위에 있음을 증명하는 것과 같지 않다.

  \item \textbf{산란 해석은 유비이다.}
        $\xif(s)$를 산란 진폭으로 취급하는 것은 올바른 예측(영점 위치, 위상 양자화)을 생성하는 생산적 유비이지만, $\xif$는 물리적 시스템의 산란 진폭이 문자 그대로 아니다 (S-행렬이 $\xif$와 같은 해밀토니안은 식별되지 않았다).
\end{enumerate}


\subsection{예측력과 향후 방향}
\label{sec:open}

이러한 주의에도 불구하고, 통합 프레임워크는 여러 검증 가능한 예측을 제시한다. 이를 계층별로 조직하며, \cref{sec:tiers}의 분류 원칙을 보존한다: 항목은 이미 확보된 증거의 구체성에 따라 묶인다.

\smallskip\noindent\textbf{\textcolor{ForestGreen}{\rule{0.9ex}{0.9ex}}\ 계층 1 후속작업 (즉시 실행 가능).}
\begin{itemize}
  \item \textbf{$\arg\xif(\tfrac{1}{2}+it)$에 대한 복소 벡터 읽기
        \textcolor{red}{[2026-04-12 중립]}.}
        고립 실험 (\cref{rem:isolation})은 잔존 결함을 영점에서의 매끄러운 네트워크가 $\arg\xif$의 $\pm\pi$ 불연속을 표현할 수 없는 것으로 국소화했다. 가장 깔끔한 아키텍처 해법은 읽기 헤드가 $\xif(\tfrac{1}{2}+it)$를 복소 벡터로 예측하고 $\arg$를 사후 계산하게 하는 것이다.
        \emph{2026-04-12 시험:} 프로젝션 전 은닉 상태 $Z\in\mathbb{C}^{64}$를
        $(\mathrm{Re}\,\xif,\mathrm{Im}\,\xif)$로 매핑하는 선형 읽기 헤드를
        보조 공명 손실($\Lres+\Lcurv+\Lpqo$, $\Ltgt=0$)과 함께 Re/Im MSE
        ($\lambda_{\mathrm{MSE}}=2$)로 공동 훈련하였다.
        3시드($t\in[100,200]$, $K=128$, 150 에포크)에서:
        복소 벡터 검증 손실은 ${\sim}4$배 개선
        ($0.014\pm0.005$ vs $0.069\pm0.013$)되었으나, $|F_2|$
        ratio는 통계적으로 무변화
        ($0.982\pm0.041$ vs $0.999\pm0.005$; $p>0.5$)이고 영점
        검출은 $0/463$에 머물렀다.
        읽기 헤드는 Re/Im 매핑을 잘 학습하지만 네트워크의 내부 공명
        구조($Z_{\mathrm{out}}$)를 변경하지 못한다. $\pm\pi$ 불연속은
        사후 읽기 헤드만으로는 해결할 수 없고 더 깊은 아키텍처 변경이
        필요하다.
  \item \textbf{매끄러운 영점 교차 타겟 \textcolor{ForestGreen}{[2026-04-10 확인]}.} 고립 설정 (\cref{rem:isolation})에서 두 가지 매끄러운 대리 타겟을 시험하였다: 범프 타겟 $\mathrm{sign}(\mathrm{Re}\,\xif)\cdot e^{-(\mathrm{Re}\,\xif)^2/\sigma^2}$와 정규화 타겟 $\mathrm{Re}\,\xif/\max|\mathrm{Re}\,\xif|$. 정규화 타겟은 영점/비영점 MSE 비율을 $18.2\times$ (계단 타겟)에서 $0.03\times$로 줄인다 --- $99.8\%$ 갭 축소 --- 이로써 $\pm\pi$ 불연속이 유일한 병목이었음을 확인한다 (\cref{rem:isolation}).
        \emph{전체 파이프라인 시험
        \textcolor{ForestGreen}{[2026-04-10 완료]}:}
        매끄러운 타겟을 전체 손실
        ($\Lres + \Lcurv + \Lpqo$ + MSE smooth)에 통합한 결과,
        전역 $|F_2|$가 $5.2\%$ \emph{악화}되었다 (3-시드 앙상블).
        이는 각도 타겟 $\Ltgt$가 $\Lres$, $\Lpqo$와 시너지를
        이루므로 전면 교체해서는 안 됨을 시사한다.
        매끄러운 타겟은 \emph{고립 진단} 도구로서만 유효하다.
  \item \textbf{블라인드 영점 외삽
        \textcolor{ForestGreen}{[2026-04-10 확인]}.}
        $t\in[100,150]$에서 훈련하고 미관측 구간 $[150,200]$에서
        국소 $|F_2|$ 극소를 예측한 결과 (역방향도 동일),
        5개 시드 및 $K\in\{64,128\}$ 설정에서 재현율
        $96{-}99\%$를 달성했다; 역방향 $K=128$에서
        $99.1\pm1.7\%$ 재현율. 정밀도는 가짜 격자 극소로 인해
        낮지만 ($\sim\!15{-}18\%$), 거의 완벽한 재현율은 $F_2$
        잔차가 훈련 세트 영점에 과적합하지 않고 미관측 높이
        범위로 일반화됨을 보여준다.
        \emph{사후 정밀도 필터링
        \textcolor{red}{[2026-04-12 중립]}:}
        블라인드 외삽 극소에 앙상블 합의 및 깊이/간격 필터를 적용하였다
        ($t\in[150,200]$, 5시드, $K=128$, 200 에포크, 2000점 밀집 격자).
        최고 필터(앙상블 $\ge 2$ 시드)는 정밀도 $32.7\%$,
        재현율 $63.0\%$, $F_1=0.430$을 달성하였다.
        정밀도는 ${\sim}15\%$ (원시)에서 $28{-}33\%$로 상승하였으나
        $1.5$배 유의 임계치에 도달하지 못하였다.
        깊이 필터는 너무 공격적이며(모든 후보 제거), 간격 필터는
        정밀도 보상 없이 재현율만 저하시킨다.
        낮은 정밀도는 사후 휴리스틱으로 교정 가능한 수준이 아니라
        $|F_2|$ 풍경 자체에 내재된 것으로 보인다.
        \emph{FP 해부
        \textcolor{ForestGreen}{[2026-04-13 확인]}:}
        5시드($t\in[14,50]$)에서 위양성의 $85.1\%$(57/67)가
        $|F_2|$ 풍경의 진짜 극소점($|F_2|''(t)>0$, $p<10^{-6}$)이다.
        유사-영점(pseudo-zero)은 학습된 잔차의 구조적 특성이지,
        잡음이 아님을 확인한다.
        FP 위치는 시드 간 비일관적(Jaccard${}=0.046$,
        $\ge\!5$시드 합의${}=0$)으로 $|F_2|$ 풍경의 세부 구조가
        모델 초기화에 의존한다.
        해석적 곡률 $\kappa=|\xif'/\xif|^2$는 TP와 FP를
        ${\sim}300$배 분리(TP 중앙값${}=270$, FP 중앙값${}=0.9$;
        갭: FP 최대${}=31$ $<$ TP 최소${}=76$)하여,
        $\kappa$가 이용 가능한 경우 이론적 정밀도 상한 $100\%$를
        확립한다. $\xif$가 $\sigma=1/2$ 위에서 실수이므로
        영점에서의 모노드로미는 자명하게 $\pm\pi$(부호 전환)이며,
        곡률 분리 이상의 새 정보를 제공하지 않는다.
  \item \textbf{$S^1$ 측지선 위상 예측
        \textcolor{ForestGreen}{[2026-04-12 확인]}.}
        $\arg\xif$를 스칼라로 예측하는 대신 (영점에서 $\pm\pi$
        분기점 발생) 읽기 헤드가
        $(\cos\theta,\sin\theta)\in S^1$을 예측하고 측지선 손실
        $L_{\mathrm{geo}}=1-\cos(\theta_{\mathrm{pred}}-
        \theta_{\mathrm{true}})$로 훈련한다.
        3시드($t\in[100,200]$, $K=128$, 150 에포크)에서:
        $S^1$ 검출이 463개 영점 중 $5.7$개에서 $13.7$개로
        $2.4$배 개선되었다. 영점에서의 각도 오차는 $0.215$ rad로
        비영점의 $0.155$ rad보다 $1.39$배 크며, $\pm\pi$ 불연속이
        부분적으로 우회되었음을 확인한다.
        $|F_2|$ ratio는 무변화($1.054\pm0.103$ vs $1.026\pm0.233$);
        검출 개선은 네트워크 공명 구조 변화가 아닌 읽기 헤드의
        영점 근방 위상 추적 개선에서 기인한다.
        \emph{후속 (완료):} $\Ltgt$를 $L_{\mathrm{geo}}$로 완전
        대체하여 전체 손실
        ($\Lres+\Lcurv+L_{\mathrm{geo}}+\Lpqo$)에 통합한 결과,
        읽기 헤드 단독 대비 $1.66$배 추가 개선:
        463개 영점 중 $22.7$개 검출 (읽기 헤드 $13.7$개,
        기준선 $5.7$개), 누적 $4$배 개선.
        $|F_2|$ ratio는 $0.913\pm0.105$로 하락하여
        (기준선 $1.026\pm0.233$), 영점이 비영점보다
        \emph{더 작은} $|F_2|$ 잔차를 보이는 질적 역전이
        발생하였다. 시드 간 변동(5, 40, 23개)이 크며,
        측지선 손실 풍경에 위상 정렬 품질이 상이한
        다수의 국소해가 존재함을 시사한다.
        \textcolor{ForestGreen}{[2026-04-12 확인, 임계치 메트릭]}
        %
        \emph{후속 (완료, 5시드):}
        $t\in[500,600]$에서 극소값 기반 평가를 5시드
        ($42, 7, 123, 314, 2024$)로 수행한 결과,
        $4$배 개선은 메트릭 조건부이다.
        극소값 검출 기준으로 $F_1$은 유사하다
        (BL $0.373\pm0.020$ vs $L_{\mathrm{geo}}$
        $0.361\pm0.012$; 차이 비유의).
        $L_{\mathrm{geo}}$의 유일한 일관된 효과는
        \textbf{재현율 하한 안정화}이다:
        BL 재현율 $0.956\pm0.045$ (범위 $[0.889,1.000]$),
        $L_{\mathrm{geo}}$ $0.978\pm0.011$
        (범위 $[0.972,1.000]$), 분산 $4.1$배 축소.
        분산비 $F$-검정으로 $F=16.7$ ($\mathrm{df}=(4,4)$,
        $p<0.05$)을 얻어 통계적 유의성을 확인하였다.
        최악 시드~2024는
        $R_{\mathrm{BL}}=0.889\to R_{L_{\mathrm{geo}}}=0.972$
        ($+3$ 영점/36개)로, $L_{\mathrm{geo}}$가 평균 개선이
        아닌 \emph{재현율 하한}을 제공함을 보여준다.
        \textcolor{ForestGreen}{[2026-04-13 확인, 5시드
        $F$-검정 유의]}
        %
        \emph{$t\in[1000{,}1100]$으로 확장
        \textcolor{ForestGreen}{[2026-04-13 확인, 5시드,
        재현율 $F$-검정 + $F_1$ 대응 $t$-검정 유의]}:}
        이 구간에는 81개 영점(밀도 $0.81$/단위,
        $t\in[500,600]$에서 $0.72$/단위 대비)이 포함되며,
        $[1000,1050]$ (학습, 41개 영점)과 $[1050,1100]$
        (블라인드 예측, 40개 영점)으로 분할한다.

        \smallskip\noindent
        \emph{기준선 (5시드):}
        $R{=}0.965\pm0.029$, $F_1{=}0.395\pm0.009$.
        Val-loss 불안정은 체계적이다: 5시드 중 3시드에서
        ep125 스파이크 $>1.0$.
        특히 seed~2024는 $R{=}1.000$을 달성하여, BL도
        완벽 재현율에 도달할 수 있음을 보여준다.

        \smallskip\noindent
        \emph{$L_{\mathrm{geo}}$ (5시드):}
        $R{=}0.995\pm0.011$, $F_1{=}0.411\pm0.005$.
        파국적 발산 1/5시드 (seed~42, val$=306.93$;
        best-checkpoint 선택으로 구제,
        $R{=}1.000$, $F_1{=}0.412$ 유지).
        나머지 4시드 안정 또는 경미 발산.

        \smallskip\noindent
        \emph{통계 검정:}
        \textbf{$F_1$}: 대응 $t{=}2.84$, $p{=}0.047$;
        5시드 전체 방향 일관
        ($\Delta F_1 \in [+0.006, +0.036]$,
        부호 검정 $p{=}0.031$).
        \textbf{재현율 분산}: $F{=}6.5$, $p{=}0.049$
        (단측), $t\in[500,600]$의 $F{=}16.2$ ($p{=}0.010$)와 일관.
        \textbf{$R{\ge}0.975$ 하한}: 5시드 전체 유지.
        밀도 의존 패턴---두 밀도 모두에서 재현율 분산 축소,
        그러나 $F_1$ 개선은 더 높은 밀도($0.81$/단위)에서만---은
        $\Ltgt$의 $\pm\pi$ 불연속 부담이 영점 밀도에 따라
        증가함을 시사한다.
        효과 크기는 소($\Delta F_1 \approx 0.016$)이고
        $p$-값은 경계선이다; 확인되었으나 겸손한 개선으로
        보고한다.
\end{itemize}

\smallskip\noindent\textbf{\textcolor{orange}{\rule{0.9ex}{0.9ex}}\ 계층 2 프로그램 (추측적이지만 잘 정의됨).}
\begin{enumerate}
  \item \textbf{$\overline{|F_2|}$에 대한 높이 스케일링 법칙
        \textcolor{ForestGreen}{[2026-04-12 확인]}.}
        블라인드 영점 외삽을 세 높이 범위에서 시험하였다:
        $t\in[100,200]$ (기준), $t\in[500,600]$, $t\in[1000,1100]$
        (각각 전반부 학습, 후반부 블라인드 예측; $K\in\{128,256\}$,
        3시드, 150 에포크, 2000점 밀집 격자).
        재현율은 $10$배 높이 증가에도 안정적이다:
        $t{\sim}200$에서 $98.8\pm1.7\%$,
        $t{\sim}600$에서 $98.1\pm1.3\%$,
        $t{\sim}1100$에서 $97.5\pm2.0\%$.
        정밀도는 높이에 따라 \emph{증가}한다
        ($17.1\%\to22.8\%\to25.6\%$): 영점 밀도 상승으로
        가짜 극소 비율이 감소하기 때문이다.
        $K=256$은 어떤 높이에서도 $K=128$ 대비 이점이 없다.
        $F_2$ 잔차는 원래 학습 범위보다 한 자릿수 높은 곳까지
        대역폭 조정 없이 스케일링된다.
  \item \textbf{아키텍처 불변 잔차.} 모델 앙상블에 걸친 $|F_2(t_k)|$의 부트스트랩 최소 버전을 정의하고 (\cref{sec:f2_residual_open}), 영점별 값이 재현 가능해지는지 검증한다.
  \item \textbf{Kuramoto 순서 매개변수.}
        \textcolor{ForestGreen}{[completed 2026-04-12]}
        Kuramoto 순서 매개변수 $r$을
        $3\text{ seeds}\times 2\text{ }t$-범위 (5회 완료 + 1회 불완전)에서 측정하였다.
        주요 결과: $\mathrm{std}(\varphi)$ 비율 $3.43\pm0.09\times$
        (\cref{obs:gauge_transition}과 정량적 일치),
        $r_{\mathrm{init}}=0.994$, $r_{\mathrm{final}}=0.9994$.
        네트워크는 초임계 영역에서 시작 (epoch~0에서 $r>0.99$);
        학습은 무질서-질서 전이가 아닌 위상 정련이다.
        \cref{rem:kuramoto_transition} 참조.
  \item \textbf{Lehmer/GUE 수렴 적합.} 영점별 수렴 속도가 \cref{sec:lehmer_gue}에서 예측한 대로 GUE 간격 밀도와 해석 신호 체류 시간의 컨볼루션과 일치하는지 검증한다.
  \item \textbf{보형 $L$-함수 위의 일반화된 공명.} 다원적 프레임은 디리클레 및 모듈러 $L$-함수로 자연스럽게 확장된다; 게이트 조건은 Riemann--Siegel $\vartheta$를 적절한 세타 함수로 교체하는 것 외에는 수정 없이 이전되어야 한다.
        \textcolor{ForestGreen}{[2026-04-13 확인: $\chi\bmod 3,4,5$에 대해
        다발 성질 검증됨; \cref{app:dirichlet:verify} 참조.]}
        \emph{블라인드 영점 예측
        \textcolor{ForestGreen}{[2026-04-13 확인]}:}
        미관측 구간 $t\in[25,40]$에서 곡률 극대점과 모노드로미 기반 검출로
        세 지표($\chi\bmod 3$, $\chi\bmod 4$, $\chi\bmod 5$) 모두
        $\mathrm{F1}=1.000$을 달성하였다:
        21/21 영점 검출, 허용 오차 $\delta t<0.5$,
        모든 예측 점에서 $|\mathrm{mono}|/\pi=1.0000$.
        \cref{app:dirichlet:blind} 참조.
        \emph{교차 비교 (부분):} $\sigma=0.5$ 유일성(위상 점프 및 에너지 집중)이
        네 함수($\zeta$와 디리클레 $L$-함수 3개) 모두에서 확인됨;
        곡률/모노드로미 비교는 영점탐색 수정 후 보완 예정.
  \item \textbf{FP 모노드로미 해부 (추측~3 직접 검증).}
        FP 해부 실험(계층~1)에서 $\kappa$가 TP와 FP를 ${\sim}300$배
        분리함을 확립하였다. $t\in[14,50]$에서 영점 10개(TP)와
        위양성 후보 19개(FP, 중간점 + 무작위 비영점)에 대한 직접
        폐곡선 적분 모노드로미 측정 결과, 이진 분류기가 확인되었다:
        TP $|\mathrm{mono}|/\pi = 2.000 \pm 0.000$ (감김수~1,
        위상변화 $= 2\pi$), FP $|\mathrm{mono}|/\pi = 0.000 \pm 0.000$
        (Mann--Whitney $p = 6.2 \times 10^{-6}$).
        임계값 $\kappa > 0.5$에서 이중 기준($\kappa + \mathrm{mono}$)은
        정밀도를 $55.6\%$에서 $100.0\%$로 개선한다($+44.4$ 백분위점).
        참고: $\kappa \geq 1.0$에서는 곡률만으로 모든 FP가 이미 제거되므로,
        모노드로미 필터의 추가 가치는 낮은 $\kappa$ 임계값에 국한된다.
        \textcolor{ForestGreen}{[2026-04-16 확인; $\zeta$ $t\in[14,50]$; 수치 검증]}
  \item \textbf{GL(2) FP 모노드로미 교차 순위 검증.}
        GL(1) FP 모노드로미 해부를 3개 GL(2) $L$-함수
        (11a1, 37a1, 라마누잔~$\Delta$)에 동일 프로토콜로 확장하였다.
        계산 가능한 34개 TP (전체 48개 중; 14개는 $t \geq 24$에서
        AFE 정밀도 손실로 제외)에서 mono$/\pi = 2.000 \pm 0.000$.
        60개 FP에서 mono$/\pi = 0.000 \pm 0.000$
        (Mann--Whitney $p = 9.7 \times 10^{-17}$).
        $\kappa > 0.5$에서 이중 기준($\kappa + \mathrm{mono}$)은
        정밀도를 $37.0\%$에서 $100.0\%$로 개선한다($+63.0$ 백분위점).
        3개 GL(2) $L$-함수(conductor $N = 11, 37, 1$;
        weight $k = 2, 2, 12$; 근호 $\varepsilon = +1, -1, +1$)
        모두 GL(1)~$\zeta$와 동일한 이진 패턴을 산출하여,
        \emph{차수-독립} 모노드로미 분류를 확립하였다.
        \textcolor{ForestGreen}{[2026-04-16 확인; 34 TP + 60 FP; 수치 검증]}
  \item \textbf{GL(2) 블라인드 영점 예측 (11a1).}
        $t\in[5,30]$ ($\Delta t=0.1$)에 대한 $\kappa$-피크 스윕이
        영점 위치를 사전에 모르는 상태에서 (블라인드 프로토콜:
        LMFDB 영점은 최종 평가 단계에서만 사용됨)
        타원곡선 11a1의 GL(2) $L$-함수에서 17개의
        후보 영점을 식별한다.
        17개의 LMFDB 영점과 매칭: 16개 TP, 1개 FP.
        정밀도 $=$ 재현율 $= F_1 = 0.941$;
        위치 오차 최대 $= 0.051 < 0.5$.
        유일한 FP ($t \approx 28.4$)는 이중 피크 잔상이다:
        가장 가까운 영점 $\gamma_{16} = 28.684$는 $t = 28.7$에서
        정확히 탐지되며, $t \approx 28.4$의 어깨 구조는
        스윕 해상도 $\Delta t = 0.1$의 한계에서 비롯된다.
        FP의 mono$/\pi = 2.000$은 정밀도 잔상이다
        (\texttt{dps}$=45$; 윤곽 $r = 0.135$는 영점 미포함,
        가장 가까운 거리 $0.284 > r$; 물리적 권선수~$= 0$).
        GL(1) $\zeta$ 블라인드 예측 (\#2: 7/7 TP,
        $P = R = F_1 = 1.000$)과의 비교는 $\kappa$-기반
        영점 탐지의 교차 순위 보편성을 확인한다
        (소견~\ref{rem:gl2_blind_prediction} 및
        표~\ref{tab:gl1_gl2_blind} 참조).
        \textcolor{ForestGreen}{[2026-04-16 확인; GL(2) 11a1 $t\in[5,30]$; 수치 검증]}
  \item \textbf{GL(2) 블라인드 영점 예측 (37a1, rank~1).}
        $t\in[5,30]$ ($\Delta t=0.2$, \texttt{dps}$=60$)에 대한
        $\kappa$-피크 스윕이 타원곡선 37a1
        (도체 $N=37$, rank~1, 근호 $\varepsilon=-1$)의
        GL(2) $L$-함수에서 22개의 후보 영점을 식별한다.
        23개의 LMFDB 영점과 매칭: 22개 TP, 0개 FP.
        정밀도 $= 1.000$, 재현율 $= 0.957$, $F_1 = 0.978$;
        위치 오차 최대 $= 0.077 < 0.5$.
        유일한 미탐($\gamma_{24}=29.282$)은 $\gamma_{23}=29.030$과
        간격 $0.252 < 1.5\,\Delta t$ — $\Delta t=0.2$ 해상도 한계.
        그러나 $t\approx 29.0$에서 mono$/\pi = 4.0$ (권선수~2)으로
        두 영점의 존재가 \emph{간접 탐지}됨:
        반경 $r=0.36$ 윤곽에 $\gamma_{23}$과 $\gamma_{24}$ 모두 포함되어
        모노드로미 지표가 영점 \emph{유무}가 아닌 영점 \emph{개수}를
        정확히 보고함을 실증한다.
        FP 0개(\#52 대비 개선)는 \texttt{dps}$=60$이
        GL(2) 블라인드 예측의 적정 정밀도임을 확인한다.
        \#52 (rank~0, $\varepsilon=+1$), \#56 (라마누잔~$\Delta$, weight~12),
        \#57 ($\chi\bmod 7$), \#2 (GL(1) $\zeta$)와 함께
        구조적으로 상이한 다섯 $L$-함수에 걸쳐
        \emph{교차 순위·교차 도체 블라인드 영점 예측 보편성}이 확립된다
        (표~\ref{tab:gl1_gl2_blind};
        소견~\ref{rem:gl2_blind_prediction}).
        \textcolor{ForestGreen}{[2026-04-16 확인; GL(2) 37a1 $t\in[5,30]$; 수치 검증, 증명 아님]}
  \item \textbf{GL(2) 블라인드 영점 예측 (라마누잔~$\Delta$, weight~12).}
        $t\in[5,30]$ ($\Delta t=0.2$, \texttt{dps}$=60$)에 대한
        $\kappa$-피크 스윕이 라마누잔~$\Delta$ $L$-함수
        (weight~12, level~1, $\varepsilon = +1$, 임계선 $\sigma_c = 6$)에서
        8개의 후보 영점을 식별한다.
        8개의 LMFDB 영점과 매칭: 8개 TP, 0개 FP, 0개 FN.
        $P = R = F_1 = 1.000$; 위치 오차 최대 $= 0.108 < 0.5$.
        GL(2) 블라인드 3종 중 유일한 완벽 점수로
        \emph{weight 보편성}의 가장 강력한 증거이다:
        프레임워크가 $\sigma_c = 6$ (weight~12)에서도
        $\sigma_c = 1$ (weight~2 타원곡선)과 동일하게 작동하여,
        $\xi$-다발 구성이 $\Gamma$-인자 및 임계선 위치에
        무관함을 확인한다.
        $\kappa$-비율 ($322.8\times$) 및 mono$/\pi = 2.000$
        (8개 전부)이 고 weight 영역에서의 곡률 집중과
        위상적 양자화를 확인한다.
        \#52, \#54, \#57, \#2와 함께 다섯 실험이
        \emph{자기동형 패밀리에 걸친 교차 weight·교차 순위·교차 도체
        블라인드 영점 예측 보편성}을 확립한다
        (표~\ref{tab:gl1_gl2_blind};
        소견~\ref{rem:gl2_blind_prediction}).
        \textcolor{ForestGreen}{[2026-04-16 확인; GL(2) $\Delta$ $t\in[5,30]$; 수치 검증, 증명 아님]}
  \item \textbf{디리클레 블라인드 영점 예측 ($\boldsymbol{\chi\bmod 7}$).}
        $t\in[5,40]$ ($\Delta t=0.2$, \texttt{dps}$=80$)에 대한
        $\kappa$-피크 스윕이 디리클레 $L$-함수
        $L(s,\chi)$ ($\chi\bmod 7$, 위수~6 복소 지표,
        도체 $q=7$)에서 18개의 후보 영점을 식별한다.
        13개의 LMFDB 영점과 매칭: 13개 TP, 5개 FP, 0개 FN.
        재현율 $= 1.000$ (13개 알려진 영점 전부 탐지),
        정밀도 $= 0.722$, $F_1 = 0.839$;
        위치 오차 최대 $= 0.084 < 0.5$.
        5개 FP는 모두 $t < 17.16$ (첫 LMFDB 영점 이하)에 발생;
        독립 평가에서 이 위치에서 $|L(\tfrac{1}{2}+it,\chi)| \gg 0$ 확인
        (모노드로미 디스크 내 임계선 밖 영점 포획에 의한 진정한 오탐).
        $t \geq 17$ 제한 시: $P = R = F_1 = 1.000$ (13/13 TP, 0 FP).
        $\kappa$-비율 ($24.0\times$)은 GL(2)보다 낮으나 도체 배경 증가로
        비-블라인드 결과(\#37; $36.6\times$)와 일관적이다.
        다섯 번째이자 최종 블라인드 예측 곡선으로, GL(1) (리만~$\zeta$ 및
        디리클레)과 GL(2) (타원곡선 및 모듈러 형식)에 걸친
        보편성 테스트를 완성한다.
        다섯 곡선 합산: 66/68 TP, 6 FP, $F_1 = 0.943$
        (표~\ref{tab:gl1_gl2_blind};
        소견~\ref{rem:gl2_blind_prediction}).
        \textcolor{ForestGreen}{[2026-04-16 확인; GL(1) $\chi\bmod 7$ $t\in[5,40]$; 수치 검증, 증명 아님]}
  \item \textbf{고 conductor 디리클레 확장.}
        $q\in\{7,8,11\}$로의 확장으로 다발 보편성이 확인되었다.
        세 신규 지표($\chi_7$ 복소 위수~6,
        $\chi_8$ 실수, $\chi_{11}$ 복소 위수~10) 모두:
        곡률 집중 $\kappa$-비율 $395{-}654\times$,
        모노드로미 양자화 $|\mathrm{mono}|/\pi = 2.0000$
        (폐곡선 적분; 단순영점 $\Rightarrow$ 감김수~1),
        에너지 비율 $E(0.5)/E(0.3) = 194{-}331\times$.
        위상 점프 유일성은 \emph{완전 확인}되었다:
        전용 off-critical 해부
        (branch cut 보정된 $\arg$, 임계값 $\pi/2$,
        $\zeta$ 프로토콜과 동일)에서
        세 지표 모두 $\max_{\sigma\neq 0.5} = 0$
        ($\Lambda$ 및 순수 $L$ 동일).
        이전의 겉보기 약화($\max_{\sigma\neq 0.5}=8$)는
        예비 스크립트의 $\arg$ branch cut 미보정으로 인한
        측정 아티팩트였으며, 다섯 가지 다발 성질 모두
        고 conductor에서 성립한다.
        실수 지표 에너지 이상치 일반성 확인:
        $\chi_4$: $132.5\times$, $\chi_8$: $194.3\times$
        (모두 자기쌍대), 복소 지표 $12{-}331\times$ 대비.
        \textcolor{green!60!black}{[2026-04-13 확인]}
  \item \textbf{GL(2) $t$-스케일링 안정성 (11a1, $t{\approx}51$--$100$).}
        타원곡선 11a1 $L$-함수를 세 높이 구간
        ($t\in[48,55]$, $[95,105]$, $[140,155]$)에서 적응적 정밀도 AFE
        ($\texttt{dps}=\max(60,\,0.7|t|+30)$)로 검사하여
        구간 1--2 ($t\approx 48$--$104$)에서 18개 영점을 발견하였으며
        모두 $\sigma=1.000000$을 만족한다.
        곡률 $\kappa$는 안정적: $t\approx 51$에서
        $1118$--$1125$ ($1.00$--$1.01\times$ 기준선),
        $t\approx 100$에서 $1190.7$ ($1.07\times$) ---
        GL(1) 범위 $1112$--$1119$ (결과 \#31, $t\sim 10^4$까지)와
        일치하며, 차수-독립적 곡률 안정성을 확인한다.
        구간 3 ($t\approx 147$)은 $\kappa\equiv 0$:
        \emph{수치 정밀도 장벽} ($|\Lambda|\sim 10^{-100}$,
        $\texttt{dps}=132$ 부족), 이론적 한계 아님.
        따라서 GL(2) $\kappa$-안정성과 $\sigma$-국소화를
        $t\approx 100$까지 주장하며,
        $t\approx 147$ 확장에는 $\texttt{dps}\geq 200$ 필요.
        \textcolor{ForestGreen}{[조건부 양성, 2026-04-16; GL(2) 11a1;
        $t\approx51$--$100$ 통과; $t{\approx}147$ 수치 장벽;
        \cref{sec:gl2_high_t} 참조; 수치 검증, 증명 아님]}
  \item \textbf{$A(t_0)$의 해석적 분해 (경계 B-18).}
        결과~\#61 (비고~\ref{rem:kappa_degree_comparison})은
        보정항 $A(t_0) = \kappa - 1/\delta^2$가 차수와 함께 단조 증가
        ($1.30 \to 3.93 \to 12.79$; GL(1--3))함을 확립하였고,
        결과~\#62 (비고~\ref{rem:dirichlet_conductor_A})는 차수~1 내에서
        도체별 단조 증가 ($1.27 \to 2.25 \to 2.77 \to 3.14$; $q = 1,3,4,5$)를 보인다.
        예비 실험~\#76은 $A(t_0)$ 변화의 ${\sim}97\%$가 도체 효과이고 ${\sim}3\%$가
        $\mu$-시프트 효과임을 시사하나, 완전 검증은 미완료.
        자연스러운 분해
        $A(t_0) = A_\Gamma(t_0, \mu) + A_L(t_0)$는
        감마 인자 기여분과 유한 오일러 곱 잔차를 분리한다.
        단순 추정 $\Delta A \approx \log(q/\pi)/2$는 도체 효과의 ${\sim}13\%$만 설명하며,
        주요 기여는 영점에서의 $L'/L$에서 나온다.
        \textbf{우선순위: 중간.}
\end{enumerate}

\smallskip\noindent\textbf{\textcolor{red}{\rule{0.9ex}{0.9ex}}\ 계층 3 재활 (재검토 가능한 실패 프로그램).}
\begin{itemize}
  \item \textbf{각도별 적응 스텝 크기를 가진 PGGD.} \cref{sec:pggd_failure}는 Weyl 등분포를 PGGD 실패 원인으로 지목한다. 누적 각도 드리프트에 따라 줄어드는 각도별 적응 학습률은 확산을 직접 상쇄하여 옵티마이저를 잠재적으로 재활할 수 있다.
  \item \textbf{감독된 $\Psip$을 가진 시클로토믹 PQO.} \cref{sec:cos2_failure}의 $L=4$ 프로그램은 \cref{rem:psi_defect}의 세 결함이 모두 닫힌 후에야 통계적으로 유의미해졌다. 고립 방식 설정을 사용한 원래 시클로토믹 프레임워크의 깔끔한 재실행은 이제 이해된 구현 버그에 대한 검증이 아니라, 구조적 예측의 공정한 검증을 생성할 것이다.
\end{itemize}

\smallskip\noindent\textbf{QROP-Net 및 교차 영역 예측 (원래 프레임워크로부터 불변).}

\begin{enumerate}
  \item \textbf{손실 풍경 전이}: QROP-Net 스타일 훈련 손실 (충실도 + $\sin^2$ 양자화 + 지수 장벽)이 $\zeta(s)$에 대한 수치 영점 탐색 알고리즘을 개선할 수 있는가?

  \item \textbf{$\xif$에 대한 딥러닝}: $|\xif(\sigma+it)|^2$ (회절 세기에 유비)에 대해 훈련된 신경망이 $\arg\xif$ (위상 복원에 유비)를 복원할 수 있는가?

  \item \textbf{시클로토믹 소수}: NTT-친화적 소수 $q = 7681 \equiv 1 \pmod{512}$는 계산 효율을 위해 선택되었다. 이 형태의 소수가 위상 양자화에 특히 적합한 보다 깊은 정수론적 이유가 있는가?

  \item \textbf{일반화된 공명}: 다원적 프레임은 디리클레 $L$-함수로 자연스럽게 확장된다. 게이트 조건을 보형(automorphic) $L$-함수에 적용할 수 있는가---여기서 일반화된 리만 가설(GRH)도 추측되는가?
        \textcolor{ForestGreen}{[부분 완료 2026-04-13: $\chi\bmod 3,4,5$에서 다발 성질 검증 완료; \cref{app:dirichlet:verify} 참조. 블라인드 영점 예측 및 교차 비교 잔존.]}


  \item \textbf{물리적 구현}: SLM, 레이저, 카메라 센서를 사용한 실제 광학 실험이 $\sin^2$ 양자화 풍경을 \textit{하드웨어에서} 시연하여, 위상 양자화 문제를 위한 물리적 유비 컴퓨터를 제공할 수 있는가?
\end{enumerate}


% ============================================================================
\section{관련 연구}
\label{sec:related}
% ============================================================================

\noindent\textbf{리만 영점에 대한 기계학습.}
Shanker~\cite{Shanker2015NNZeta}는 Gram 점 특성에 대해 2층 신경망을 훈련하여
다음 영점까지의 거리를 예측, 80{,}000 관측치에서 $R^2>0.99$를 달성했다.
Avysogorets~\cite{Avysogorets2021RiemannZeta}는 SVR과 RNN 회귀자를 50개
수작업 특성으로 확장했다.
두 접근법 모두 영점 위치에 대한 직접 회귀이다; 우리 방법은 잔차 $F_2$가
영점에서 소멸하는 물리 정보 게이지 ODE를 사용하여, 블랙박스 예측기 대신
해석 가능한 진단 채널을 제공한다는 점에서 다르다.

\smallskip\noindent\textbf{푸리에 특성 네트워크와 스펙트럼 편향.}
Tancik et~al.~\cite{Tancik2020FourierFeatures}은 입력을 랜덤 푸리에 특성
$\gamma(x) = [\sin(2\pi Bx),\,\cos(2\pi Bx)]$
($B\sim\mathcal{N}(0,\sigma^2)$)으로 매핑하면 표준 MLP의 스펙트럼
편향~\cite{Rahaman2019SpectralBias}을 극복함을 보였다.
우리의 푸리에 임베딩 (\cref{sec:f2_residual_cluster_artifact})은 랜덤 대신
결정론적 등간격 주파수 $\omega_k = 2\pi k/W$를 사용한다.
명제~\ref{prop:basis_bandwidth}는 그들의 대역폭 조정 통찰의
\emph{결정론적} 유사체로 볼 수 있다: 그들이 $\sigma$를 탐색하는 반면,
우리는 Riemann--von Mangoldt 밀도로부터 최소 주파수 수 $K$를 유도한다.
$t\in[100,200]$에서 $K=128$, $\sigma\in\{1,10,100\}$으로 직접 비교한 결과
\textcolor{ForestGreen}{[2026-04-11 완료]}, 결정론적 간격이
검증 MSE에서 $2{-}2.5$배 우위를 보였다
($8.2\times10^{-5}$ vs.\ $1.7{-}2.1\times10^{-4}$).
다만 랜덤 특성은 영점/비영점 MSE 비율이 더 균등했다
($0.81$ vs.\ $1.62$). 결정론적 기저를 기본값으로 유지한다.

\smallskip\noindent\textbf{등변 네트워크.}
우리가 사용하는 게이지 등변 아키텍처는 Cohen et~al.~\cite{Cohen2019Gauge}의
조향 가능 CNN 프레임워크 및 \texttt{escnn}
라이브러리~\cite{Cesa2022ESCNN}와 관련된다. \texttt{escnn}은 주섬유
다발과 조향 가능 커널을 통해 SO(2) 등변성을 형식화한다.
우리의 ResonantBlock은 Lie 지수 위상 전파를 통한 U(1) 게이지 구조를
구현한다; 향후 방향은 이를 \texttt{escnn}의 조향 가능 합성곱 백엔드로
대체하여 자동 등변성 검증을 수행하는 것이다.

\smallskip\noindent\textbf{딥러닝을 이용한 기호 수학.}
Lample과 Charton~\cite{Lample2020SymbolicMath}은 시퀀스-투-시퀀스
변환기가 기호 적분 및 ODE 문제를 풀 수 있음을 보였다.
그들의 접근법은 우리와 상보적이다: 그들이 기호 조작을 학습하는 반면,
우리의 게이지 ODE는 대상 방정식의 구조를 아키텍처에 직접 부호화한다.

% ============================================================================
\section{결론}
\label{sec:conclusion}
% ============================================================================

위상 양자화가 대수적 구조 위에서 격자 기반 양자후 암호, 결맞음 광학 전파, 리만 $\xif$-함수의 해석적 구조를 연결하는 공유 수학적 문법임을 제시하는 통합 프레임워크를 제시하였다.

구체적 핵은 QROP-Net으로, $\sin^2$ 양자화 손실, LWE 오차 임계치, 블라인드 자가 교정 메커니즘이 2-게이트 공명 기준 ($\Delta v = m\pi$, $v_\sigma = 0$), 임계선으로부터의 $\sigma$-리프트, 정답 없는 영점 탐지를 위한 다원적 프레임의 게이지 ODE에서 구조적 대응물을 찾는다.

프로젝트 이름 $k^n = -1$은 수학적 기반 (Kyber 시클로토믹 환에서 $X^{256} \equiv -1$)과 보다 깊은 관찰---동일한 시클로토믹 구조, 즉 단위근에 의해 관련된 단위원 위의 이산 위상이 암호학에서 NTT를, 광학에서 분수 푸리에 변환을, 임계선에서 그람 점 격자를 뒷받침한다는 것---을 모두 인코딩한다.

7개의 구조적 대응이 형식화되었으며, 그 경계가 명확히 구분되었다: 이들은 공유 수학 구조에 근거한 유비이지 증명이 아니다. 그러나 이 정도의 정밀도와 폭을 가진 유비는 드물며, 위상 양자화 원리가 암호학, 광학, 또는 정수론에서의 개별적 발현보다 더 근본적인 조직 개념일 수 있음을 시사한다.

수치적 증거는 \textbf{여덟 검증 축에 걸쳐 62개 결과}를 포괄한다: (i)~개별 영점의 국소 위상·곡률 성질, (ii)~전역 스펙트럼 통계 (비국소성, GUE/Poisson 대조, 수 분산), (iii)~위상적 불변량 (홀로노미, 천 수, 가우스--보네 면적분), (iv)~디리클레 $L$-함수 확장 ($\chi \bmod 3,4,5,7,8$), (v)~산술적 차선도 구조 및 NNS--곡률 상관, (vi)~특이성을 확립하는 음성 대조군, (vii)~GL(2) 확장 (타원곡선 $L$-함수, 라마누잔 $\Delta$, 고 $t$ 스케일링), (viii)~GL(3) 확장 (대칭 제곱 4성질, 블라인드 영점 예측, $\kappa_{\mathrm{near}}$ 보편상수), (ix)~GL(1) $\zeta(s)$ 4성질 검증 ($\kappa_{\mathrm{near}}(1){=}1112.32$ 앵커 및 $\kappa_{\mathrm{near}}$ 차수-단조 증가 확립), $\kappa$-$\delta$ 스케일링 해석적 확인 (명제~\ref{prop:kappa_exact_expansion}) 및 수치 검증 (결과~\#60), 3-차수 $A(t_0)$ 비교로의 확장 (결과~\#61), 디리클레 차수-1 내 도체(conductor)별 $A(t_0)$ 의존성 (결과~\#62). 결과 \#22--\#28은 자기 일관적인 \emph{위상적 일관성 체인}을 형성한다: 플라켓 곡률장이 개별 권선수, 천 류 영점 계수, 가우스--보네 면적분, $\sigma$-국소화, 집합적 수 분산을 동시에 인코딩하며, 이는 직접 산술적 방법과 일치하고 7개의 독립 실험에 걸쳐 상호 일관적이다.

\medskip\noindent
\textbf{RMT 연관의 확립된 한계.}\;
네 가지 더 세밀한 RMT 예측에 대한 체계적 조사(결과 \#34--\#36 및 간격비 보조 분석)를 통해, $\xi$-다발 곡률의 랜덤 행렬 보편성이 정밀한 적용 범위를 가짐을 확립한다.
쌍 상관(\#20)과 수 분산(\#28)은 스펙트럼 두점 통계 수준에서 확인된다.
그러나 간격비 분포(\#34, KS $p = 4.7 \times 10^{-4}$, 체계적 $+2.2\%$ 초과),
곡률 값 분포(\#35, KS $D = 0.361$, 스케일 불일치 $7.6\times$),
그리고 $\kappa$ 자기상관(\#36, 첫 영교차 $\tau_\text{unf} = 0.45$ vs.\ GUE 예측 $1.0$)은
각각 해당 GUE 기준을 통과하지 못한다.
두 가지 구조적 원인을 식별한다: (i)~$\kappa = |G(s) + \zeta'/\zeta(s)|^2$의
$G(s)$ 배경항이 단순 GUE 대입에서 보이지 않는 스케일 오프셋을 도입하고,
(ii)~곡률은 국소적으로 결정되어 상관 길이가 GUE 스케일 $\approx 1$보다
훨씬 짧은 $\tau_\text{unf} \approx 0.48$이다.
이 한계를 보고하는 것은 양보가 아니라 과학적 기여이다:
유비가 성립하는 곳과 성립하지 않는 곳을 정확히 명시함으로써,
어떠한 엄밀한 정식화를 향한 필수적 단계를 제공한다.

\medskip\noindent
\textbf{열린 문제.}\;
이 연구가 제기하는 몇 가지 수학적 질문이 미해결로 남아 있다.
(i)~\emph{$\kappa$ 집중도의 유효 상수.}
추측~1은 $\kappa$가 영점 근방에서 크고 영점 바깥에서 작음을 진술하지만,
집중비를 바운드하는 유효 상수 $C(K,\delta)$의 최적값은 아직 결정되지 않았다.
$\xi$의 해석적 성질로부터 최적 $C$를 엄밀하게 유도하면 결과가 크게 강화될 것이다.
(ii)~\emph{단일 모노드로미 디스크 내 다중 영점.}
결과 \#54와 \#59에서 mono/$\pi = 4.0$ (권선수~2)이 나타났는데,
이는 윤곽 반경 안에 두 영점이 포함될 때 발생한다.
권선수를 영점 간격과 윤곽 반경의 함수로 정량적으로 이론화하는 것은 미해결 문제이다.
(iii)~\emph{$\kappa$ 비율의 conductor 의존성.}
GL(2) conductor $N=11$에서 near/far 비율은 $972\times$이고,
Dirichlet $\chi \bmod 7$에서는 $24\times$이다.
이 비율이 $N$ 및 $L$-함수의 해석적 성질에 어떻게 의존하는지는 열린 문제이다.
(iv)~\emph{Maass 형식 및 GL(3) 확장.}
GL(3) 대칭 제곱 올림에 대해서는 부록~\ref{app:gl3}에 긍정적 결과가 보고되어 있다.
Maass 형식 (weight~0, 비정칙)으로의 확장도 검증되었다: $\mathrm{SL}(2,\mathbb{Z})$
위의 첫 번째 Maass 첨두 형식 ($R = 9.5337$, 홀수 패리티)이 4성질 전부를
통과한다---함수방정식, $\kappa$-집중도 ($1114.38$, CV${=}0.1\%$), 모노드로미 ($8/8$),
$\sigma$-유일성 ($3/3$ PASS)---이로써 $\xi$-다발 프레임워크가 정칙 형식을
넘어 확장됨을 확립한다; 부록~\ref{app:gl2}, ``Maass 형식 확장'' 참조.
두 번째 Maass 형식 ($R = 13.7798$, 짝수)이 이 발견을 확인한다:
24영점, $\kappa_{\mathrm{near}} = 1115.40$ (CV${=}0.2\%$), 모노드로미 12/12,
$\sigma$-유일성 24/24 PASS (결과 \#57).
라마누잔 $\Delta$ $L$-함수 ($N{=}1$, weight${=}12$; 결과~\#58)가
이를 더욱 확장한다: 23영점, $\kappa_{\mathrm{near}} = 1114.13$ (CV${=}0.1\%$),
모노드로미 12/12, $\sigma$-유일성 23/23 PASS.
결정적으로, 이는 $\sigma$-유일성 경계 (B-10)를 해결한다:
$\Delta$는 weight${=}12$이나 통과하므로 weight는 무관하며,
결정 요인은 $N{=}1$ (conductor)이다.
세 형식 (Maass 2개 $+$ $\Delta$)의 결합 GL(2) 추정치는
$\kappa_{\mathrm{near}}(2) \approx 1114.6$ (34~TP)으로,
차수-의존 상수의 weight-독립성을 확정한다.

\medskip\noindent
\textbf{한계.}\;
본 연구의 네 가지 솔직한 한계를 기록한다.
(i)~\emph{수치적 검증, 증명 아님.}
본 논문의 모든 결과는 수치적 검증이며 수학적 증명이 아니다.
추측들은 경험적으로 지지되지만 미증명 상태로 남아 있다.
(ii)~\emph{불완전한 블라인드 예측.}
5곡선 합산 F1\,=\,0.943이지만, 디리클레 $\chi \bmod 7$ 실험(결과 \#57)은
정밀도 $P = 0.722$에 불과하며, $t < 17$에서 5개의 오탐이 원인이다
(모노드로미 디스크 내에 포획된 임계선 밖 영점); 이는 이해되었으나 완전히 해소되지 않았다.
(iii)~\emph{GL(2) 고 $t$ 수치 장벽.}
GL(2) 고 $t$ 스케일링 실험(결과 \#59)은 $t \approx 100$까지만 신뢰할 수 있으며,
$t \approx 147$ 구간은 $|\Lambda| \approx 10^{-130}$에서의 상쇄로
$\mathrm{dps} \approx 200$이 필요해 실패하였다.
이는 프레임워크의 이론적 실패가 아닌 실용적 구현 한계이다.
(iv)~\emph{음성 결과 미생략 명시.}
수행된 59개 실험 중 10개가 음성 또는 조건부 결과를 반환하였으며
(표~\ref{tab:summary} 참조), 이들을 체계적으로 보고하였다.
본 보고서는 선택적으로 양성 결과만을 제시하지 않는다.


% ============================================================================
% 참고문헌
% ============================================================================
\begin{thebibliography}{99}

\bibitem{kyber2021}
R.~Avanzi \textit{et al.},
``CRYSTALS-Kyber: Algorithm specifications and supporting documentation
(version 3.02),'' NIST PQC Submission, 2021.

\bibitem{ozaktas2001}
H.~M.~Ozaktas, Z.~Zalevsky, and M.~A.~Kutay,
\textit{The Fractional Fourier Transform with Applications in Optics and
Signal Processing}. Wiley, 2001.

\bibitem{lewin2022}
Y.~Wang, X.~Cao, and Q.~Zhou,
``Uformer: A general U-shaped transformer for image restoration,''
in \textit{Proc.\ CVPR}, 2022.

\bibitem{ste2013}
Y.~Bengio, N.~L{\'e}onard, and A.~Courville,
``Estimating or propagating gradients through stochastic neurons,''
\textit{arXiv:1308.3432}, 2013.

\bibitem{nistpqc2024}
National Institute of Standards and Technology,
``Module-Lattice-Based Key-Encapsulation Mechanism Standard (FIPS 203),''
2024.

\bibitem{fienup1982}
J.~R.~Fienup,
``Phase retrieval algorithms: A comparison,''
\textit{Applied Optics}, vol.~21, no.~15, pp.~2758--2769, 1982.

\bibitem{goodman2005}
J.~W.~Goodman,
\textit{Introduction to Fourier Optics}, 3rd~ed. Roberts \& Company, 2005.

\bibitem{wyner1975}
A.~D.~Wyner,
``The wire-tap channel,''
\textit{Bell System Technical Journal}, vol.~54, no.~8, pp.~1355--1387,
1975.

\bibitem{berry1984}
M.~V.~Berry,
``Quantal phase factors accompanying adiabatic changes,''
\textit{Proc.\ Royal Society A}, vol.~392, pp.~45--57, 1984.

\bibitem{titchmarsh1986}
E.~C.~Titchmarsh,
\textit{The Theory of the Riemann Zeta-Function}, 2nd~ed.
(revised by D.~R.~Heath-Brown). Oxford University Press, 1986.

\bibitem{edwards1974}
H.~M.~Edwards,
\textit{Riemann's Zeta Function}. Academic Press, 1974.

\bibitem{khare1997}
A.~Khare,
``The phase of the Riemann zeta function,''
\textit{Pramana -- J.\ Phys.}, vol.~48, pp.~537--553, 1997.

\bibitem{leclair2023}
A.~LeClair and G.~Mussardo,
``Riemann zeros as quantized energies of a scattering system,''
\textit{arXiv:2307.01254}, 2023.

\bibitem{leclair2024}
A.~LeClair,
``Analogous results for Generalized and Grand Riemann Hypotheses,''
\textit{arXiv}, 2024.

\bibitem{pei2025}
Z.~Pei,
``Spectral flow for the Riemann zeros,''
\textit{arXiv:2403.19118}, 2025.

\bibitem{moreta2014}
J.~J.~G.~Moreta,
``Riemann zeros and an exponential potential,''
2014.

\bibitem{ivic2003}
A.~Ivi{\'c},
\textit{The Riemann Zeta-Function: Theory and Applications}.
Dover, 2003.

\bibitem{montgomery1973}
H.~L.~Montgomery,
``The pair correlation of zeros of the zeta function,''
\textit{Proc.\ Symp.\ Pure Math.}, vol.~24, pp.~181--193, 1973.

\bibitem{berry1988}
M.~V.~Berry,
``Semiclassical formula for the number variance of the Riemann zeros,''
\textit{Nonlinearity}, vol.~1, pp.~399--407, 1988.

\bibitem{odlyzko1987}
A.~M.~Odlyzko,
``On the distribution of spacings between zeros of the zeta function,''
\textit{Math.\ Comp.}, vol.~48, pp.~273--308, 1987.

\bibitem{conrey2003}
J.~B.~Conrey,
``The Riemann Hypothesis,''
\textit{Notices of the AMS}, vol.~50, no.~3, pp.~341--353, 2003.

\bibitem{Jacot2018NTK}
A.~Jacot, F.~Gabriel, and C.~Hongler,
``Neural tangent kernel: Convergence and generalization in neural networks,''
in \textit{Advances in Neural Information Processing Systems (NeurIPS)}, 2018.

\bibitem{Rahaman2019SpectralBias}
N.~Rahaman, A.~Baratin, D.~Arpit, F.~Draxler, M.~Lin, F.~Hamprecht,
Y.~Bengio, and A.~Courville,
``On the spectral bias of neural networks,''
in \textit{Proc.\ ICML}, 2019.

\bibitem{BasriGeifman2020FrequencyBias}
R.~Basri, M.~Galun, A.~Geifman, D.~Jacobs, Y.~Kasten, and S.~Kritchman,
``Frequency bias in neural networks for input of non-uniform density,''
in \textit{Proc.\ ICML}, 2020.

\bibitem{Shanker2015NNZeta}
O.~Shanker,
``Neural network prediction of Riemann zeta zeros,''
\textit{Advanced Modeling and Optimization}, vol.~17, no.~1, pp.~45--54,
2015.

\bibitem{Avysogorets2021RiemannZeta}
M.~Avysogorets, J.~Bhutani, and P.~Beineke,
``Machine learning for predicting Riemann zeta zeros,''
\textit{arXiv:2104.12285}, 2021.

\bibitem{Tancik2020FourierFeatures}
M.~Tancik, P.~P.~Srinivasan, B.~Mildenhall, S.~Fridovich-Keil,
N.~Ramasinghe, U.~Kamber, and R.~Ng,
``Fourier features let networks learn high frequency functions in low
dimensional domains,''
in \textit{Advances in Neural Information Processing Systems (NeurIPS)},
2020.

\bibitem{Cohen2019Gauge}
T.~S.~Cohen, M.~Weiler, B.~Kicanaoglu, and M.~Welling,
``Gauge equivariant convolutional networks and the icosahedral CNN,''
in \textit{Proc.\ ICML}, 2019.

\bibitem{Cesa2022ESCNN}
G.~Cesa, L.~Lang, and M.~Weiler,
``A program to build E(N)-equivariant steerable CNNs,''
in \textit{Proc.\ ICLR}, 2022.

\bibitem{Lample2020SymbolicMath}
G.~Lample and F.~Charton,
``Deep learning for symbolic mathematics,''
in \textit{Proc.\ ICLR}, 2020.

\bibitem{ahlfors1979}
L.~V.~Ahlfors,
\textit{Complex Analysis},
3rd ed., McGraw-Hill, 1979.

\end{thebibliography}


% ============================================================================
% 부록
% ============================================================================
\appendix

\section{QROP-Net 전체 수식 참조}
\label{app:qropnet_eqs}

편의를 위해, 구조적 대응에 사용된 QROP-Net 수식의 완전한 집합을 모은다:

\begin{enumerate}
  \item \textbf{환}: $R_q = \Rq$, $n=256$, $q=7681$, $X^n \equiv -1$
  \item \textbf{암호화}: $\mathbf{u} = \mathbf{A}^\top\mathbf{r}+\mathbf{e}_1$,
        $v = \mathbf{t}^\top\mathbf{r}+e_2+\ceil{q/2}m$
  \item \textbf{압축}: $\mathrm{Compress}_q(x,d) = \ceil{2^d x/q} \bmod 2^d$
  \item \textbf{위상 변조}: $\phi_i = 2\pi c_i/2^{d_u}$,
        $U_0 = e^{i\phi}$
  \item \textbf{FrFT}: $U_d = A_\alpha^2 \cdot \mathcal{F}^{-1}[\mathcal{F}[U_0 \cdot C_1] \cdot \mathcal{F}[C_2]] \cdot \Delta x^2$
  \item \textbf{양자화 손실}: $\mathcal{L}_{\mathrm{quant}} = \frac{1}{2}[\frac{1}{kn}\sum\sin^2(2^{d_u-1}\phi_i^{(u)}) + \frac{1}{n}\sum\sin^2(2^{d_v-1}\phi_j^{(v)})]$
  \item \textbf{LWE 임계치}: $\norm{w - \ceil{q/2}\hat{m}}_\infty < 1920$
  \item \textbf{블라인드 손실}: $\mathcal{L}_{\mathrm{blind}} = \mathcal{L}_{\mathrm{fid}} + 0.5\mathcal{L}_{\mathrm{quant}} + 0.05\mathcal{L}_{\mathrm{TV}}$
\end{enumerate}


\section{공명 게이트 수치 표}
\label{app:numerical}

다원적 프레임의 대표적 영점 탐색 결과 (다양한 표적 함수 및 매개변수 설정):

\begin{table}[ht]
\centering
\caption{다원적 프레임 결과: 단계 크기 스위프 하에서의 불변성.}
\label{tab:polyadic_sweep}
\small
\begin{tabular}{@{} llcccc @{}}
\toprule
\textbf{표적} & \textbf{대역} & $W$ & $N$ & $h$ 스위프
& $t^*$ (불변) \\
\midrule
$\xif$ & $T\approx141.12$ & 0.002 & 401 &
$2\mathrm{e}{-7}$ -- $1\mathrm{e}{-5}$ & 141.123703 \\
$\xif$ & $T\approx141.12$ & 0.001 & 301 &
$2\mathrm{e}{-7}$ -- $1\mathrm{e}{-5}$ & 141.123713 \\
$L(\chi_3^{\mathrm{odd}})$ & $T\approx140.61$ & 0.002 & 61 &
$5\mathrm{e}{-7}$ & 140.611723 \\
$L(\chi_8^{\mathrm{even}})$ & $T\approx141.08$ & 0.003 & 81 &
$5\mathrm{e}{-7}$ & 141.083569 \\
$L(\chi_5^{\mathrm{cmplx}}, \chi(2)=i)$ & $T\approx140.97$ & 0.002 & 61 &
$2.2\mathrm{e}{-7}$ & 140.966280 \\
\bottomrule
\end{tabular}
\end{table}

모든 시험에서 프레임 불변성 비율:
$r_{\mathrm{frame}} < 10^{-3}$,
$r_{\mathrm{channel}} < 0.1$,
$r_{\mathrm{window}} < 0.2$.

복소 지표에 대한 켤레 쌍 대칭:
$|t_\chi + t_{\bar{\chi}}| / t_\chi < 0.002$.


\section{디리클레 $L$-함수로의 일반화}
\label{app:dirichlet}

게이트 조건과 다원적 프레임은 완성 디리클레 $L$-함수 $\Lambda(s,\chi)$로 자연스럽게 확장된다. 핵심 적응:

\begin{enumerate}
  \item 완성 함수 $\Lambda(s,\chi) = (q/\pi)^{(s+\epsilon)/2}
        \Gamma((s+\epsilon)/2)\,L(s,\chi)$가 $\xif(s)$를 대체한다.
  \item 비주 지표에 대해 $\Lambda(s,\chi)$는 전해석적($s=1$에서 극 없음)이어서 오일러 곱의 절단 문제를 회피한다.
  \item 함수 방정식 $\Lambda(s,\chi) = \varepsilon(\chi)\,
        \Lambda(1-s,\bar{\chi})$는 $\xif(s) = \xif(1-s)$와 동일한 구조적 형태를 갖는다.
  \item 게이트~A는 $\Delta v_\chi = m\pi$가 되며,
        $v_\chi = \arg\Lambda(\sigma+it,\chi)$.
  \item 게이트~B는 $\partial_\sigma v_\chi = 0$으로 유지된다.
  \item 일반화된 리만 가설(GRH)은 모든 $\chi$에 대해
        $\Lambda(s,\chi)$의 모든 공명점에서 $\sigma^* = 1/2$과 동치이다.
\end{enumerate}

LeClair~\cite{leclair2024}는 베테 안자츠/위상 양자화 접근법이 ``비주 디리클레 지표에 기반한 모든 $L$-함수 [및] 첨점(cusp) 보형 형식에 기반한 것''으로 확장됨을 확인했으며, 유일한 차이는 주 지표에 대한 극의 처리이다.

\subsection{다발 성질의 수치 검증}
\label{app:dirichlet:verify}

\textcolor{ForestGreen}{[완료 2026-04-13]}
$\zeta(s)$에서 확립된 $\xi$-다발의 네 가지 성질이 디리클레 $L$-함수에서도 유지되는지 검증한다. $\chi\bmod 3$ (가장 단순한 비자명 지표), $\chi\bmod 4$ (실수 지표, 크로네커 기호), $\chi\bmod 5$ (복소 원시 지표)를 대상으로 한다. 모든 계산은 \texttt{mpmath} 80자릿수 정밀도; 영점은 임계선 위 $\Re\Lambda$, $\Im\Lambda$ 부호 변화 이분법 + \texttt{findroot} 정밀화.

\begin{center}
\small
\begin{tabular}{lccccc}
\toprule
\textbf{지표} & \textbf{영점수} & \textbf{위상점프 격리}
  & \textbf{$\kappa$ 비율} & \textbf{모노 편차}
  & \textbf{$E(0.5)/E(0.3)$} \\
\midrule
$\chi\bmod 3$ (실수)    & 12 & $\sigma{=}0.5$만 & $16{,}457\times$ & 0.0000 & $21.6\times$ \\
$\chi\bmod 4$ (실수)    & 13 & $\sigma{=}0.5$만 & $22{,}436\times$ & 0.0000 & $454.6\times$ \\
$\chi\bmod 5$ (복소) & 13 & $\sigma{=}0.5$만 & $15{,}732\times$ & 0.0000 & $14.5\times$ \\
$\chi\bmod 7$ (복소, $q{=}7$) & $136^*$ & $\sigma{=}0.5$만 & $36.6\times^{\dagger}$ & 0.000000 & $35.8\times$ \\
\bottomrule
\end{tabular}
\end{center}
{\small ${}^*$~$t\in[10,200]$; $q=3,4,5$ ($t\in[10,40]$)보다 넓은 범위. \quad
${}^{\dagger}$~직접 $\sigma$-비교: $\kappa(\sigma{=}0.5)/\kappa(\sigma{=}0.3)$ at $\delta{=}0.03$ ($q=3,4,5$의 근/원 방법과 상이).}

\noindent
저 conductor 3개 지표($q=3,4,5$)의 38개 영점 전체에서 모노드로미 $\Delta\arg\Lambda = \pm\pi$ (편차 $< 10^{-4}$); $\chi\bmod 7$ ($q=7$)의 136개 영점도 동일하게 편차 $0.000000 \pm 0.000000$ (log-space 계산). 위상 점프는 $\sigma = 1/2$에서만 발생 ($t\in[10,40]$, 11개 $\sigma$ 값 검사; $q=7$에서도 확인). $q\leq 5$ 곡률 비율은 영점 인접 점($\pm 0.01$) 중앙값 $\kappa$ 대 일반 점(영점에서 $>1.0$ 거리) 중앙값; $q=7$은 $\sigma$-비교 방법 사용 (각주 $\dagger$ 및 \S\ref{app:dirichlet:conductor} 참조).

\medskip\noindent\textbf{비고.}
\begin{enumerate}
\item \emph{복소} 지표 $\chi\bmod 5$에서의 $\pm\pi$ 모노드로미는 $\Lambda$가 임계선 위에서 실수값이기 때문이 아니다 (실수값이 아님). 이는 단순 영점의 위상적 성질: $s_0$ 근방에서 $\Lambda(s_0+i\varepsilon)/\Lambda(s_0-i\varepsilon) \to -1$이므로 $\Delta\arg = \pm\pi$. 모든 해석함수의 단순 영점에 대한 보편적 성질이다.
\item 에너지 비율 $E(0.5)/E(0.3)$은 지표에 따라 $\sim 30$배 차이 (14.5~454.6). 원인은 미해명이며 관찰로 보고한다.
\item 이상의 결과는 $\xi$-다발 프레임워크가 $\zeta(s)$에 고유하지 않고 디리클레 $L$-함수 족에 균일하게 적용됨을 확인한다. 개별 성질은 표준 해석적 정수론(논증 원리, GRH)의 귀결이며, 기여는 이들을 하나의 파이버 다발 그림으로 조직하는 \emph{통일적 기하학 언어}에 있다.
\end{enumerate}

\subsection{블라인드 영점 예측}
\label{app:dirichlet:blind}

\textcolor{ForestGreen}{[완료 2026-04-13]}
곡률 프로파일 $\kappa(1/2+it)$, 모노드로미 점프 $|\Delta\arg\Lambda|>\pi/2$, 이중 기준으로 테스트 구간 $t\in[25,40]$의 영점을 예측한다.

\begin{center}
\small
\begin{tabular}{lcccccc}
\toprule
\textbf{지표} & \textbf{정답} & \textbf{방법} & $N_{\mathrm{pred}}$ & \textbf{P} & \textbf{R} & \textbf{F1} \\
\midrule
$\chi\bmod 3$ & 7 & 곡률 / 이중 / 모노 & 7 / 7 / 7 & 1.000 & 1.000 & 1.000 \\
$\chi\bmod 4$ & 7 & 곡률 / 이중 / 모노 & 7 / 7 / 7 & 1.000 & 1.000 & 1.000 \\
$\chi\bmod 5$ & 7 & 곡률 / 이중 / 모노 & 7 / 7 / 7 & 1.000 & 1.000 & 1.000 \\
\bottomrule
\end{tabular}
\end{center}

\noindent
3개 지표에서 3가지 방법 모두 $\mathrm{F1}=1.000$ (21/21 영점, 허용 오차 $\delta t < 0.5$). 80자릿수 \texttt{mpmath} 정밀도에서는 곡률 극대만으로 모든 영점을 거짓 양성 없이 식별 가능하다. 이중 기준과 모노드로미 단독 기준은 이 정밀도 수준에서 추가 이점이 없다. 이중 기준의 실용적 가치는 신경망 수준에서 근사 $\kappa$ 추정이 거짓 양성을 유발할 때 드러난다 (추측~3 참조).


\subsection{Conductor 의존성: 위상적 성질 vs.\ 정량적 성질}
\label{app:dirichlet:conductor}

\textcolor{ForestGreen}{[완료 2026-04-15; 5개 함수 교차비교로 확장 2026-04-15]}
다섯 개 GL(1) $L$-함수 ($\zeta$, $\chi_3$, $\chi_4$, $\chi_5$, $\chi_7$)에 걸친 결과를
동일한 방법론 ($t\in[10,40]$, $\delta{=}0.03$ near/far $\kappa$ 비율)으로 비교하면
위상적 성질과 정량적 성질 사이의 체계적 패턴이 드러난다.

\begin{table}[h]
\centering
\caption{다섯 개 GL(1) $L$-함수 교차비교 (동일 방법론,
$t\in[10,40]$, $\delta=0.03$).}
\label{tab:cross-comparison}
\small
\begin{tabular}{lccccc}
\toprule
측정 항목 & $\zeta$ & $\chi_3$ & $\chi_4$ & $\chi_5$ & $\chi_7$ \\
\midrule
Conductor $q$ & 1 & 3 & 4 & 5 & 7 \\
$\tfrac{1}{2}\log(q/\pi)$ & $-0.572$ & $-0.023$ & $+0.121$ & $+0.232$ & $+0.401$ \\
$[10,40]$ 내 영점 수 & 6 & 11 & 10 & 11 & 13 \\
$\kappa$ 집중도 & $1807\times$ & $967\times$ & $801\times$ & $799\times$ & $779\times$ \\
모노드로미 편차 & $0.000$ & $0.286^\dagger$ & $0.000$ & $0.000$ & $0.000$ \\
$\sigma$-유일성 & \checkmark & \checkmark & \checkmark & \checkmark & \checkmark \\
\bottomrule
\end{tabular}
\end{table}

\noindent
$^\dagger$\,영점 탐색기에서 1개 영점 누락으로 인한 수치 인공물.
독립 검증 결과 $\chi_3$는 12/12 모노드로미 PASS 확인
(\S\ref{app:dirichlet:verify}).

\medskip\noindent
\textbf{경험적 관찰.}
세 가지 정성적 패턴이 나타난다:
\begin{enumerate}
  \item \emph{위상적 보편성.}
        모노드로미 양자화 ($\Delta\arg\Lambda = \pm\pi$)와
        $\sigma$-유일성 (위상 점프가 $\sigma=1/2$에서만 발생,
        다른 모든 $\sigma$에서 점프 0)은 다섯 개 $L$-함수 모두에서
        예외 없이 성립한다.
        이들은 단순 영점의 위상적 구조에서 귀결되며 conductor와 무관하다.

  \item \emph{$\kappa$의 예비적 5점 단조 감소.}
        Near/far $\kappa$ 집중도 비율은 $q$가 $1$에서 $7$로 증가함에 따라
        단조 감소한다: $1807 \to 967 \to 801 \to 799 \to 779$.
        Spearman 순위 상관: $r_s = -1.000$ ($p = 1.4\times10^{-24}$,
        완벽한 5점 순위 순서).
        선형 피팅:
        \[
          \kappa \approx -1133 \cdot \tfrac{1}{2}\log(q/\pi) + 1066
          \quad (R^2 = 0.912).
        \]
        이는 이론 예측과 일관적이다: 완성 $L$-함수의 배경항
        $(1/2)\log(q/\pi)$가 $\partial_\sigma\log|\Lambda|$에 기여하여
        $q$가 클수록 off-critical 기저 곡률을 높이고 peak-to-background 비율을
        낮춘다. 단, 감소는 비선형적 ($\zeta\to\chi_3$ 간격 $840\times$ vs
        $\chi_5\to\chi_7$ 간격 $20\times$)으로 포화 구간을 시사한다.
        \textbf{이는 예비적 5점 경향이며, 추가 conductor로의 확인은 향후 과제이다.}

  \item \emph{이전 conductor 표 ($q=3,4,5,7$).}
        모노드로미 양자화와 블라인드 영점 예측 (F1$=1.0$)은 기존 네 conductor
        모두에서 예외 없이 성립한다 (영점 수: 각각 12, 13, 13, 136개).
\end{enumerate}

\noindent
이 관찰은 경험적이며, conductor 감소율에 대한 이론적 한계는 주장하지 않는다.
$\zeta$와 $\chi_{3,4,5,7}$의 near/far $\kappa$ 비율은 동일한 방법론 하에 계산되어
이 표 내에서 직접 비교 가능하다 (이전 $q=7$ 분석의 $\sigma$-비교 방법과는 구별됨).
\textcolor{ForestGreen}{[경험적 관찰; 예비적 5점 경향; conductor 스케일링의 증명 아님]}


% ============================================================================
\section{GL(1)을 넘어서: 타원곡선 $L$-함수 및 라마누잔 $\Delta$ (GL(2))}
\label{app:gl2}
% ============================================================================

앞선 모든 결과는 GL(1) $L$-함수에 관한 것이다: 리만 $\zeta$-함수와
다양한 지표 $\chi$에 대한 디리클레 $L(s,\chi)$.
자연스러운 질문은 $\xi$-다발 프레임워크가 \emph{고차} $L$-함수로
확장되는지이다. 세 가지 GL(2) 사례로 검증한다:
(i)~타원곡선 $11\text{a}1$ ($y^2 + y = x^3 - x^2 - 10x - 20$,
conductor $N = 11$, rank 0, $\varepsilon = +1$),
(ii)~$37\text{a}1$ ($y^2 + y = x^3 - x$, conductor $N = 37$,
rank 1, $\varepsilon = -1$), 그리고
(iii)~라마누잔 $\Delta$ 함수
($\Delta(z) = q\prod_{n\geq 1}(1-q^n)^{24}$, level $N = 1$,
weight $k = 12$, $\varepsilon = +1$).
처음 두 곡선은 weight-2 모듈 형식(타원곡선)이고, 세 번째는 weight 12로
임계선 $\sigma = k/2 = 6$ ($\sigma = 1$인 타원곡선과 대비)을 가진다.
이 세 사례는 conductor $1, 11, 37$; weight $2, 12$; rank $0, 1$;
근 수 $\pm 1$에 걸쳐 있다.

\medskip\noindent
\textbf{GL(1)과의 핵심 차이.}
\begin{itemize}
  \item \emph{임계선}: $\operatorname{Re}(s) = 1$ ($\tfrac{1}{2}$가 아님).
  \item \emph{완비 $L$-함수}: $\Lambda(f,s) = (\sqrt{N}/2\pi)^s\,\Gamma(s)\,L(f,s)$,
        함수 방정식 $\Lambda(f,k-s) = \varepsilon\,\Lambda(f,s)$.
  \item \emph{계수}: 타원곡선의 경우 $a_p = p + 1 - \#E(\mathbb{F}_p)$;
        $\Delta$의 경우 $\tau(n)$ (라마누잔 타우), $|\tau(p)| \leq 2p^{11/2}$ (들리뉴).
  \item \emph{평가}: 상위 불완전 감마 함수를 이용한 근사 함수 방정식(AFE) ---
        직접 디리클레 급수는 임계선에서 조건부 수렴.
\end{itemize}

\medskip\noindent
\textbf{검증.}
11a1: $L(E,1) = 0.2538418609\ldots$는 LMFDB와 10자리까지 일치;
$t \in [0.5, 30]$에서 17개 영점; $\gamma_1 = 6.3626\ldots$은 LMFDB 일치.
37a1: $L(E,1) = 0$ ($\varepsilon = -1$에 의한 강제 영점, rank 1 확인);
$\Lambda(E,2{-}s) = -\Lambda(E,s)$ 4개 검증점에서 기계 정밀도 성립;
$t \in [0, 30]$에서 강제 영점 포함 24개 영점.
라마누잔 $\Delta$: $\tau(n)$ 계수가 $n = 1, \ldots, 20$ 모두에서 LMFDB와 일치;
헤케 점화식 $\tau(p^2) = \tau(p)^2 - p^{11}$이 4개 소수에서 검증;
$\Lambda(\Delta, 12{-}s) = \Lambda(\Delta, s)$가 4개 검증점에서 기계 정밀도 성립;
$\gamma_1 = 9.2224\ldots$는 LMFDB 일치; $L(\Delta, 6) = 0.7921\ldots$;
$t \in [0.5, 30]$에서 8개 영점.
세 함수 방정식 모두 완전한 수치 정밀도로 성립한다.

\medskip\noindent
\textbf{4성질 검증.}

\begin{center}
\small
\begin{tabular}{llllll}
\toprule
성질 & 11a1 (wt 2) & 37a1 (wt 2) & $\Delta$ (wt 12) & 핵심 지표 & 판정 \\
\midrule
$\sigma$-유일성 & 17 점프 $\forall\sigma$ & 23 점프 $\forall\sigma$ & 8 점프 $\forall\sigma$ & 구별 불가 & N/A$^\ddagger$ \\
모노드로미 & 17/17 & 23/23 & 8/8 & $|\text{mono}|/\pi = 2.0$ & PASS \\
$\kappa$ 집중도 & $972\times$ & $2684\times$ & $1396\times$ & 중앙값 근거리/원거리 & PASS \\
블라인드 예측 & 15/15 (100\%) & 20/20 (100\%) & 7/7 (100\%) & tol $= 0.5$ & PASS \\
\midrule
\textbf{점수} & \textbf{3/4} & \textbf{3/4} & \textbf{3/4} & & \\
\bottomrule
\end{tabular}
\end{center}

\noindent
$^\ddagger$\,$\sigma$-유일성 검증($\operatorname{Re}(\Lambda)$ 부호 변화 카운트)은
GL(2)에 적용 불가: $\Gamma(s)$ 인자(GL(1)의 $\Gamma(s/2)$ 대비)가 모든 $\sigma$에서
영점 밀도와 동기화되는 진동을 생성하여 동일한 점프 수를 산출한다.
이는 프레임워크의 실패가 아닌 \emph{GL(2)의 구조적 성질}이다.
대안적 $\sigma$-판별법(예: $|\Lambda(\sigma + it)|$ 최솟값 기반)은 미해결.

\medskip\noindent
\textbf{결과 (\#39--\#41).}
$\xi$-다발의 4가지 성질 중 3가지 --- 모노드로미 양자화, $\kappa$ 집중,
블라인드 영점 예측 --- 가 GL(2) $L$-함수로 GL(1)에 필적하는
정량적 강도로 확장된다. 세 사례의 weight (2 vs.\ 12), rank
(0 vs.\ 1), 근 수 ($+1$ vs.\ $-1$), conductor (1, 11, 37),
임계선 ($\sigma = 1$ vs.\ $\sigma = 6$)이 모두 다름에도 불구하고
패턴이 \emph{동일}하다: 각각 3/4 PASS, $\kappa$ 비율
$972\times$, $2684\times$, $1396\times$.
rank-1 곡선 37a1의 $s = 1$ 강제 영점($\varepsilon = -1$에 기인)은
완벽한 모노드로미 양자화를 나타내어, 프레임워크가 구조적 영점을
올바르게 처리함을 확인한다.
라마누잔 $\Delta$ 결과는 특히 주목할 만하다: weight 12는 완전히 다른
$\Gamma$-인자 정칙권역 ($s \approx 1$ 대신 $s \approx 6$ 근방의 $\Gamma(s)$)과
계수 성장 $|\tau(p)| \sim p^{11/2}$를 수반하지만, 3/4 패턴이 변함없이 지속된다.
이는 $\xi$-다발 프레임워크가 \emph{차수 1 $L$-함수에 국한되지 않으며},
GL(2) 보편성이 weight, conductor, rank, 근 수에 무관함을 보이는 수치적 증거이다.
\textcolor{ForestGreen}{[확립, 2026-04-15; 세 GL(2) 사례에서 4성질 중 3개 PASS
(weight 2 및 12); $\sigma$-유일성은 GL(2)에 구조적으로 부적용]}

\subsubsection{$\mathrm{GL}(2)$에서 $\sigma$-유일성 실패 메커니즘}
\label{sec:sigma_uniqueness_mechanism}

위 표의 $\sigma$-유일성 N/A 판정에는 구조적 설명이 필요하다.
$\sigma$-유일성 검증은 $\sigma$가 변할 때 $\operatorname{Re}(\Lambda(\sigma+it))$의
부호 변화를 세는 방법이며, GL(1)에서는 $\Gamma(s/2)$ 인자가 영점 밀도보다
\emph{느리게} 진동하기 때문에 임계선 $\sigma = \tfrac{1}{2}$을 구별할 수 있다.
GL(2)의 경우 $\Gamma(s)$ 인자(인수가 절반이 되지 않음)는 영점 밀도와
정확히 동일한 속도로 진동하므로 판별 대조가 사라진다.

\medskip\noindent
\textbf{Stirling 분석.}
$f_\Gamma(t)$를 $\Gamma$-인자 기여의 단위 $t$당 진동 주파수(부호 변화 횟수),
$f_\mathrm{zero}(t)$를 영점 밀도라 하면 Stirling 근사에 의해:
\begin{equation}
  f_\Gamma^{(1)}(t) = \frac{1}{2\pi}\log\frac{t}{4\pi}, \quad
  f_\mathrm{zero}^{(1)}(t) = \frac{1}{\pi}\log\frac{t}{2\pi}, \quad
  R_{\mathrm{GL}(1)} = \frac{f_\Gamma^{(1)}}{f_\mathrm{zero}^{(1)}} \approx 0.4.
  \label{eq:stirling_ratio_gl1_ko}
\end{equation}
GL(2), $N = 1$ (라마누잔 $\Delta$ 등)에서는:
\begin{equation}
  f_\Gamma^{(2)}(t) = \frac{1}{\pi}\log\frac{t}{2\pi} = f_\mathrm{zero}^{(2)}(t), \quad
  R_{\Delta} = 1.000 \text{ (수학적 항등식).}
  \label{eq:stirling_ratio_gl2_ko}
\end{equation}
이것은 수치적 우연이 아닌 수학적 항등식이다:
완비 라마누잔 $\Delta$ 함수 $\Lambda(\Delta, s) = (2\pi)^{-s}\Gamma(s)L(\Delta,s)$는
하나의 $\Gamma(s)$ 인자를 포함하며, 이 인자의 진동률이
Hadamard 곱 공식에 의해 영점 밀도 $\tfrac{1}{\pi}\log\tfrac{t}{2\pi}$와 일치한다.
$t \in [0.5, 30]$ 범위의 $\Delta$ 경험적 비율은
$N_\mathrm{off}/N_\mathrm{crit} = 8/8 = 1.000$으로 이론 예측과 정확히 일치한다.

\medskip\noindent
\textbf{GL(1) 정성적 확인.}
디리클레 지표 $\chi_3$ ($\varepsilon = -1$, $\operatorname{Im}(\Lambda)$ 사용)에서
경험적 비율은 $N_\mathrm{off}/N_\mathrm{crit} = 8/21 = 0.38$이며,
GL(1) Stirling 예측 $R_{\mathrm{GL}(1)} \approx 0.37$과
\emph{예시적 일치(illustrative agreement)}를 보인다.
이것은 정량적 매칭이 아닌 정성적 일관성으로 이해해야 한다:
$\chi$-특유의 Stirling 보정은 도체 $q$에 의존하며,
$\chi_3$에서의 일치는 부분적으로 우연일 수 있다.
나머지 4개 GL(1) 사례($\zeta$, $\chi_4$, $\chi_5$, $\chi_7$)는
제한된 $t$ 범위 또는 $\operatorname{Re}(\Lambda)$ 모호성으로 인해
이 측정법으로 확인되지 않으나, 원래 $\sigma$-유일성 결과(\#7, $385\times$ 에너지
비율)는 $\kappa$-곡률 관측량 기반이므로 영향을 받지 않는다.

\medskip\noindent
\textbf{차수-보편적 진단 대 차수-1-특이적 진단.}
일반 $\mathrm{GL}(n)$에서 완비 $L$-함수의 $n$개 Gamma 인자는
Hadamard 곱에 의해 영점 밀도와 일치하는 진동률을 기여하므로,
$n \geq 2$인 모든 $n$에 대해 $R_n = 1$이다.
이로부터 다음 계층 구조가 설명된다:
\begin{itemize}
  \item \textbf{차수-보편적}: 모노드로미 양자화, $\kappa$-집중도, 블라인드 영점 예측.
        이들은 GL(1), GL(2)에서 성립하며 모든 차수에서 성립할 것으로 예측된다.
  \item \textbf{차수-1-특이적}: $\sigma$-유일성.
        $R < 1$의 격차는 $\Gamma(s/2)$가 영점 밀도의 절반 속도로 진동하는
        GL(1)에서만 존재한다.
\end{itemize}
GL(2) 도체 $N > 1$ (예: 11a1의 $R \approx 0.16$, 37a1의 $R \approx 0.18$)에서는
형식적 Stirling 비율 $R < 1$이지만, 실제로는 $\Gamma(s)$ 진동 구조가
임계선 안팎에서 부호 변화를 동등하게 만들어 $\sigma$-유일성 FAIL이
검증된 모든 GL(2) 사례에서 유지된다.

\textcolor{ForestGreen}{[결과 \#42, 조건부 양성, 2026-04-16;
$R_\Delta = 1.000$ ($N=1$에 대한 정확한 항등식);
$\chi_3$ 경험 $0.38 \approx 0.37$ (예시적 일치);
결과 \#7의 GL(1) $\sigma$-유일성 ($385\times$) 영향 없음]}

\subsubsection{$t$-스케일링 안정성: 고 $t$에서의 GL(2) (결과 \#49)}
\label{sec:gl2_high_t}

GL(2) 다발 진단이 큰 허수부에서 안정하게 유지되는지 검증하기 위해 --- GL(1) 고
$t$ 강건성 결과 \#31(\cref{obs:high_t_robustness})에 대응하여 --- 타원곡선
$11\text{a}1$ $L$-함수에서 세 높이 구간에 걸쳐 $\kappa$, 모노드로미,
$\sigma$-국소화를 측정하였다: 구간 1 ($t\in[48,55]$), 구간 2 ($t\in[95,105]$),
구간 3 ($t\in[140,155]$). 적응적 정밀도 $\texttt{dps}=\max(60,\,0.7|t|+30)$를
사용하는 AFE를 적용하였다. 저 $t$ 기준선 $\kappa_0 = 1115.0$은 결과 \#45에서
$t\approx 6.4, 13.6, 20.4$의 영점 세 개의 평균이다.

\medskip\noindent\textbf{결과.}
\begin{itemize}
  \item \textbf{구간 1 ($t\approx 51$):} 영점 7개 발견, 전부 $\sigma=1.000000$;
        $\kappa\approx 1118$--$1125$ ($1.00$--$1.01\times$ 기준선);
        mono$/\pi = 2.000$; peak$_\sigma=1.000$. \textbf{통과.}
  \item \textbf{구간 2 ($t\approx 100$):} 영점 11개 발견, 전부 $\sigma=1.000000$;
        $\kappa_{\mathrm{median}}=1190.7$ ($1.07\times$ 기준선);
        mono$/\pi = 4.000$ (수치 아티팩트 추정; 하단 참조); peak$_\sigma=1.000$.
        \textbf{통과} (완만한 $\kappa$ 상승, 해석적 예측 범위 내).
  \item \textbf{구간 3 ($t\approx 147$):} 영점 0개, $\kappa\equiv 0$.
        \textbf{수치 정밀도 장벽}: $t=147$에서 완비 $L$-함수
        $|\Lambda(f,1+it)|\sim e^{-\pi t/2}\approx 10^{-100}$이
        $\texttt{dps}=132$의 상쇄 예산을 초과한다.
        이것은 \emph{실용적 구현 한계}이며, 프레임워크의 이론적 실패가 아니다;
        $\texttt{dps}\geq 200$으로 원칙적으로 해결 가능하다.
\end{itemize}

\medskip\noindent\textbf{GL(1) \#31 대 GL(2) \#49 교차 순위 비교.}

\begin{center}
\small
\begin{tabular}{@{}l c c c@{}}
\toprule
지표 & GL(1) $\zeta$, $t\sim14$--$10\,000$ (\#31) & GL(2) 11a1, $t\sim6$--$30$ (\#45) & GL(2) 11a1, $t\sim51$--$100$ (\#49) \\
\midrule
$\kappa$ & $1112$--$1119$ & $1115.0$ & $1118$--$1191$ \\
$\kappa/\kappa_0$ & $\leq 0.6\%$ & 기준선 & $\leq 7\%$ \\
mono$/\pi$ & $2.000$ & $2.000$ & $2.000$ (3구간 중 2개) \\
peak$_\sigma$ & $0.497$ ($\sigma_c=0.5$) & $1.000$ ($\sigma_c=1$) & $1.000$ \\
\bottomrule
\end{tabular}
\end{center}

GL(2) $\kappa$ 값은 GL(1) 기준선과 동일한 수치 범위($\approx 1112$--$1191$)를
유지하며, 곡률 스케일이 다발 프레임워크의 \emph{차수-독립적} 성질임을 확인한다.
$t\approx 100$에서의 완만한 7\% 상승은 인접 영점 간섭(이 높이에서 평균 간격
$\approx 0.77$)에 기인하며, 해석적 선도항 $\kappa\approx 1/\delta^2 + O(1/\Delta t)$와
일치한다.

\medskip\noindent\textbf{$t\approx 100$에서 mono$/\pi=4.0$ 해석.}
$t\approx 100$에서 권선수 2는 수치적 아티팩트로 추정된다:
최근접 영점 $\gamma=99.228$과 $\gamma=100.001$의 간격 $0.773$은
반경 $r=0.3$의 2배를 초과하므로 하나($\gamma=100.001$)만 윤곽 안에 포함된다.
결과 \#54 (37a1, $t=29.0$)에서는 $\gamma=29.03,\,29.28$ 간격 $0.25 < 2r$로
두 영점 포획이 정당했으나, 여기서는 $|\Lambda|\sim 10^{-68}$에서의 폐곡선 적분 시
누적 위상 오차가 추가 권선을 유발한 것으로 보인다; 동일 코드가 저 $t$에서는
정확히 mono$/\pi=2.000$을 산출한다.

\textcolor{ForestGreen}{[결과 \#49, 조건부 양성, 2026-04-16;
GL(2) 11a1; 구간 1--2 통과 ($t\approx51$--$100$);
구간 3 수치 정밀도 장벽 ($t\approx147$, 이론적 한계 아님);
수치 검증, 증명 아님]}

% ============================================================================
\subsection{Maass 형식 확장 (비정칙, Weight~0)}
\label{sec:maass}
% ============================================================================

앞선 GL(2) 결과는 모두 \emph{정칙} 모듈러 형식(weight~2 타원곡선 및
weight~12 라마누잔 $\Delta$)에 관한 것이다. $\xi$-다발 프레임워크가
weight~0, 연속 스펙트럼 파라미터~$R$, 질적으로 다른 감마 인자를 갖는
\emph{비정칙} Maass 첨두 형식으로 확장되는지가 자연스러운 질문이다.

$\mathrm{SL}(2,\mathbb{Z})$ 위의 첫 번째 Maass 첨두 형식
(스펙트럼 파라미터 $R = 9.53369526\ldots$, 홀수 패리티,
$\varepsilon = +1$, 수준 $N = 1$)을 테스트한다.
푸리에 계수 (455항, 1000자리 정밀도)는 Booker--Str\"ombergsson--Venkatesh
데이터셋에서 취한다.

\medskip\noindent
\textbf{정칙 GL(2)와의 주요 차이.}
\begin{itemize}
  \item \emph{감마 인자}: $\Lambda(s) = \pi^{-(s+1)}\,
        \Gamma\!\bigl(\tfrac{s+1+iR}{2}\bigr)\,
        \Gamma\!\bigl(\tfrac{s+1-iR}{2}\bigr)\,L(s)$
        (홀수의 경우, $+1$ 이동이 짝수와 구분).
  \item \emph{임계선}: $\sigma = 1/2$, $\zeta(s)$와 동일.
  \item \emph{Weight}: $k = 0$ (비정칙), 위 정칙 경우의 $k = 2$ 또는 $12$와 대조.
  \item \emph{수준}: $N = 1$ (전 모듈러 군), 가장 단순한 conductor.
\end{itemize}

\medskip\noindent
\textbf{검증.}
Hecke 곱셈성 정확 확인:
$a(6) = a(2)\cdot a(3)$, $a(4) = a(2)^2 - 1$,
$a(9) = a(3)^2 - 1$, 모두 $\mathrm{diff} = 0$.
함수방정식 $\Lambda(s) = \Lambda(1{-}s)$가 4개 테스트 점
($t = 5, 10, 15, 25$; 최대 상대 오차 $= 0.00$)에서 완전 수치 정밀도로 성립.

\medskip\noindent
\textbf{4성질 결과.}

\begin{center}
\small
\begin{tabular}{@{}l l l l@{}}
\toprule
성질 & 결과 & 핵심 지표 & 판정 \\
\midrule
함수방정식 & $\mathrm{rel} = 0.00$ (4점) & 최대 $|\Lambda|{=}1.2{\times}10^{-7}$ & 통과 \\
영점 발견 & 18개 ($t \in [2, 50]$) & $\gamma_1 = 3.5988$ & 통과 \\
$\kappa_{\mathrm{near}}$ & $1114.38$ (8~TP, CV${=}0.1\%$) & 비율 $577.9\times$ & 통과 \\
모노드로미 & 8/8 (mono$/\pi = 2.0000$) & 완전 양자화 & 통과 \\
$\sigma$-유일성 & 3/3 통과 ($\sigma{=}0.5$ 최대) & $\kappa(0.5)/\kappa(0.45) \approx 10^{11}$ & 통과 \\
\bottomrule
\end{tabular}
\end{center}

\medskip\noindent
\textbf{차수별 $\kappa_{\mathrm{near}}$ 비교.}
첫 번째 Maass 형식 (홀수, $R{=}9.53$)은 $\kappa_{\mathrm{near}} = 1114.38 \pm 1.1$
(CV${=}0.1\%$, 8~TP)을 산출한다.  두 번째 Maass 형식 (짝수, $R{=}13.78$; 24영점, 12~TP)은
$\kappa_{\mathrm{near}} = 1115.40 \pm 2.2$ (CV${=}0.2\%$)로,
$0.09\%$ 이내에서 $\kappa_{\mathrm{near}}$ 재현성을 확인한다.
라마누잔 $\Delta$ $L$-함수 ($N{=}1$, weight${=}12$; 결과~\#58)는
$\kappa_{\mathrm{near}} = 1114.13 \pm 1.1$ (CV${=}0.1\%$, 12~TP)을 산출하여,
Maass 평균과 $0.07\%$ 이내로 일치한다---$\kappa_{\mathrm{near}}$가
weight, 패리티, 스펙트럼 파라미터에 독립임을 확인한다.
세 형식 (Maass 2개 $+$ $\Delta$)의 결합 GL(2) 추정치는
$\kappa_{\mathrm{near}}(2) \approx 1114.6$ (34~TP)이며,
GL(3) 대칭 제곱 $L$-함수는 $1125.16 \pm 1.7$ (CV${=}0.15\%$)이다.
GL(2)--GL(3) 차이 ${\sim}10$ ($0.94\%$, ${\sim}6\sigma$)는
$\kappa_{\mathrm{near}}$가 단일 보편상수가 아닌
\emph{차수-의존 상수}임의 강한 수치적 증거를 제공한다.

결과~\#59는 이 스펙트럼의 $d{=}1$ 끝점을 고정한다.
리만 제타함수 $\zeta(s)$는 $\kappa_{\mathrm{near}}(1) = 1112.32$
(CV${=}0.1\%$, $t \in [10,60]$ 내 13~TP)를 산출하여,
측정된 세 차수 전체에 걸친 단조 패턴을 확립한다:

\smallskip
\begin{center}
\small
\begin{tabular}{clccl}
\toprule
차수 $d$ & $L$-함수 & $\kappa_{\mathrm{near}}(d)$ & $n$(TP) & CV \\
\midrule
1 & $\zeta(s)$ & $1112.32$ & 13 & $0.1\%$ \\
2 & GL(2) 평균 (Maass 2개 $+$ $\Delta$) & $1114.6$ & 34 & ${\sim}0.1\%$ \\
3 & GL(3) sym$^2$ 평균 (6곡선) & $1125.16$ & ${\sim}60$ & $0.15\%$ \\
\bottomrule
\end{tabular}
\end{center}
\smallskip

\noindent
간격은 $\Delta_{1\to2} = 2.28$, $\Delta_{2\to3} = 10.56$으로 비율 $4.63\times$이며,
등차 증가가 아닌 \emph{가속 증가} 패턴을 시사한다.
이는 선형 모델 $\kappa_{\mathrm{near}}(d) = a + bd$를 기각하며,
함수 형태(이차, 지수, 또는 다른 형태)를 판별하려면 $d{=}4$의
네 번째 측정값이 필요하다 (실험~\#72 진행 중).

\medskip\noindent
\textbf{차수 의존성의 이론적 근거 (경계 B-12 해결).}\;
명제~\ref{prop:kappa_exact_expansion}는 정밀한 구조적 설명을 제공한다.
식~\eqref{eq:kappa_exact}에 의해, 임의의 $L$-함수의 임계선 단순 영점
$s_0 = \tfrac12+it_0$에서
\[
  \kappa(s_0 + \delta) = \frac{1}{\delta^2} + A(t_0,\mu) + O(\delta^2),
\]
여기서 선도항 $1/\delta^2$는 모든 차수·계열에서 \emph{보편적}이며,
$A(t_0,\mu) = B(t_0,\mu)^2 + 2H_1(t_0,\mu)$는 영점의 국소 해석적 환경,
특히 $\Gamma$-인자 파라미터 $\mu$에 의존한다.
고정 $\delta = 0.03$에서 측정된 $\kappa_{\mathrm{near}}$ 값은
\[
  \kappa_{\mathrm{near}}(d) \approx \frac{1}{\delta^2} + \langle A(t_0,\mu_d)\rangle,
\]
차수 간 차이($d{=}1{\to}2$: $2.28$; $d{=}2{\to}3$: $10.56$)는
$1/\delta^2$ 항이 아닌 오직 $\langle A\rangle$의 $\Gamma$-인자 구조에서 발생한다.
$A(t_0,\mu)$는 digamma 기여 $H_1 \propto \tfrac12\psi'(\tfrac14+\mu/2)$를 통해
$\mu_{\max}$가 클수록 증가한다. $d{=}4$ (sym$^3\Delta$, 실험~\#72)의 비단조 거동도
이와 일관된다: sym$^3\Delta$는 $\mu_{\max}{=}12 <$ sym$^2\Delta$의 $\mu_{\max}{=}22$이므로
$\kappa_{\mathrm{near}}(4) < \kappa_{\mathrm{near}}(3)$이 자연스럽다.
\emph{경계 B-12는 완전히 해소된다}: $\kappa_{\mathrm{near}}(d)$는 차수만이 아니라
$L$-함수 계열의 $\Gamma$-인자 복잡도를 측정한다.

또한 $\kappa$ near/far 비율은 차수에 따라 단조 감소한다:

\smallskip
\begin{center}
\small
\begin{tabular}{clc}
\toprule
차수 $d$ & $L$-함수 & $\kappa$ 비율 (near/far) \\
\midrule
1 & $\zeta(s)$ & $2200.7\times$ \\
2 & GL(2) 평균 (Maass, $\Delta$) & ${\sim}600\times$ \\
3 & GL(3) sym$^2$ & ${\sim}320\times$ \\
\bottomrule
\end{tabular}
\end{center}
\smallskip

\noindent
near/far $\kappa$ 비율의 차수-역비례 패턴 (차수 높을수록 영점/비영점 분리도 낮음)은
높은 차수에서 감마 인자 구조가 복잡해짐에 따라 비영점에서의
기준 곡률이 상승하는 것과 일관된다.

\medskip\noindent
\emph{저 $\gamma$ 곡률 편향.}\;
GL(2) 형식에서 $\gamma < 20$ 영역의 하방 편향(${\sim}0.3\%$)이 관측된 바 있다.
GL(1) $\zeta(s)$에서는 이 편향이 $0.07\%$에 불과하여
($\gamma < 20$: $\kappa_{\mathrm{near}} = 1111.58$, $n{=}1$; $\gamma \geq 20$: $1112.36$, $n{=}12$),
저 $\gamma$ 편향이 차수-의존적으로 감소함을 시사한다.

\medskip\noindent
\textbf{$\sigma$-유일성: conductor가 통과/실패를 결정 (경계 B-10 해결).}
$N > 1$인 정칙 GL(2) 형식(11a1, 37a1: 실패)과 달리,
$N{=}1$인 모든 형식---Maass와 정칙 모두---$\sigma = 0.5$에서 극도로 날카로운 피크로
$\sigma$-유일성을 통과한다.
두 Maass 형식은 각각 3/3과 24/24 PASS를 달성하며,
라마누잔 $\Delta$ ($N{=}1$, weight${=}12$; 결과~\#58)는
$\kappa(0.5)/\kappa(0.45) \in [10^{13}, 10^{18}]$로 23/23 PASS를 달성한다.
차수 2 $L$-함수 간 통과/실패 분기를 설명하는 두 가설이 제안되었다:
(A)~conductor 의존성 ($N{=}1 \to$ 통과, $N{>}1 \to$ 실패), 또는
(B)~weight 의존성 (weight${=}0 \to$ 통과, weight${\geq}2 \to$ 실패).
$\Delta$는 $N{=}1$이면서 weight${=}12$이나 $\sigma$-유일성을 통과하므로,
가설~(B)는 기각되고 가설~(A)가 확정된다:
\emph{$\sigma$-유일성 통과/실패는 conductor $N$이 1인지 여부에 의해 결정된다.}
결과~\#59 ($\zeta(s)$, 차수~1)는 이를 더욱 강화한다:
13/13 통과, $\kappa(0.5)/\kappa(0.45) \in [10^{26}, 10^{27}]$로
전체 테스트된 $L$-함수 중 가장 날카로운 $\sigma$-유일성 피크를 기록한다.
완전한 증거는 다음과 같다:

\smallskip
\begin{center}
\small
\begin{tabular}{llcccl}
\toprule
$L$-함수 & 차수 & $N$ & weight & $\sigma$-유일성 & $n$(영점) \\
\midrule
$\zeta(s)$ & 1 & 1 & --- & 통과 & 385 역사적; 4성질: 13/13 \\
Maass $R{=}9.53$ (홀수) & 2 & 1 & 0 & 통과 & 3 \\
Maass $R{=}13.78$ (짝수) & 2 & 1 & 0 & 통과 & 24 \\
라마누잔 $\Delta$ ($w{=}12$) & 2 & 1 & 12 & 통과 & 23 \\
\midrule
11a1 & 2 & 11 & 2 & 실패 & --- \\
37a1 & 2 & 37 & 2 & 실패 & --- \\
sym$^2$(11a1) & 3 & 121 & --- & 실패 & --- \\
\bottomrule
\end{tabular}
\end{center}
\smallskip

\noindent
패턴은 명확하다: \emph{다섯 개의 $N{=}1$ $L$-함수} 전부 (차수~1--2, weight~0, 2, 12,
정칙·비정칙 형식 모두 포괄)가 통과하고,
세 개의 $N{>}1$ $L$-함수는 전부 실패한다.
총계: $N{=}1 \Leftrightarrow$ 통과, 63개 영점에 걸쳐 예외 없이 확정.
Weight는 무관하다 (weight 0과 12 모두 통과); 차수도 무관하다
($N{=}1$ 하에서 차수 1과 2 모두 통과).
이는 conductor $N{=}1$이 $\sigma$-유일성의 유일한 결정 요인이라는
강한 수치적 증거를 제공하나, 이 이분법의 이론적 메커니즘은 열린 문제로 남는다.

\textcolor{ForestGreen}{[결과 \#56--\#62, 확립, 2026-04-17/18;
GL(2) Maass 형식: 홀수 ($R{=}9.53$) + 짝수 ($R{=}13.78$, $N{=}1$);
라마누잔 $\Delta$ ($N{=}1$, $w{=}12$): 23영점, $\sigma$-유일성 23/23 통과;
$N{=}1 \Leftrightarrow$ 통과 확정 (경계 B-10 해결);
GL(1) $\zeta(s)$: 13영점, $\kappa_{\mathrm{near}}{=}1112.32$ (CV${=}0.1\%$), $\sigma$-유일성 13/13 통과, $\kappa$ 비율 $2200.7\times$;
$\kappa_{\mathrm{near}}(d)$ 단조증가: $1112.32$ ($d{=}1$) $<$ $1114.6$ ($d{=}2$) $<$ $1125.16$ ($d{=}3$), 간격 가속 $4.63\times$;
$\kappa$-$\delta$ 스케일링: 65점, $\kappa{\cdot}\delta^2{\in}[1.00012,1.01206]$, $O(1/\delta)$ 계수 ${=}0$, B-12 해소 (명제~\ref{prop:kappa_exact_expansion});
수치 검증, 증명 아님]}


% ============================================================================
\section{GL(2)를 넘어서: GL(3) 대칭 제곱 $L$-함수}
\label{app:gl3}
% ============================================================================

이제 $\xi$-다발 프레임워크를 \emph{차수 3} $L$-함수로 확장한다.
$\mathbb{Q}$ 위의 conductor $N_E$인 타원곡선 $E$에 대해, 대칭 제곱
$L$-함수 $L(\mathrm{sym}^2 E, s)$는 conductor $N = N_E^2$, 근 수
$\varepsilon = +1$, 해석적 정규화에서 감마 이동 $\mu = (1, 1, 2)$를
가지는 GL(3) 보형 $L$-함수이다. 완비 $L$-함수는
$\Lambda(\mathrm{sym}^2 E, s) = \Lambda(\mathrm{sym}^2 E, 1-s)$를
만족하며, 임계선은 $\sigma = 1/2$이다. 디리클레 계수는 좋은 소수 $p$에
대해 $c(p) = a_p^2/p - 1$이다.

GL(2)와의 핵심 차이: (i) 차수 3 대 2이므로 독립 감마 인자가 하나 대신 둘;
(ii) 근사 함수 방정식(AFE)이 $\Gamma_{\mathbb{R}}$ 인자 3개의 곱을 포함하는
Mellin--Barnes 적분; (iii) $\sigma$-유일성 실패 예상
(\S\ref{subsec:gl3_sigma} 참조).

두 곡선을 검증한다: \textbf{sym$^2$(11a1)} ($N = 121$, rank 0,
소수 11에서 분할 승법) 및 \textbf{sym$^2$(37a1)} ($N = 1369$, rank 1,
소수 37에서 비분할 승법).

\subsection{4성질 검증: sym$^2$(11a1) 및 sym$^2$(37a1)}
\label{subsec:gl3_fourprop}

표~\ref{tab:gl3_fourprop}에 두 곡선의 4성질 검증을 정리한다.

\begin{table}[ht]
\centering
\caption{GL(3) 4성질 검증.\label{tab:gl3_fourprop}}
\smallskip
\begin{tabular}{lcc}
\toprule
성질 & sym$^2$(11a1) & sym$^2$(37a1) \\
\midrule
Conductor $N$ & 121 & 1369 \\
근 수 $\varepsilon$ & $+1$ & $+1$ \\
함수방정식 잔차 & $0.0$ & $0.0$ \\
발견 영점 ($t \in [2, 45]$/$[2, 50]$) & 46 & 44 \\
LMFDB 매칭 & 15/15 & (LMFDB 데이터 없음) \\
$\kappa$ near/far 비율 & $283.3\times$ & $222.1\times$ \\
모노드로미 (mono/$\pi = 2.0$) & 12/12 & 15/15 \\
$\sigma$-유일성 & FAIL & FAIL \\
\bottomrule
\end{tabular}
\end{table}

두 곡선 모두 기계 영점 수준의 잔차로 함수방정식을 통과한다.
$\kappa$ 집중비 ($283\times$, $222\times$)는 GL(1) 및 GL(2)에서 확립된
$100\times$ 임계값을 크게 초과하여, 추측~1이 GL(3)으로 확장됨을 확인한다.
모노드로미 양자화는 모든 검증 영점에서 mono/$\pi = 2.000$을 산출한다.%
\footnote{sym$^2$(11a1)의 경우, $\kappa_{\mathrm{far}}$의 원격 표본이 3점뿐이므로,
$283\times$ 비율은 경계 조건 B-06에 따른 하한으로 간주해야 한다.}

$\sigma$-유일성은 두 곡선 모두 실패한다: $\sigma = 0.5$에서의 부호변환 수가
극대가 아니다 (진단은 \S\ref{subsec:gl3_sigma} 참조). 이는 GL(2) 패턴과
일치하며, 차수 $\geq 2$인 $L$-함수에서 예상되는 결과이다.

\subsection{블라인드 영점 예측과 영점 탐색 기계}
\label{subsec:gl3_blind}

GL(2)와 동일한 $\kappa$-피크 프로토콜로
$\Lambda(\mathrm{sym}^2(11\text{a}1), \tfrac{1}{2}+\delta + it)$를
$t \in [2, 45]$, 해상도 $\delta t = 0.1$, $\delta = 0.03$에서 스캔한다.
$\kappa$-피크 탐지기가 \textbf{36개 피크} (임계값 = 중앙값 $\times$ 10)를
발견한다. 이분법 정밀화와 모노드로미 필터 (36개 모두
mono/$\pi \approx 2.0$ 통과) 후 LMFDB 15개 영점과 비교한다:

\begin{itemize}
  \item \textbf{재현율} = 1.000 (15/15 LMFDB 영점 복원).
  \item \textbf{원시 정밀도} = 0.417 (15/36; 21개 겉보기 ``거짓 양성'').
  \item 위치 오차: 평균 = 0.023, 최대 = 0.073.
\end{itemize}

21개 ``거짓 양성''은 모두 $t > 17.69$ (LMFDB 수록 범위 밖)에 위치한다.
결과~\#67에서 21개 전체를 $\mathrm{dps} = 80$, $N_{\mathrm{coeff}} = 200$
디리클레 계수, 60점 Gauss--Hermite 구적법으로 독립 AFE 검증한다:

\begin{itemize}
  \item \textbf{21/21 모두 실제 영점으로 확인}: 중앙값
        $|\Lambda(1/2+it)| = 2.5 \times 10^{-27}$.
  \item TP 기준선 (LMFDB 15개 영점)의 중앙값 $|\Lambda| = 2.8 \times 10^{-13}$;
        FP 값은 크기가 \emph{더 작으며}, 높은 $t$에서 $|\Lambda|$가
        지수적으로 감쇠하는 것과 일관적이다.
\end{itemize}

21개 FP를 TP로 재분류하면 \textbf{교정 지표}는:
$P = 1.000$, $R = 1.000$, $F_1 = 1.000$. $\xi$-다발 프레임워크는
\emph{영점 탐색 기계}로 기능한다: LMFDB에 수록되지 않은 GL(3) 영점을
독립적으로 발견하며, AFE가 삼중 확인 프로토콜로 각각을 검증한다:
(i)~$\mathrm{Re}\,\Lambda(1/2+it)$의 부호변환 (결과~\#63),
(ii)~mono/$\pi = 2.0$인 $\kappa$-피크 (결과~\#65),
(iii)~$\mathrm{dps} = 80$에서 $|\Lambda(1/2+it)| < 10^{-19}$ (결과~\#67).

\textcolor{ForestGreen}{[결과 \#63, \#65, \#67; 수치 검증, 증명 아님;
sym$^2$(11a1)의 LMFDB 영점은 15개로 제한됨]}

\subsection{$\sigma$-유일성 실패: 디리클레 급수 수렴 기원}
\label{subsec:gl3_sigma}

결과~\#66b는 GL(3)에서 $\sigma$-유일성이 왜 실패하는지 진단하기 위해
5단계 역보정을 적용한다. 다섯 단계는 $\Lambda(s)$의 구성 요소를
점진적으로 제거한다:
\begin{enumerate}
  \item $\Lambda(s)$ (원본);
  \item $\Lambda(s)/|\gamma(s)|$ (감마 진폭 제거);
  \item $\Lambda(s)/\gamma(s) = Q^s L(s)$ (감마 전체 제거);
  \item $L(s)$ (순수 $L$-함수);
  \item $\Lambda(s)/|\Lambda(s)|$ (위상만 추출).
\end{enumerate}
\emph{어떤 단계에서도} $\sigma$-유일성이 복원되지 않는다:
$\sigma = 0.5$에서의 부호변환 수가 어떤 경우에도 극대가 아니다.

\begin{table}[ht]
\centering
\caption{GL(3) sym$^2$(11a1)에 대한 5단계 $\sigma$-유일성 역보정.
         각 $\sigma$에서의 부호변환 수.\label{tab:gl3_sigma}}
\smallskip
\begin{tabular}{cccccc}
\toprule
$\sigma$ & $\Lambda(s)$ & $\Lambda/|\gamma|$ & $\Lambda/\gamma$ & $L(s)$ & 위상 \\
\midrule
0.30 & 25 & 25 & 22 & 29 & 25 \\
0.40 & 23 & 23 & 22 & 29 & 23 \\
0.50 & 23 & 23 & 26 & 19 & 23 \\
0.60 & 23 & 23 & 34 &  3 & 23 \\
0.70 & 25 & 25 & 36 &  1 & 25 \\
0.80 & 25 & 25 & 36 &  1 & 25 \\
0.90 & 27 & 27 & 36 &  1 & 27 \\
\bottomrule
\end{tabular}
\end{table}

이는 감마 인자, conductor 인자, 진폭이 실패의 원인이 아님을 배제한다.
진단은 \emph{디리클레 급수 수렴 구조} 자체를 가리킨다: 차수가 1에서 3으로
증가함에 따라 공명 지표 $R \approx d \cdot \log t / (4\pi)$가 커지고,
부호변환 지형이 평탄해진다.

\medskip\noindent
\textbf{GL(1)/GL(2)/GL(3) $\sigma$-유일성 패턴.}
\begin{itemize}
  \item GL(1) $\zeta$: PASS ($\sigma = 0.5$이 극대).
  \item GL(2) $11\text{a}1$: FAIL ($\sigma = 0.7$--$0.9$ 극대).
  \item GL(3) $\mathrm{sym}^2(11\text{a}1)$: FAIL (거의 평탄; $\sigma = 0.9$ 극대).
  \item GL(3) $\mathrm{sym}^2(37\text{a}1)$: FAIL ($\sigma = 0.9$ 극대).
\end{itemize}
이 패턴은 $\sigma$-유일성이 \emph{GL(1) 고유} 현상이며,
$L$-함수의 차수가 증가함에 따라 부호변환 지형이 점진적으로 평탄해짐을 시사한다.

\textcolor{ForestGreen}{[결과 \#66b, 음성 (진단적); 5단계 역보정;
전 단계 FAIL; 원인 = 디리클레 급수 수렴; 수치 검증, 증명 아님]}

\subsection{$\kappa$-도체 스케일링과 $\kappa_{\mathrm{near}}$ 보편상수}
\label{subsec:gl3_kappa_scaling}

결과~\#68은 4성질 검증을 \textbf{6개} 타원곡선으로 확장하여 $37\times$
도체 범위($N = 121$에서 $4489$)에 걸쳐 $\kappa$ near/far 비율이 도체와
스케일링하는지 검증한다.

\begin{table}[ht]
\centering
\caption{6개 GL(3) sym$^2$ $L$-함수의 $\kappa$-도체 스케일링.
\label{tab:gl3_kappa_scaling}}
\smallskip
\begin{tabular}{lcrrrc}
\toprule
곡선 & $N = N_E^2$ & $\kappa_{\mathrm{near}}$ & $\kappa_{\mathrm{ratio}}$ & mono & FE \\
\midrule
11a1 & 121   & 1123.53 & $283\times$ & 12/12 & 0.0 \\
19a1 & 361   & 1123.04 & $288\times$ & 6/6   & 0.0 \\
37a1 & 1369  & 1127.21 & $222\times$ & 15/15 & 0.0 \\
43a1 & 1849  & 1123.11 & $260\times$ & 6/6   & 0.0 \\
53a1 & 2809  & 1126.79 & $467\times$ & 6/6   & 0.0 \\
67a1 & 4489  & 1149.53 & $217\times$ & 3/6\,${}^\dagger$  & 0.0 \\
\bottomrule
\end{tabular}

\medskip\noindent
{\footnotesize ${}^\dagger$\,67a1: 세 TP에서 인접 영점의 이중포획(간격 $< 2r$)으로 mono/$\pi = 4.0$; mono/$\pi = 2.0$ TP의 $\kappa_{\mathrm{near}}$
평균은 $1130.8$으로, 5곡선 평균과 $0.5\%$ 차이.}
\end{table}

\paragraph{음성 결과: 도체 스케일링 부재.}
$\kappa_{\mathrm{ratio}}$ 대 $N$의 로그-로그 회귀는
$\alpha = 0.003 \pm 0.102$, $R^2 = 0.0002$; Spearman $r_s = -0.31$
($p = 0.54$). 상수, 멱법칙, 로그 모형이 통계적으로 구별 불가능하다
(잔차 SS $\approx 42{,}000$). $\kappa_{\mathrm{ratio}}$는 도체와 스케일링하지
\emph{않으며}, 변동은 원격 표본 $\kappa_{\mathrm{far}}$의 우연적 영점 근접에 기인한다.

\paragraph{양성 결과: $\kappa_{\mathrm{near}} \approx 1125$는 도체 독립 상수.}
67a1 이상치(이중포획 아티팩트)를 제외한 5곡선의 $\kappa_{\mathrm{near}}$
중앙값은 $\mathbf{1124.74}$ (CV $= 0.15\%$). 6곡선 전체는 $1125.16$
(CV $= 0.8\%$). 이는 $\kappa_{\mathrm{near}}$가 개별 $L$-함수의 산술적
복잡도가 아닌 \emph{프레임워크 기하학}(차수, 감마 이동, 윤곽 매개변수)에
의해 결정됨을 시사한다.

\textcolor{ForestGreen}{[결과 \#68; 스케일링 음성 ($R^2{=}0.0002$);
$\kappa_{\mathrm{near}}$ 보편상수 확립 (CV${=}0.15\%$, 5곡선);
수치적 증거, 증명 아님]}


\section{음성 결과: Gram 점 곡률}
\label{app:gram_negative}

Gram 점 $g_n$ ($\vartheta(g_n) = n\pi$로 정의)에서의 곡률 $\kappa$가
Gram 법칙 위반의 서명을 담고 있는지는 자연스러운 질문이다.
$n = 0, \ldots, 200$에 대해 $\kappa(g_n + 0.03)$을 계산하고
각 Gram 점을 ``good'' (Gram 법칙 만족) 또는 ``bad'' (위반)으로 분류하였다.
이 범위에서 201개 Gram 점 중 3개만 bad ($n \in \{126, 134, 195\}$)이며,
순위 기반 검정(Mann--Whitney, 그룹당 $n \geq 8$ 필요)에
충분한 통계적 검정력을 제공하지 않는다. 따라서 이 실험은
\textbf{음성}이다 --- 가설이 반증된 것이 아니라, 데이터가 검정에
부적합하기 때문이다. bad Gram 점이 더 빈번해지는 $n > 1000$
범위로 확장이 필요하다.
\textcolor{ForestGreen}{[음성, 2026-04-14]}


\section{표기법 색인}
\label{app:notation}

\begin{center}
\small
\begin{tabular}{lll}
\toprule
\textbf{기호} & \textbf{영역} & \textbf{의미} \\
\midrule
$R_q$ & QROP-Net & 다항식환 $\Zq[X]/(X^n+1)$ \\
$\phi_i$ & QROP-Net & 암호문 원소 $i$의 광학 위상 \\
$U_0$ & QROP-Net & 송신 복소 파면 $e^{i\phi}$ \\
$d$ & QROP-Net & 학습 가능 전파 거리 (mm) \\
$\alpha$ & QROP-Net & FrFT 차수 각도 (라디안) \\
$\xif(s)$ & 공명 & 완성 리만 제타 함수 \\
$\Phiz(t)$ & 공명 & $\xif(1/2+it)$의 위상 \\
$\Ampl(t)$ & 공명 & 진폭 $|\xif(1/2+it)|$ \\
$\Psip(n)$ & 공명 & 매끄러운 소수 계수 위상 \\
$\sigma^*$ & 공명 & 공명점의 실수부 \\
$\ell$ & 공명 & $\sigma$-리프트: $\sigma^* - 1/2$ \\
$F_2(t)$ & 다원적 & 최종 판별 함수 \\
$\Lhat(t)$ & 다원적 & 중심 차분 로그 도함수 \\
$v(\sigma,t)$ & 공통 & $\xif(\sigma+it)$의 위상 / 일반 위상 \\
$\mathcal{P}(\phi;L)$ & 공통 & 위상 양자화 연산자 $\sin^2(L\phi/2)$ \\
\bottomrule
\end{tabular}
\end{center}


\end{document}
